\documentclass[
    a4paper,
    12pt,
    twoside,
    openany,
    UTF8,
    fontset=windows
]{ctexbook}

\usepackage{bookmark}
\usepackage{geometry}
\usepackage{xcolor}
\usepackage{graphicx}
\usepackage{tikz}
\usepackage{setspace}
\usepackage{hyperref}
\usepackage{tocloft}
\usepackage{fontspec}
\usepackage{microtype}
\usepackage{amsmath}
\usepackage{amssymb}
\usepackage{fancyhdr}
\usepackage{titletoc}
\usepackage{textcomp}       % 提供 \textdegree 命令
\usepackage{tocloft}

% 全局英文设为 Times New Roman
\setromanfont{Times New Roman}
\setsansfont{Times New Roman}
\setmonofont{Times New Roman}

\geometry{
    left=2.5cm,
    right=2.5cm,
    top=2.5cm,
    bottom=2.5cm
}

\hypersetup{
    colorlinks=true,
    linkcolor=black,
    bookmarksnumbered=true,
    pdftitle={Therefore, My First Love Was Proven}
}

% 页码设置：正文页脚居中，章首页也居中
\pagestyle{fancy}
\fancyhf{}
\cfoot{\thepage}
\renewcommand{\headrulewidth}{0pt}

\fancypagestyle{plain}{
    \fancyhf{}
    \cfoot{\thepage}
    \renewcommand{\headrulewidth}{0pt}
}

% 万能全页插图（防止插图后空白页出现页码）
\newcommand{\fullpageimage}[2][]{%
    \clearpage
    \thispagestyle{empty}           % 插图页无页码
    \begin{tikzpicture}[remember picture, overlay]
        \node[inner sep=0pt] at (current page.center) {%
            \includegraphics[#1, width=\paperwidth, height=\paperheight, keepaspectratio]{#2}%
        };
    \end{tikzpicture}%
    \clearpage
    \thispagestyle{empty}           % 插图后空白页也无页码
}

% 全局行距：每页约30行
\linespread{1.3}
\setlength{\parindent}{1em}
\setlength{\parskip}{0.3em}

\begin{document}

% ---------- 插图（无页码） ----------
\fullpageimage{1.jpg}
\fullpageimage{2.jpg}
\fullpageimage[angle=90]{3.jpg}
\fullpageimage[angle=90]{4.jpg}


% ---------- 目录（经典小说风格） ----------
\clearpage
\pagenumbering{gobble}
\ifodd\value{page}\else
    \mbox{}\thispagestyle{empty}
    \clearpage
\fi
\thispagestyle{empty}

\vspace*{2.5cm}

\begin{center}
    {\Large\scshape\bfseries Contents}
\end{center}

\begin{flushleft}
\normalsize
\begin{list}{}{
    \setlength{\leftmargin}{0pt}
    \setlength{\labelwidth}{0pt}
    \setlength{\itemindent}{0pt}
    \setlength{\parsep}{5pt}
    \setlength{\topsep}{0pt}
    \setlength{\itemsep}{2pt}
}
    \item Prologue \dotfill \pageref{sec:prologue}
    \item \textsc{Chapter 1} \dotfill \pageref{sec:ch1}
    \item \textsc{Chapter 2} \dotfill \pageref{sec:ch2}
    \item \textsc{Chapter 3} \dotfill \pageref{sec:ch3}
    \item \textsc{Chapter 4} \dotfill \pageref{sec:ch4}
    \item \textsc{Chapter 5} \dotfill \pageref{sec:ch5}
    \item \textsc{Chapter 6} \dotfill \pageref{sec:ch6}
    \item \textsc{Chapter 7} \dotfill \pageref{sec:ch7}
    \item \textsc{Chapter 8} \dotfill \pageref{sec:ch8}
    \item \textsc{Chapter 9} \dotfill \pageref{sec:ch9}
    \item Epilogue \dotfill \pageref{sec:epilogue}
    \item Afterwards \dotfill \pageref{sec:afterwards}
\end{list}
\end{flushleft}

\vspace{1.5cm}

\begin{center}
    \footnotesize\color{gray}
    \textit{A story that begins with tears and ends with tears.}
\end{center}

\vspace{2cm}
\clearpage


% ---------- 正文开始，恢复页码 ----------
\pagenumbering{arabic}          % 恢复阿拉伯数字页码
\setcounter{page}{1}            % 正文从第1页开始

% ==============================
%        Prologue
% ==============================
\clearpage
\thispagestyle{plain}
\phantomsection
\begin{center}
    \large\bfseries
    Prologue
\end{center}
\label{sec:prologue}
\addcontentsline{toc}{chapter}{Prologue}
\vspace{1cm}

I am writing this story in order to head to the bottom of the sea, and to the past——

I came to the port of Soma City in Fukushima Prefecture.\footnote{Soma City, Fukushima Prefecture, located in the Tohoku region of Japan, was severely damaged by the 2011 Great East Japan Earthquake and tsunami.}

The moment I stepped out of the car, I caught the scent of the sea. The wind carried the sound of countless waves, folding over one another, from far and near. Against the overcast sky, a single white speck---a black-tailed gull crying lonesomely---and a fleck of lingering snow: it was still a cold March wind. The wind passed through the desolate port, where the scars of the disaster still remained.

On a small boat floating in the harbor, my childhood friend Shimizu was waiting. Even from a distance, I could tell it was him at a glance---a huge man, making the boat look oddly small. He had filled out even more than before, almost like a lumbering bear that had wandered into the port. Only when I drew near did he finally notice me, and he broke into a grin. It was a smile just like Ebisu, one of the Seven Gods of Fortune\footnote{Ebisu, one of the Seven Gods of Fortune in Japanese mythology, is often depicted as a smiling fisherman and is associated with fishing, commerce, and agriculture.}---a kind, gentle face that belied his burly frame.

``Yacchan!''

``Shimizu!''

Shimizu called out to me in a loud voice. With a grunt, he climbed from the boat onto the pier and rushed over. Then, with full force, he squeezed me tight---a perfect illustration of ``Ordinary Man Attacked by Bear.'' Shimizu had always been like this, intense with his physical affection. But for me, this exaggerated yet warm display of love felt perfectly natural, and I was truly glad for it.

``Long time no see, Shimizu.''

I hugged him back with all my strength and gave his happy-fat sides a soft, pudgy pat.

Shimizu finally released me and said, with a distant look in his eyes.

``It's already been four years since then, hasn't it?''

``Yeah, four years already...''

When I thought of the length of time that had passed, my eyes grew warm with tears.

Shimizu sat down in the pilothouse of the small boat, which had a roof. I sat behind him, using my hard-shell luggage case as a chair, hugging the box that contained fragile items tightly. The engine growled, and the boat began to move. It cut through the sea---a dull, blue-gray expanse that gave off a heavy, metallic gleam---kicking up white spray. As the harbor receded, it grew thin and lost its outline against the horizon; in its place, the majestic Abukuma Highlands\footnote{The Abukuma Highlands, a range of low mountains extending from eastern Fukushima to southern Miyagi Prefecture, known for its gently undulating ridges.} revealed its ridgeline like black waves.

After about forty minutes, the boat stopped. There was nothing around---just the vast, boundless ocean stretching in every direction. I checked the latitude and longitude on my smartphone's GPS.

37\textdegree 49.99' N, 141\textdegree 9.41' E.\footnote{This coordinate lies off the coast of Soma City, Fukushima, in an area heavily affected by the 2011 tsunami.}

It matched exactly the location I had inquired about beforehand with the Fukushima Fishery Cooperative.\footnote{The Fukushima Prefectural Federation of Fisheries Cooperative Associations, which manages local fishing affairs.}

We began preparing for the dive. We put on wetsuits, fitted our masks, strapped on fins, and shouldered our tanks. I had just earned my Advanced Open Water Diver certification in Okinawa specifically for this purpose. Shimizu, the more experienced diver, checked my gear.

Shimizu plunged into the sea with a splash. It reminded me of swimming lessons---even as a child, he was the type to dive in headfirst, bold as an angel. I, on the other hand, was as cautious as a devil---or rather, simply a coward---and eased in slowly from my toes. The water was warmer than I'd expected, but my thin frame soon began to shiver. Shimizu swam over and asked with concern:

``Yacchan, you okay? Your lips are already purple.''

``Already?'' My physique was so pathetic it was almost embarrassing. ``Well, I'm fine. Probably.''

Even through his goggles, I could tell Shimizu wore an expression that said ``Really?''

``All right, let's go. I'll dive first, so follow me.''

Shimizu put his regulator\footnote{A diving regulator, connected to the air tank, supplies breathable air to the diver.} in his mouth and descended gracefully. I followed clumsily after him.

A deep blue world---

The ripples on the surface grew distant, and only the sound of breathing and exhaled bubbles grew louder.

Chasing after Shimizu, who was already five meters down, I left the faintly shimmering surface behind. At about one meter deep, I equalized my ears\footnote{Equalizing ear pressure during descent to prevent pain by adjusting middle ear pressure.} for the third time---the first had been upon entry, the second just after submerging.

Because of the overcast sky, visibility was poor. No fish were in sight. Shimizu, descending deeper and deeper, looked like a solitary giant grouper.\footnote{``Kue'' is a large species of grouper commonly found in Japanese coastal waters.} Afraid of losing him, I carefully followed the silver bubbles that rose swayingly---like Hansel and Gretel\footnote{Hansel and Gretel, protagonists of a Grimms' fairy tale, who use pebbles and breadcrumbs to mark their path.} tracing pebbles that glowed in the moonlight.

As I went deeper, the blue of my surroundings deepened as well.

In the dim silence, detritus fell like upside-down snow.

It looked like salt crystals, I thought.

Salt---for most people, it is nothing more than a seasoning. Those salty white grains in a small bottle, making bland salads tasty and bringing out the sweetness of watermelon---a daily necessity.

But for me, salt is something more special---it is death, it is time, it is life.

The extraordinary twists of my fate had given me this unique image of salt.

As I dove deeper and deeper, that past naturally came back to me. My mind returned to fifteen years ago---when I was still in the third grade of elementary school. Music began to drift through the sea.

A beautiful piano melody---

Frédéric Chopin, Étude Op. 10 No. 3, ``Tristesse.''\footnote{Chopin's Étude Op. 10 No. 3, commonly known as ``Tristesse'' (Sadness), is a lyrical and melancholic piano piece.}

This is a story that begins with tears and ends with tears.


% ==============================
%        Chapter 1
% ==============================
\fullpageimage{5.jpg}

\clearpage
\thispagestyle{plain}
\phantomsection
\begin{center}
    \large\bfseries
    Chapter 1
\end{center}
\addcontentsline{toc}{chapter}{Chapter 1}
\label{sec:ch1}
\vspace{1cm}

\subsection*{1}
``Your mother has salt disease,'' said the doctor, seated on a round stool. He was a man who looked to be in his late thirties, yet also in his early twenties. It was the childish look around his eyes that made his age so hard to guess. Behind square black-rimmed glasses, his eyes were perfectly round, and his thick eyebrows drooped at the ends, as if troubled.

``It's a disease in which the body is gradually replaced by sodium chloride, starting from the extremities.''

Unable to grasp the meaning, I looked up at the nurse standing behind the doctor. She leaned down, matching my eye level, and said:

``Your body turns into salt, little by little, starting from your fingers and toes, and then it crumbles away.''

The beautiful, lightly made-up nurse made a gesture---her right hand like a knife, slowly cutting away her left hand from the fingertips upward. Her right hand finally stopped over her heart.

Dazed, I stared at her hand. When I finally caught up with what she meant, I asked:

``Is Mother... going to die?''

The doctor's face grew even more troubled. He stuck out his lower lip, looking almost like a fish. It was a silent affirmation. Unable to accept reality, I pressed on:

``Her whole body will turn into salt? But how... why...?''

Still looking troubled, the doctor rubbed his lower lip with his right middle finger.

``The human body is made mainly of hydrogen, oxygen, carbon, nitrogen, phosphorus, and sulfur. Well, according to one theory, atoms---no, but why would they turn into sodium chloride---''

The nurse interrupted.

``Nobody knows yet. It's an extremely rare disease---there are only a handful of cases in the entire world. There are many diseases whose causes remain unknown.''

``Then... it can't be cured?''

Silence fell. The doctor didn't even blink---like a fish pretending to be asleep. I returned to my mother's hospital room on the first floor. She was looking out the west-facing window. As if the entire window were a frame, a flowering dogwood was in bloom. A breeze that shook the white flowers drifted gently in through the window. The warm afternoon sunlight, at three o'clock, faintly tinted my mother's thin, light-colored hair.

She noticed me and turned around. She had the expression of a little girl about to be scolded. I sat on the round stool by the bed, clenched my fists on my knees, and said, my voice angry:

``Why didn't you tell me sooner?''

It came out as anger. I must have been angry. It had all happened so suddenly, I was so confused, I didn't even know my own feelings.

``...I'm sorry.''

That was all she said. I understood, at least, that she had wanted to hurt me as little as possible. She was a delicate, gentle person. She would rather feed me something delicious than eat it herself. She would rather disappear than hurt anyone---that was the kind of person she was. Gentle to the point of cruelty.

``Show me your arms.''

My mother rolled up the sleeves of her hospital gown. I caught my breath. From about the middle of her forearms onward, they were gone. The cross-sections were covered with something like clusters of crystal. The fingers that touched the bed felt rough to the touch. Looking closely, I saw white grains on her fingertips.

Salt crystals...

It finally sank in---my mother was turning into salt. My mother was turning into salt, and soon she would be gone. She would become roughness on this bed. The wind that shook the dogwood blossoms would carry her away somewhere...

I cried. I cried as I clung to her stomach and said:

``It hurts, doesn't it... Mother, it hurts, doesn't it...''

I could tell from the muffled sound coming from her belly that she was crying too. Cold teardrops fell on the back of my neck.

``It doesn't hurt... it doesn't hurt...''

She said this as she twisted her body several times. It was a sorrowful motion. She tried to hug me back---but her arms couldn't reach.


\subsection*{2}

I went to Sakuranoshita Elementary School, a public school in Koriyama City, Fukushima Prefecture.\footnote{Koriyama City is a major city in central Fukushima Prefecture, located in the inland part of the Tohoku region.} Cherry trees were planted around the schoolyard, and every spring they bloomed with such vivid color that local people came to view them.

After school, when everyone had gone home and the schoolyard was quiet, I walked alone. Beneath the cherry blossoms that seemed to burn against the blue sky, I circled the schoolyard over and over. The color of the flowers, the call of the bush warbler---none of it entered my absent-minded heart. Only the dappled light and shade between sun and shadow, and the slight differences in scent, flickered faintly like distant memories.

I don't know how many laps I'd made when, suddenly, I heard music begin. A piano's tone---

Perhaps it had just started, or perhaps I had finally noticed something that had been playing all along. Against the clear blue sky, the school building's walls gleamed white. I looked up at the third-floor music room window. Cherry blossom petals danced in the wind.

I thought it was a beautiful performance. For the first time in my life, I felt music was beautiful. It was as if the mud that had clogged my eyes and ears had fallen away all at once, and the world felt vivid. The rustling of the burning cherry blossoms, the aching sense of sound being sucked into the blue sky, the speed of the falling petals---it felt as though something hidden within those things was being drawn out, resonating lightly with the piano's tones.

I stood there in a daze, as if stunned, for a while. Eventually, I headed toward the music room. I went around to the entrance, changed into my indoor shoes, and climbed the stairs warmed by the sun. Strangely, there was no one around. As if a lively sea had suddenly emptied out.

I crossed the dim hallway and stood before the music room. The small window on the sliding door was covered by a pitch-black blackout curtain. I put my hand on the door---and hesitated. I was an uninvited guest. But I had to know who was playing. I slid the door open quietly, so as not to make a sound.

The grand piano was in the back left; the player was hidden in shadow, invisible. All I could see were slender feet pressing the pedals. I crept closer---the piece was in a climactic section, and it felt somewhat unsettling. When the melody softened again, I could see the performer.

In an instant, I was captivated.

She was a beautiful girl. Her bangs were cut straight at her eyebrows, and she was lost in the music, her long lashes lowered dreamily. A breeze carrying the scent of cherry blossoms stirred her glossy black hair. The faint light from the window shone through her pale white skin, and only her cherry-colored lips slid like tiny pearls. On her slender body, a vivid sky-blue dress---it seemed as if a scrap of spring sky had whimsically descended to earth.

\---As if the wind had stopped, the performance ended.

A sunlit moment passed. The bush warbler's call began again. The girl opened her eyes with a start and looked at me. Large almond-shaped eyes---as if two flames had suddenly kindled, as if two flowers had suddenly bloomed---they were intensely vivid eyes.

It was as if time had stopped. How long had we been looking at each other?

``That was really, really good,'' I said, coming back to my senses. That was all I could find to say.

``...Thank you,'' she said, a little puzzled, and then suddenly smiled. I smiled back. She leaned forward slightly from her chair and said:

``I was watching you. I thought, 'What a strange person.' ''

``A strange person?''

``You were walking around the schoolyard in circles, weren't you?''

I gave a wry smile. My cheeks felt a little warm, and I said, to cover it up:

``I was lost.''

``You're quite the directionally challenged person.''

She laughed with amusement. When she laughed, her big eyes narrowed, and her puffy tear troughs were adorable. She seemed to have taken a sudden interest, leaning forward to ask:

``Hey, what were you really doing? You were collecting something in your pockets, weren't you?''

Caught by her sparkling, curious eyes, I gave in and said:

``I'll warn you in advance---there's a deep reason for this. So please don't laugh.''

``I won't laugh, I won't laugh.''

A mischievous smile on her face, she held out both hands as if to scoop water. I sighed, walked up to her, and dropped the contents of my pockets into her palms.

Cherry blossom petals scattered onto her snow-white palms.

The girl looked puzzled, and met my eyes.


\subsection*{3}

To explain that deep reason, I have to go back in time to when I was three years old. It was then that my young, soft sensibility twisted and never returned---like heated glass cooling and hardening into a strange shape.

The trigger was my father---Saegusa Ryunosuke.\footnote{Saegusa Ryunosuke shares his name with the famous Japanese literary giant Akutagawa Ryunosuke (1892--1927), known for works such as \textit{Rashomon} and \textit{In a Grove}. The shared name is a deliberate, ironic choice by the author.} He is a novelist, praised for his unique style and unpredictable storylines, born from his distinctive sensibility. Despite sharing a name with the great predecessor Akutagawa Ryunosuke, he has the audacity to work under his real name. From the fact that he named his son Yakumo\footnote{Yakumo (八雲) is a name with literary and historical resonance in Japanese culture, often associated with the poet and scholar Lafcadio Hearn (who took the name Koizumi Yakumo).} alone, his arrogant, irreverent personality seeps through.

One evening, we were walking together along the Abukuma River.\footnote{The Abukuma River flows through Fukushima and Miyagi prefectures, one of the major rivers in the Tohoku region.}

``Dad, why do you only have one eye?'' I asked, at age three. My father stroked his stubbly chin and let out an eerie ``Heh heh heh.''

``Because it was in the way, so I gouged it out myself. Pop.''\footnote{The Japanese onomatopoeia \textit{zuppo} (ズポッ) suggests a wet, sudden plucking or pulling sound.}

``...Where did your eye go?''

``I ate it myself.''

I froze in horror and shouted: ``That's a lie!''

My father turned, walked back to me, crouched to my eye level, and said: ``It's true, kiddo.''

Then he flipped up the black eyepatch over his right eye.

A small, pitch-black void gaped open. The blood-red evening sun did not illuminate that hole at all. The murmur of the Abukuma River, the light dancing on the water's surface---it all felt as if it were being sucked into that hole, never to return. In that instant, my soft sensibility twisted.

I felt a strange pain. I felt pain in the darkness of the lost eye. The ``absence'' of my right eye was ``pain''---not the pain of a wound, but the absence of what should be there. That ``emptiness'' began to bring me pain. For example, when the tail of my Godzilla figure broke off, I would cry from the pain of the ``emptiness'' left by the lost tail.

\---A few months later, I fell down the stairs of our apartment building. For a while, I lay curled up on the landing between the second and first floors. Eventually, I endured the pain, got up, made my way back to our third-floor apartment, and looked into the bathroom mirror. My left temple was split open, blood gushing out. When I wiped it away, I could see bone.

Yet I remained calm. Composed, I pinched the wound with my fingers and pressed the gap tightly together. I thought: I see. I'm not missing a single thing. The split came together perfectly. There was no ``emptiness'' anywhere. It felt to me like it wasn't real pain.

After covering it with several bandages to stop the bleeding, I felt relieved and started watching anime on TV. My temple throbbed, but it felt distant, like someone else's pain.

When my mother came home, she screamed. She took one look at the wound on my temple and burst into tears. I didn't quite understand why she was crying. But it saddened me to see her cry, so I cried too. Mother cries, and I'm confused---this pattern was reversed when my mother fell ill with salt disease. The emptiness left by the loss of my mother's limbs caused me intense pain.

``It hurts, doesn't it... Mother, it hurts, doesn't it...''

And my mother, not understanding my pain, cried simply because she was sad to see me cry.

``It doesn't hurt... it doesn't hurt...''

She kept saying this, trying to hug me, but unable to reach me...

I didn't know how to handle this strange, mysterious pain---like a phantom limb pain---but one that belonged to others, to things that didn't exist. I took it into my heart and didn't know how to deal with a wound I couldn't see. A special wound needed a special bandage. Eventually, I found a way to ease that pain.

It was to collect things to fill the wound. Anything would do: a twig, a pretty pebble, a shard of glass.

What mattered was to pray. To pray that the collected things would properly fill the wound that was the source of the pain. That they would heal the pain. So I searched for something to fill the emptiness left by my mother's lost limbs. At first, I walked around the classroom. How about the big triangle ruler the teacher used on the blackboard? No, that would make my mother look like a Gundam.\footnote{Gundam (ガンダム) is a famous Japanese mecha anime franchise featuring giant humanoid robots.} Chalk? Someone's forgotten pencil case?---that's when I noticed the cherry trees in full bloom in the schoolyard.

Against the cool spring sky, the cherry blossoms' color burned in my eyes. When the wind blew, they looked like sparks flying. I turned my eyes from the sky to the ground. Petals that had fallen in the shade of the trees lay scattered like distant wildfires on a mountain. When I timidly touched them, contrary to expectation, they were cool. When I placed them in my palm, it seemed as though warmth slowly seeped out from within that quiet, cool outline.

I felt that the cherry blossoms petals could fill my mother's emptiness. That they would become new limbs, along with the wind. That they would warm the cold pain.

So I collected cherry blossom petals. Walking around the schoolyard, one by one, little by little. I kept walking and kept collecting, praying that my mother might be healed, even just a little.

% ===== 以下为 Chapter 1 新增部分（第4～10节） =====
% 请将此代码插入到您原有 Chapter 1 的末尾（第3节之后）

\subsection*{4}

I sat on a desk right next to the piano, staring at my folded hands for a long time. I was a sensitive child, and I felt ashamed to talk about myself. When I finished speaking, I finally looked up. The girl had lowered her face. For a moment, I thought I saw something glistening on her cheek. Was she crying---?

But she quickly rubbed her face and looked straight at me. Her eyes seemed a little red, but in the end, I couldn't tell whether she had really shed tears.

``You're---strange. Feeling pain in emptiness is especially strange, but your other sensibilities are also unusual. Normal children don't perceive cherry blossoms as flames.''

Being called ``strange'' made my body grow hot. I was at that age when being different from others was embarrassing. Seeing me like that, the girl said hurriedly:

``Oh, but I understand. Passionate flowers---for example, if I tried to express a vivid red rose or bougainvillea on the piano, I would play them like flames. So maybe you're just more sensitive than ordinary people? Like how I have perfect pitch.''

``You have perfect pitch?''

I stood up and walked over to a guitar lined up against the wall opposite the piano. I sounded an A note---I only knew how to play the notes Do, Re, Mi, Fa, Sol, La, Si, Do. ``Do you know what note this is?''

The girl smiled. ``It's 'La.' Maybe slightly leaning toward 'Sol♯.' Even with perfect pitch, for some reason the region around 'La' is a bit fuzzy for me. I get pulled toward 'La.'''

``Wow, that's amazing!''

``Hey, what did you think of my playing?''

``It was really good.''

``Not that---in your own, different words, what did you think?''

I thought for a while, locking eyes with her captivating gaze. Then I said:

``...I felt no pain at all. Every single note was beautiful, and it felt like each was in its proper place. As if it had been decided from the very beginning... what do you call that?''

``...Fate?''

``Yes. It was as if fate itself was playing.''

``As if fate itself was playing...''

She repeated, looking somewhat surprised, as if turning over something found in the dark to examine its true nature.

Suddenly, like a flower bursting into bloom, her face broke into a smile. Then, looking a little shy, she said:

``I'm Igarashi Yuzuki. Written with characters meaning 'swaying moon.' And you?''

``Saegusa Yakumo. 'Eight clouds.'---What was the name of that piece you just played?''

``'Tristesse.' Composed by Frédéric Chopin.''

That was how Yuzuki and I met.


\subsection*{5}

After promising to meet again tomorrow, I parted with Yuzuki and headed to Sato General Hospital, where my mother was admitted.\footnote{Sato General Hospital is a fictional name, but typical of Japanese community hospitals.} I rode my bicycle for about forty minutes, arriving around five in the afternoon. The sunset was dyeing the hospital's white walls orange. I parked my bicycle in the lot and went through the automatic doors. The familiar smell of disinfectant stung my nose. After checking in at the reception, I headed to Room 108, where my mother was---

My hand on the door froze. I heard voices from inside. I slid the door open a crack and peered through the gap.

A ``shadow'' stood in the twilight hospital room. My father---Saegusa Ryunosuke. After he divorced my mother when I was five, he had become a shadow to me. A tall, lanky, stooped, unreliable shadow...

Perhaps it was because he always wore black clothes. But the decisive reason was that when he left home, I didn't feel much pain. Even though a person had vanished, there was only a vague, itchy sensation. That was why I thought of him as a shadow from the start. A formless, blurry shadow.

Only the void of his right eye, shown to me on the banks of the Abukuma River, was real to me. I hid in a corner until night came and the shadow went somewhere else. I didn't want to see him. I had no words to say to that inscrutable shadow.


\subsection*{6}

The next day, I was distracted throughout class. Thinking about my mother's illness made me unbearably sad, and thinking about Yuzuki, whom I was to meet after school, made my heart race and left me restless. Time, on one hand, was stealing my mother's life; on the other hand, it was bringing Yuzuki and me together. Caught between these conflicting emotions, I couldn't focus on my studies.

During break, I asked Shimizu, who sat behind me: ``Hey, do you know Igarashi Yuzuki from Class 3?''

Shimizu looked puzzled. I thought I might have asked the wrong person. Shimizu had always been big, slow, and far from being well-informed. He looked genuinely surprised and said: ``Yacchan, you don't know about Igarashi? Really?''

It turned out I was the one out of touch. Yuzuki was quite famous. She had started piano at age two, won the gold prize in the first and second grade elementary division of the Chopin International Piano Competition in ASIA at seven, and the following year won the gold prize in the Concerto A category at the same competition, becoming the youngest winner in history---in short, she was a piano prodigy.\footnote{The Chopin International Piano Competition in ASIA is a regional variant of the famous Warsaw competition, held in various Asian countries.} She was expected to become a world-class pianist without a doubt. Of course, since Shimizu couldn't even remember how many home runs he'd hit, at this point he only knew that ``she'd apparently won some amazing prize.''

I was surprised, but at the same time, I felt a sense of ``I knew it.'' Yuzuki's playing was overwhelmingly brilliant, even to a layman like me.

In my head, I built a fanciful image of Yuzuki: a ``recluse heiress.'' Living in a white mansion, playing the piano through open windows, occasionally coughing delicately---a frail, elegant young lady.

That stupid stereotype was shattered at lunchtime. I was playing dodgeball in the schoolyard with Shimizu and the others when Yuzuki suddenly appeared. She was leading a group of girls from Class 3, striding across the schoolyard with a frown, as if heading into battle.

I watched in astonishment as the girls' coalition stopped in front of the boys of Class 3, who were playing soccer. At the center, standing dumbfounded, was Sakamoto---a boy with good athletic skills, the class bully. Yuzuki stepped forward and said to Sakamoto:

``You, stop bullying Koyocchan!''

Sakamoto's eyes widened as he looked at Kobayashi Koyomi, who was hiding behind Yuzuki. She was a small, cute girl, gripping the hem of her skirt, looking down. Sakamoto stiffened and said:

``Hah? Why would I have to bully someone like her...''

His voice lacked conviction. Yuzuki said in a clear, sharp voice: ``You like Koyocchan, don't you!''

Sakamoto's face said ``Gah!'' Kobayashi turned bright red and her mouth opened and closed helplessly. By now, they were the center of attention across the entire schoolyard. Sakamoto yelled, clearly flustered:

``...Sh-shut up! I don't like such an ugly girl!''

Kobayashi's eyes instantly filled with tears and she started sobbing. Yuzuki's eyebrows shot up in anger: ``Hey, you apologize to Koyocchan right now!''

``Hah? Why do I have to---owowowow!''

Yuzuki had grabbed Sakamoto's nose with her right hand and twisted it upward. You should never underestimate a pianist's grip strength. Sakamoto screamed, grabbing her right arm with both hands and struggling. I half-expected his feet to leave the ground---that's how formidable Yuzuki was.

``For a guy, you're so uncool!''

Just then, a teacher who had heard the commotion arrived. He calmed things down and marched everyone off.

``Igarashi is scary...'' Shimizu muttered, shuddering.

This incident shattered my image of Yuzuki as a ``recluse heiress.'' ``Genius pianist'' and ``Amazon warrior''---that was her new image.


\subsection*{7}

After school, at the appointed time, I headed to the music room. I could hear the beautiful sound of the piano---powerful and delicate. It was strange to think that those same fingers had twisted Sakamoto's nose.

Even when I entered the music room, Yuzuki didn't notice. She was completely absorbed in her playing. I brought a chair and sat down, watching her play. With her eyes nearly closed, listening intently, pouring emotion into each note, she looked almost like she was praying.

When the piece ended and the last reverberation faded, I clapped. Yuzuki gasped in surprise.

``If you were here, you should have said something!''

``I didn't want to interrupt.''

``Geez...'' Her cheeks flushed pink. She looked so charming that I was again unsure whether the recluse heiress or the Amazon was her true nature.

Yuzuki closed the piano lid and abruptly said: ``Well then, let's go.''

``Where?''

``My house.''

I didn't quite understand how we got to that, and I tilted my head about fifteen degrees, but Yuzuki ignored that and tugged me along. After about three minutes of walking, she finally said:

``There was that trouble today, so I can't concentrate much at school...''

I thought with my head still tilted and said: ``Trouble? You just thrashed him one-sidedly, and you seemed pretty focused a while ago.''

``You were watching---?'' Yuzuki pouted with embarrassment. Then, as if making an excuse: ``I hate guys like that. They're such a pain, you know?''

``I think all humans are a pain.''

Yuzuki looked at me with a sour expression. What was that face for...?

It was a beautiful day. Sakuranoshita, where we lived, was a new town with many well-kept gardens. Violets, azaleas, lilacs---spring flowers were blooming everywhere.

Suddenly, a dog barked along the way. Yuzuki groaned and swapped positions with me. A golden retriever was tied up by a front door, wagging its tail furiously and barking at Yuzuki. ``Does that dog hate you?'' I asked.

``No. It likes me too much and gets so excited it pees on me.''

Using me as a shield, Yuzuki hugged my body and reached out to pet the dog's head. I was torn between the closeness of Yuzuki and the fear of being peed on, and I ended up experiencing a strange kind of heart-pounding.


\subsection*{8}

Yuzuki's house was much larger than the surrounding ones. Surrounded by high walls and hedges, it was barely visible from the outside.

Passing through the stylish double gates, we entered a well-kept garden. In the corner stood figurines of the Seven Dwarfs and their dwelling, as if they tended the garden every night for Snow White. On the large botanical-style terrace, white wisteria flowers hung like curtains, positioned just right to soften the evening sun.

We went inside. The house had large windows. Grimm's fairy tale ornaments were tastefully arranged: the rose from Beauty and the Beast, Cinderella's clock, Little Red Riding Hood's basket for visiting her grandmother...

The piano was in a soundproofed room on the west side of the house. It was a ten-tatami room separated by large glass doors, with bright sunlight streaming in through big double-paned windows.\footnote{A tatami is a traditional Japanese mat, roughly 1.8 m × 0.9 m; ten tatami is about 16.2 square meters.}

Yuzuki set a CD of Maurizio Pollini\footnote{Maurizio Pollini (1942–2024) was an Italian pianist renowned for his flawless technique and interpretations of Chopin and other composers.} into the player. On top of the speaker sat a small ornament of the Bremen Town Musicians.\footnote{The Bremen Town Musicians are characters from a Grimm's fairy tale about a donkey, dog, cat, and rooster who become musicians.}

Chopin's \textit{Barcarolle} began to play.---When it ended, Yuzuki asked: ``What did you think?''

``It was astonishingly beautiful.''

She nodded with satisfaction. Then she sat down at the Steinway\footnote{Steinway \& Sons is a prestigious American-German piano manufacturer, widely considered one of the finest in the world.} piano. On the piano sat a stuffed Anpanman doll.\footnote{Anpanman is a popular Japanese children's anime character with a bean paste bun for a head, created by Takashi Yanase.} I thought it was an oddly childish taste. Yuzuki began to play the \textit{Barcarolle} herself.

Once again, I felt Yuzuki's piano was extraordinarily beautiful. Each note was pure, like particles of light. Her \textit{Barcarolle} had a quaint charm, as if the particles of light drifting on Venetian canals gradually took the shape of flowers, swirling and flowing with a fragrant scent.

I managed to convey this in clumsy words. Yuzuki said ``thank you,'' but didn't seem very pleased. ``Not that---compared to Pollini's performance, how was it?''

Pollini is one of the world's greatest pianists, known for his flawless playing. The 1972 recording of \textit{Chopin: 12 Études, Op. 10 \& 25} even bore the tagline ``What more could you want?'' To ask me to compare myself to such a master was terrifying. Of course, I knew none of this at the time, and gave an honest, offhand opinion:

``...It felt kind of flat? Like there's a boat and a water surface, but no sense of what's beneath it?''

It was a terribly abstract and vague opinion. Yet Yuzuki nodded as if convinced.

``I thought so too. It's too neatly put together. I need more maturity and profundity.''

``Err, maturity and profundity...''

It was a sensibility hardly befitting a third-grader.

``After all, I don't have enough life experience to play the \textit{Barcarolle} well. They say you can't play it properly without going through three heartbreaks. Sigh, I want to fall in love soon.''

I couldn't help but chuckle.

``I might be strange, but you're quite something too.''

``I'm seriously troubled here.''

Just then, a woman in her early thirties appeared on the other side of the glass door. Since it was soundproofed, I hadn't noticed any outside sound.

``Oh, my mother's back...''

I observed the woman: round sunglasses, loosely waved long hair dyed a flashy brown, a yellow knit jacket over a navy skirt. She had a great figure, and in hindsight, she had a stylish, Milanese air. Igarashi Ranako---that was her name. Shimizu had told me she was a professional pianist. When she took off her sunglasses, her beauty was revealed. She certainly resembled Yuzuki, but with a sharp, prickly kind of beauty. She glanced at me as if I were some kind of nuisance, then strode briskly to another room. I was perplexed, wondering if I'd done something wrong.

``Sorry---'' Yuzuki said apologetically. ``My mother can be harsh. I have a lesson now, so let's call it a day.''

I nodded and put on my shoes at the entrance. As I said goodbye to Yuzuki, I saw Ranako behind her. I said, ``Thank you for having me,'' and she replied curtly, ``Yes.'' Then she waved her hand dismissively, as if shooing me away.


\subsection*{9}

My mother gradually turned into salt. It was like an hourglass. As the sand flows down and never returns, so my mother slowly disappeared and never came back.

During this time, whenever I woke up, I was always crying. I had some empty, frightening dream, and I woke up sad and weeping. When visiting hours ended and I had to leave, I would cling to my mother and cry like a baby. I didn't want to go home alone to a dark room, nor did I want to leave my mother alone in that dark hospital room. At such times, my mother would stroke my back and keep saying:

``It's okay... it's okay... it's okay...''

I kept gathering flowers for my mother. When heavy rain scattered the cherry blossoms, I picked wildflowers whose names I didn't know. Picking wildflowers made them disappear. The emptiness left by the flowers vanishing from my commuting path and around the house brought me an indescribable guilt and pain—it was hard. Yet I had to continue, to fill the greater pain inside me.

I listened to Yuzuki's piano and gave opinions that I wasn't sure were useful. Yuzuki, in turn, listened to me---about my mother. When I told Yuzuki my memories of my mother, I felt the pain of losing her ease slightly. It felt as if Yuzuki remembering her somehow supplemented the mother who was fading away.

One day, when I told her about my guilt over picking wildflowers, Yuzuki was surprised.

``Aren't you being a little too kind?''

Then she took my hand and led me to the back of the Igarashi house. We walked along a narrow animal trail on a gentle slope. The sky opened up, and I gasped involuntarily. In a small basin-like meadow, colorful wildflowers bloomed in profusion. Their colors, never mixing, became a vivid water that filled my heart with an indescribable warmth.

``From now on, you can pick flowers here---'' Yuzuki said with a gentle smile. ``No matter how many you pick, the flowers won't disappear. The love the world gives is far greater than you think, and flowers are much stronger. It's like filling a bucket with seawater easily. Filling just your own heart is no big deal. The only ones who feel unsatisfied are those with holes in their buckets.''

Hearing that, I felt saved all at once. It was okay to pick flowers. It was okay to fill my sadness with flowers. The strength of the world would forgive me. That thought settled into my heart, not as logic, but as a conviction.

Yuzuki had taught me, in an instant, what was too obvious for anyone to put into words. I felt a deep affection, something like respect, for Yuzuki.


\subsection*{10}

Sad and happy days passed.

I studied at school, became a shield on the way home to protect Yuzuki from the dog that peed out of sheer joy, listened to Yuzuki's piano, and let her listen to my stories about my mother. I picked flowers behind Yuzuki's house and arranged them in the vase in my mother's hospital room. And I went home crying with sadness.

Yuzuki listened to my stories about my mother with visible delight and envy.

``Lucky you---I wish I had a mother like that.''

``But you have a mother, a beautiful pianist.''

I will never forget Yuzuki's face at that moment. A faint smile hovered on her lips, but her eyes were open in a confused, wounded way. It was the face of a lonely, lost girl, unbearably sad.

I realized the meaning of that expression much later.

One holiday, passing near the Igarashi house, I wondered what Yuzuki was doing. Come to think of it, we had never met on a holiday. So I stopped by and rang the doorbell. There was no answer.

At first I thought they were out, but then I thought she might be in the soundproof room, where the bell couldn't be heard. I went around to the back of the house and peered through the double windows into the soundproof room.

Yuzuki and Ranako were there. I immediately knew I was seeing something I shouldn't.

Ranako's face was bright red, and she was shouting like a demon. Ruthless slaps flew. A sharp smack hit Yuzuki's head, scattering her black hair. Her slender shoulders trembled.

I stood frozen, eyes wide open. Yuzuki wiped her tears and started playing again, her hunched back quivering. Ranako shouted. A slap flew. ``Why can't you do it!''---over and over. I felt as if I could hear every sound: the piano note that always stopped at the same place, the slap, the shouting, the sobbing, the compulsive ticking of the metronome.

Anpanman's smiling face on the piano looked as if it was about to cry.

Just then, I heard the front door click open. Yuzuki's father, I thought. Probably he had been upstairs and finally responded to the doorbell. I peeked through another window into the living room. As expected, a gentle-looking, slightly overweight man in glasses---Igarashi Sosuke---stood there. He held a newspaper in one hand, looked toward the soundproof room, and seemed to want to say something. Beyond his gaze, Ranako's abusive lesson was surely taking place. I wished he would say, ``Stop that.'' Stop the violence, stop making her cry, stop making her hate the piano.

But Sosuke didn't. With a resigned look, he hunched his shoulders and trudged back upstairs. I felt utterly hopeless. My heart turned cold and gray. Unable to do anything, I left the Igarashi house and wandered around.

``Lucky you---I wish I had a mother like that...''

Now I understood the meaning of those words. Yuzuki had wanted Ranako to be kinder to her. Not to place excessive expectations on her, but to accept her gently and shower her with love, even if she couldn't play the piano well.

At last, I felt I had seen Yuzuki's true nature. She was neither a recluse heiress nor an Amazon. She was that sad piano note that always stopped at the same place and faded away---

Suddenly, I noticed the usual golden retriever looking at me, wagging its tail. I walked over casually and patted its head. It was overjoyed and soaked my clothes with pee.

``So you've fallen for me too, huh---'' I muttered with a wry smile, and suddenly tears overflowed uncontrollably.

I understood why Yuzuki always petted this dog even though she risked being peed on. Yuzuki loved this dog too. She was happy that it would wet itself with joy just from her coming to see it, without her having to play the piano. She was saved by its pure, straightforward love.

Thinking of lonely Yuzuki, I couldn't help but cry. The pee was cold and smelly, yet it made me happy in a way I couldn't explain. I hugged the dog and sobbed. The dog kept peeing on me while licking my tears with its tongue.

The dog's name was Melody.


\fullpageimage[]{6.jpg}
% ==============================
%        Chapter 2
% ==============================
\clearpage
\thispagestyle{plain}
\phantomsection
\begin{center}
    \large\bfseries Chapter 2
\end{center}
\addcontentsline{toc}{chapter}{Chapter 2}
\label{sec:ch2}
\vspace{1cm}

\subsection*{1}

There was a phone call early in the morning. It was the call that my mother had passed away. August 2nd---right in the middle of summer vacation. That day was so hot it felt like it would melt you. Against the deep blue sky, thick cumulonimbus clouds piled up, as if a gigantic bomb had fallen beyond Mount Adatara.\footnote{Mount Adatara (安達太良山) is a stratovolcano in Fukushima Prefecture, located west of Koriyama City. It is a well-known landmark in the region.} My hollow heart thought: the world had ended silently, and only the cicadas' cries echoed...

Even though my mind was spinning in vain, the dry, rattling wheels still gripped the ground, and before I knew it, I had arrived at the hospital. I parked my bicycle in the lot, flapped my sweat-soaked T-shirt, entered the cool building, and went to the reception desk.

To be honest, my memory from there on is vague. I don't remember what expression the usual receptionist had on her face. She might have widened her eyes in surprise, or perhaps she was perplexed by my unexpectedly calm demeanor.

Before I knew it, I was in the mortuary for a viewing. Using the word ``viewing'' felt deeply unnatural, because my mother was no longer in human form. On the long mortuary bed sat a large, round glass vase, filled with pure white salt. It was the vase I had been putting flowers in. My mother had probably prepared it for herself, knowing she would become salt, but since I suddenly started bringing flowers, she must have used it as a vase in the meantime. That's how kind she was.

My mother had passed away in the middle of the night, as if falling asleep, and apparently had completely turned to salt within one night. After losing life, the salinification progressed rapidly.

It was only when I stood before the vase that I finally thought, ``Ah, my mother is dead.'' Until then, I had felt a sense of disbelief. Perhaps it was a lie or a joke. There was even a possibility that when the magician pulled back the curtain, my mother would be there with her arms and legs restored, wearing a mischievous smile.

I was a child.

My mother's death pierced my chest with unbearable pain. I hugged the cold vase and cried until I broke. Yet at the same time, I thought, ``I'm glad.''

I was glad that my mother had been freed from her suffering.

As salt disease progresses, the body eventually begins to turn to salt internally. The membranes wrapping the internal organs become salt, and the organs rub against each other, causing excruciating pain. She had reached the point where even the slightest movement made her grit her teeth in pain, and in the end she was given morphine.

Despite all that pain, my mother said she didn't want to die. She didn't want to leave me alone. That made me feel guilty and miserable. It felt as if I was the one making her suffer, and I wanted to disappear from the world.

I prayed that my mother's soul would be in a place without pain, a warm place full of sunlight. A place where cats naturally gather and lazily start sunbathing. A place with a perfectly comfortable lounge chair, with oranges growing nearby. A place where the world's most interesting novel is always within reach. A place where a pleasant breeze always blows---

I wanted my mother's soul to be saved in such a place, filled with light.


\subsection*{2}

Since my mother had no relatives at all, my father came and held a small, private funeral at home. A monk chanted sutras before the altar. I glanced beside me and saw my father crying. Tears spilled from his one remaining eye. I felt conflicted.

After all the pain he had put my mother through, how dare he cry now?

But his tears looked so genuinely sad, and there seemed to be love in them, that I almost forgave him---and hurriedly shut that feeling away with reason.

After the funeral, I rode in my father's black Mercedes heading for Iwaki City.\footnote{Iwaki City (いわき市) is a coastal city in southern Fukushima Prefecture, facing the Pacific Ocean. It is known for its beaches and aquamarine park.} It was a journey of two---me and a shadow. My father talked about nothing but elusive topics.

``Do you know why foreigners are all so tall?''

I stared blankly at the scenery passing outside the window. My father continued without caring.

``It's because their facial features are deep.''

I couldn't help but say, ``Huh?'' I had been caught on the hook.

``Why would deep features make you tall?''

My father grinned.

``If your features are deep, the bones of your forehead make it hard to see upward, right? That makes you vulnerable to attacks from above. If you're tall, it's easier to attack enemies from above and harder for them to attack you from above. In other words, your chances of survival increase. Eventually, through natural selection, the tall ones survive.''

``Oh... so that's why...''

I was honestly impressed. My father grinned again.

``You didn't really believe that, did you, kiddo?''

I was furious, but I felt that showing my feelings would be a loss, so I stayed silent.

---Soon, we saw the sea.

When we got out of the car, it was cooler than in Koriyama City. In the moderately calm summer heat, the sea glittered with a gentle light. We walked to the end of an inconspicuous pier and took out my mother's salt from a furoshiki.\footnote{A furoshiki (風呂敷) is a traditional Japanese wrapping cloth, often used for carrying items.}

``You can't just scatter ashes anywhere,'' my father said. ``But if it's salt, I suppose it's all right.''

It was my mother's dying wish that when she became salt, it be scattered in the sea. My father and I gently took the salt in our hands and sprinkled it little by little. A soft wind blew, and the salt sparkled like jewels. Then I scattered the flowers I had brought. They fell, spinning like parachutes, as vivid as colorful lotuses blooming on the water's surface. Eventually, the waves carried them away. I hoped that someday they would reach my mother. She would surely display them in a new vase.

My father and I stared out at the sea for a while. I think we had some ordinary, trivial conversation---the kind any parent and child might have.

As dusk fell, my father gradually turned back into a shadow again. An elusive, expanding and contracting, unreliable shadow... The shadow said:

``So, are you going to live with me from now on, kiddo?''

My heart wavered like the flowers in my line of sight. But I still had some defiance left, so I asked:

``Why did you divorce my mother?''

My father was silent for a while, then said:

``Remember when I told you why I don't have a right eye? That was the truth.''

I caught my breath. ``Why on earth...?''

``I told you: it was in the way. Nobody believes me, but at some point, my right eye became unbearably annoying. I don't know why. But truly, it was so annoying I couldn't stand it. So I gouged it out myself. I felt bad for my dead parents, so I ate it.''

It was a story I didn't even want to imagine. I felt as if reality itself was warping.

``As soon as I lost that eye, miraculously, I gained literary talent. Things I hadn't understood became clear, and things I hadn't seen became visible. That's how I became a novelist.''

I thought carefully about this incredible story and asked:

``What does that have to do with divorcing my mother?''

My father hesitated briefly, then said:

``I loved your mother. ...But at some point, even your mother became unbearably annoying. I couldn't write novels anymore. I had no choice but to separate from her to survive.''

I was dumbfounded, and for a while I lost all sense of emotion. Then, from the bottom of that hollow heart, anger began to well up, as if slowly starting to boil.

``...You're out of your mind, aren't you?''

My voice trembled. My father's tone didn't change.

``...I'm sorry. That's how it is. I'm a defective human being.''

``...Don't screw with me. Do you have any idea how much Mother suffered?''

``I truly regret it. But I did my best with what I had.''

``This isn't about effort. ...You're... you're... go to hell.''

My father looked hurt. I grabbed my mother's vase and walked toward the land. I kept walking, heading for Koriyama. At night, I walked, crying alone. The emptiness inside the large vase I held to my chest throbbed with pain. When I could no longer take another step, I spent the night on a bench at a bus stop. The next morning, I woke to the chirping of sparrows and walked again, still crying.

In the end, I arrived at my apartment that night. I had walked for over twenty hours.


\subsection*{3}

I remember that summer as being bitterly cold, even though the temperature was very high.

The emptiness of the room where my mother would never again be became a terrible pain that tormented me. The household chores I had been managing daily, and all other activities, came to a complete halt. I had to go pick flowers to suppress the pain. But I couldn't move at all.

Even though I had a fever, I didn't sweat at all; I just shivered and wrapped myself in blankets. I couldn't move a step from there. Occasionally, tears would overflow like a melting snowman. I no longer even knew if I was sad. I was frightened by the slightest sound, and my mind became a jumbled mess.

I think I went about five days without eating or drinking anything. If things had continued that way, I might have died. Like a candle burning out, leaving only a wisp of smoke.

But then, the doorbell rang.

I couldn't move. Like a squirrel waiting out a storm, I hid under my blanket, hoping the visitor would go away. But the visitor kept ringing the bell persistently.

``Yakumo-kun, you're there, aren't you!''

It was Yuzuki's voice. Startled, I suddenly regained some awareness and crawled out of bed. I managed to stand up, but a terrible dizziness hit me. My whole body felt like it was made of heavy clay. When I tried to speak, only a faint breath of air passed through my throat.

I clung to the wall, finally made it to the front door, and unlocked it. The door was immediately pulled open. The brightness of the sunlight I hadn't seen in so long made me squint. Yuzuki saw me and was stunned.

``Yakumo-kun... are you alive?''

Perhaps it was a joke, but I was in no state to laugh.

Without asking anything, Yuzuki left for a while, then returned carrying shopping bags in both hands. She had bought pudding for me. I had long passed hunger and had no appetite, but when I forced myself to eat it, I felt a little better. I must have been suffering from low blood sugar.

Yuzuki stood in the kitchen. I could hear the rhythmic sound of a knife: \textit{sutoton-toto}. It was like the sound of time beginning to move again. She had made rice porridge for me. The vibrant colors of the mitsuba\footnote{Mitsuba (三つ葉), also known as Japanese honewort or trefoil, is a herb commonly used in Japanese cuisine for its mild flavor and decorative value.} and crushed pickled plum immediately whetted my appetite.

When I had eaten it all, a spoonful of tears remained at the bottom of the clay pot.


\subsection*{4}

I told Yuzuki that my mother had died. Yuzuki cried with me, saying she wished she could have met her just once.

From that day on, Yuzuki came to my apartment whenever she had time. She cleaned my room and cooked delicious meals for me. With her beautiful black hair tied back, wearing a bright red apron as she did housework, she looked somehow grown-up and made my heart pound. The only thing that was incongruously childish was that the apron's pattern was Anpanman.

``...Why do you do so much for me?''

When I asked, Yuzuki, still making the \textit{sutoton-toto} sound with her knife, said:

``Because if I don't take care of you, you'll die, right? Like a goldfish.''

Perhaps that was true. Though I didn't know about the goldfish part.

The summer days with Yuzuki passed. We often listened to music together, playing the CDs she brought.

Yuzuki's favorite pianist was Tanaka Kiyoko.\footnote{Tanaka Kiyoko (田中希代子, 1932--1990) was a Japanese pianist who became the first Japanese prize-winner at the 5th International Chopin Piano Competition in 1955, placing 10th.} She was a woman who, in 1955, became the first Japanese recipient of an award at the 5th International Chopin Piano Competition---she placed tenth. It is said that the judge, Arturo Benedetti Michelangeli,\footnote{Arturo Benedetti Michelangeli (1920--1995) was an Italian pianist considered one of the greatest of the 20th century, known for his flawless technique and demanding standards.} was dissatisfied with the ranking, argued for Ashkenazy\footnote{Vladimir Ashkenazy (b. 1937) is a Russian-born pianist and conductor, who won first prize at the 5th Chopin Competition in 1955.} to be first and Tanaka Kiyoko to be second, refused to sign the certificate, and walked out.

At the time, recording technology was not advanced, and the medium was records, so the sound quality was, to put it kindly, not good. Yet Tanaka Kiyoko's piano sounded astonishingly beautiful. The notes were lovely and mellow, like something a baby would instinctively want to put in its mouth.

``I feel like she plays as if she's praying---'' Yuzuki said, praising her piano highly. ``As if a pure, selfless prayer is sounding directly.''

I felt I understood her words, yet I didn't, so I asked:

``'Prayer' is the same as 'wishing,' isn't it? Since 'wishing' can't be separated from self-interest, can 'prayer' be separated from self-interest?''

Yuzuki looked at me with surprise and said:

``Didn't you offer a 'prayer' that was separated from self-interest?''

``When?''

``When you were collecting flowers for your mother.''

I was taken aback---but my thoughts took another unnecessary turn.

``That was to erase the pain inside me. That's still self-interest.''

``I don't think so---'' Yuzuki wore a gentle smile. ``I think it was also for your mother. For example, if you had continued praying for a thousand years, your body would have turned to dust, right? Dust without self or desire. But I think your strong prayer will continue for a thousand years. Just like Tanaka Kiyoko's performance reached us fifty years later. That prayer is already free from self-interest. Just as the meltwater from Mount Fuji is polished by the earth and becomes clear, prayer too is polished by time and flows clearly.''

I was dumbfounded. Even now, writing this story, I cannot help but be amazed by those words. It was hardly the kind of insight an ordinary eight-year-old girl could reach. I felt as if an angel or a Buddha resided within the girl before me.

Tanaka Kiyoko's piano began to sound different. Like the meltwater from Mount Fuji seeping into the body, it penetrated my heart with a nostalgic clarity---it was the sound of a pure prayer.


\subsection*{5}

Summer ended. The new semester began as if nothing had happened. I had felt as if I would be left behind forever in that summer when the world ended. Yet, somehow, I was carried along by time, living on shamefully. Like a cigarette butt carried by the wind. In the pool where cherry blossom petals had once floated, summer festival goldfish had taken up residence. Eventually, autumn leaves became their roof. The water grew muddy and green.

I began to feel a faint pain in Yuzuki's piano. At first, it was a slight discomfort, like biting into a grain of sand hidden in a clam's flesh. Even when I pointed it out, Yuzuki herself didn't notice, which I found strange. I could even point to the notes where the grains of sand were hidden, but Yuzuki would only tilt her head in puzzlement.

The sand gradually eroded Yuzuki's score. A desert-like loneliness covered her music.

``Lately, my mother has been telling me the same thing. But why? I don't understand it myself...''

Yuzuki lowered her eyes sadly and said that. A single tear fell from her left eye. At that spot, a bluish-purple bruise had formed. A bruise made by Ranako.

Worried that Yuzuki's piano tone was becoming increasingly murky, one Saturday I peered through the double windows into the soundproof room.---Yuzuki and Ranako were there. Ranako was a raging storm. She pulled Yuzuki's hair and slapped her cheeks. She was trying desperately to restore Yuzuki's sound, failing, and her anger was mounting. It was like adding new paint to revive a dirty color on a palette. The color only grew murkier and blacker, never becoming clean. This is pointless, I thought. It only hurts Yuzuki and accomplishes nothing. Isn't it you who's muddying Yuzuki's sound? Frustrated, feeling unbearably sorry for Yuzuki, unable to stand it any longer, I punched the window without thinking. With a crack---! the window unexpectedly shattered.

I came to my senses. Just before Ranako turned around, I quickly pulled my head back and hurriedly circled around to the left side of the house. The double window opened. Ranako leaned out and looked left and right.

``That's strange...''

I could hear my heart pounding. The double window closed. ...It seemed practice had resumed. Relieved, I looked up.

I froze. Someone was looking down at me from the second-floor balcony.

It was Yuzuki's father, Sosuke. He had both elbows resting on the second-floor railing, leaning forward as he looked down at me. I instinctively knew he must have seen everything. The sun was at its zenith, shining directly behind him, obscuring his expression in shadow.

I shifted my body little by little. Sosuke's face gradually came into view. He wore an unexpected expression. There was no anger in it. What was there was a deep, sorrowful emotion.

With the corners of his mouth turned down, Sosuke said nothing, simply looking at me.

I fled the scene. That expression of Sosuke's floated in the darkness of my heart and would not fade. Like the pale moon hanging in the winter sky.


\subsection*{6}

In mid-October, I found a pair of indoor shoes. They were lying abandoned in the gap between school buildings. The name ``Igarashi Yuzuki'' was written on them.

I made an indifferent face and returned them to Yuzuki. She also wore an indifferent face.

``Thank you. I was worried about losing them.''

Hmph, is that so? I wasn't so stupid as to accept that at face value.

``Yuzuki, are you being bullied?''

Yuzuki let out a breath. It was a blue-gray sigh that felt less like sadness and more like weariness that had piled up deep inside her body.

``...Or rather, I'm letting them bully me.''

``Letting them?''

Yuzuki smoothly took my confused hand.

``Let's talk at the secret base, shall we?''

If you walk a little east from the elementary school, the area becomes nothing but rice paddies, and houses become sparse. A small mountain draws near, and at its foot stand several mysterious abandoned old buildings. Among them is a small abandoned factory, which you can enter through a hole in the chain-link fence.

Inside, the factory is chillingly empty. In summer, the windows without glass cut out pieces of blue sky and let in cool breezes. Our secret base was, for some reason, an abandoned bus left on the premises. Its body was rusted a blue-iron color, with a rounded cuteness that somehow reminded me of a tin toy. One of the headlights was completely missing, giving it a goofy charm. I loved the face of this bus.

Inside, behind the driver's seat, there were two forward-facing single seats on each side. Further in were facing long benches. At the very back was a large forward-facing seat. The seat covers were green, torn in places with foam sticking out. The floor was wooden, and it made a pleasant tapping sound when you walked.

The interior was surprisingly clean for an abandoned bus. Yuzuki cleaned it regularly. The long benches were covered with orange geometric-patterned covers, with cushions for lying down.

It seemed this had been Yuzuki's secret base for a long time. She probably needed a personal space where Ranako's reach couldn't extend.

We sat side by side on the long bench. Yuzuki had avoided sitting face-to-face. There was a long silence.

``The trigger was music class---'' Yuzuki suddenly said. ``The teacher asked me to play the piano in front of everyone. After that, Sakamoto-kun and Aida-kun started talking to me all the time. I couldn't just treat them rudely, so I talked to them normally...''

``So the girls got jealous. Since you're both popular.''

Sakamoto was the one who had been hitting on Kobayashi Koyomi back in April, and almost had his nose twisted off by Yuzuki. The thought that he had now switched targets to Yuzuki and was causing trouble again made me irritated.

``Just twist off the bully's nose.''

``I don't know who's bullying me. My indoor shoes and pencil case get stolen and sabotaged without me knowing. All the girls in the class keep silent. If someone tells, they'll be the next target as a traitor.''

``...I hate that kind of thing. What about Kobayashi Koyomi? Isn't she on your side?''

``Koyocchan can't do anything, so she pretends not to see. But she's smart and kind, so she suffers from self-loathing and gets hurt too. Just for that, I can forgive Koyocchan.''

It took me a while to follow Yuzuki's thought process. She was far too sharp, too delicate, too bold, and too kind. It was hard to believe she was the same age as me.

``Lately---'' Yuzuki said in a slightly higher voice. ``I can understand the feelings of a girl in love. I can understand jealousy too. And then I thought, humans are all so weak. I decided to forgive everyone who bullies me. I decided to let them bully me. Everyone is really kind deep down, so I think they'll stop bullying eventually.''

``...And if they don't?''

``Then I'll twist their noses off properly.''

Yuzuki said this and smiled. I couldn't help but smile too. Then Yuzuki lay down and rested her head on my thigh. With a somewhat lonely voice, she said:

``...I'm amazing, aren't I?''

Confused by the sudden act, I still felt a great affection for Yuzuki.

``Yuzuki is amazing.''

``...Then pat my head. Say, 'Good girl, good girl.'''

``Eh...?''

After a moment's hesitation, I placed my palm on Yuzuki's head. Her shoulder twitched. I stroked her head. Her hair was soft and smooth. There was a gentle, sleeping breath---Yuzuki's peaceful breathing.


\subsection*{7}

However, the bullying toward Yuzuki only escalated instead of subsiding. It was as if she were attending school alongside a band of merciless bandits. All her belongings were stolen, thrown away, flushed, or burned.

Yuzuki couldn't twist any noses. She was destined to lose from the start. Just as those who fall in love are destined to lose, bullying is designed so that the gentler side loses.

One day after school, I heard three girls from Class 3 laughing and talking loudly.

``Serves her right. Yuzuki was getting too full of herself just because she can play the piano a little.''

It was a trash-talking session. I happened to be walking the same way home and was forced to listen to it endlessly.

``I don't know if she won some competition or something, but there are plenty of people who can play that well. I could easily surpass her with a little practice.''

``No way.''

The words slipped out. The three jumped and turned to look at me.

``You don't know how hard Yuzuki works. You don't know what feelings she puts into her playing. You're too stupid to understand how amazing Yuzuki is.''

The three were dumbfounded, their faces pale. One of them said, ``Let's go,'' and pulled the other two across the road. I could hear them whispering.

``Who is that guy...? What the hell is he...?''

``You know... that guy... the one who's been clinging to Yuzuki like a shadow lately...''

For some reason, being described as a ``shadow'' made me irrationally angry. I thought about yelling back, but suddenly felt empty and stopped. It felt as if the inside of my body had become hollow.


\subsection*{8}

With a sound that had become murky and murkier still, Yuzuki passed the regional round of the Chopin International Piano Competition in ASIA with a landslide victory. If you have overwhelming ability, you won't lose easily even when out of form. The national competition was in January. If she did well there, next would be the Asian competition.

But Yuzuki had stopped playing the piano altogether. It was the beginning of December. Light snow that wouldn't accumulate was falling. We spent our time holed up in the secret base. The air had grown quite cold, but I remember the bus being strangely warm. We bundled up in thick winter clothes and huddled together under puffy blankets. Yuzuki wore a pure white coat and knit hat, with a red scarf wrapped around her neck.

During this time, Yuzuki had a tendency to escape from reality. She wanted to distance herself from the piano, from her class, from her mother, and even from my room. For Yuzuki, there was probably no place gentler than the abandoned bus. She sought healing in the kindness that only ruins could offer.

The secret base was piled high with comics and books, and Yuzuki was always reading. Some were mine, but most were hers.

``You sure have a lot of money.''

When I said that casually one day, Yuzuki, without looking up from her new comic, replied:

``Dad keeps giving me money.''

``Keeps giving you money?''

``It's probably meant as a fine.''

I didn't fully understand what she meant, but those words sounded somehow sad. In my mind, I recalled Sosuke's face---like that pale moon---looking down at me from the balcony.

A fine---a penalty for what?

I thought about how my father continued to deposit living expenses for me. Perhaps that was also a kind of fine?

Before we knew it, the sun was setting. Winter days are short. Night was trying to drive us out of the secret base. Yuzuki said she didn't want to go home.

``I don't want to go back to that house. I don't want to go to school anymore. I don't want to play the piano. Everything has become meaningless. Yakumo-kun, let's stay here forever. Let's be together...''

Against the deep purple air between dusk and night, I let out a single white breath.

``...I can't do that.''

When I tried to stand up, Yuzuki clung to my arm, wrapping her palm around my left hand. Her hand was as cold as a piano keyboard.

My heart skipped. Yuzuki rested her head on my shoulder. A sweet fragrance drifted.

It felt as though a very long time had passed. Yet the color of the sky hadn't changed, as if we were trapped between day and night. Yuzuki asked in a whisper:

``...Yakumo-kun... do you have someone you like...?''

I could feel my heart pounding. My face was probably red too.

``...I'm not sure.''

``Won't you tell me?''

``I don't want to tell.''

``...Yakumo-kun might have someone he likes, or he might not. ...Then, if you do have someone you like... is it someone other than me?''

Before I knew it, her cold hand had grown warm. I said:

``...If I did, it wouldn't be someone other than you.''

Yuzuki chuckled softly.

``Me too... if I had someone I liked, it wouldn't be someone other than you.''

I looked at Yuzuki without thinking. She looked back at me, wearing a mischievous, almost impish smile. Even in the dimness, I could see her white cheeks faintly flushed.


\subsection*{9}

The next day, my head was foggy all day. I hadn't slept at all the night before. Even in class, my mind was elsewhere, filled with the colors of last night's sky.---On the way home yesterday, Yuzuki had said:

``...Hey, Yakumo-kun, let's go far away together. First to Lake Inawashiro,\footnote{Lake Inawashiro (猪苗代湖) is a large lake in central Fukushima Prefecture, the fourth largest in Japan by surface area. It is a popular tourist destination.} and then, as we please, as far away as possible. And then we'll live happily for months.''

I was shocked and replied: ``That's impossible. First of all, we don't have money.''

``It's okay. I'll steal it.''

Startled, I looked at Yuzuki. It was already dark, and I couldn't see her expression.

``Dad hides cash in the back of his closet. About 1.5 million yen. When I tell him I need money, he takes one bill from there and gives it to me. ...So that's a fine. Even if I take it all, Dad probably won't get angry.''

``1.5 million yen...'' I thought that would be enough to live on for a while. Yuzuki said:

``It's okay if we don't get away completely. I really just want to go far away from everything for a little while... Please, come with me? I'll be waiting at the secret base tomorrow---''

Yuzuki waved goodbye and went inside. I thought she would probably be yelled at by Ranako for coming home late.

---During class, I was absentmindedly thinking. Should I go with Yuzuki? If we left our lives for months, various problems would arise. Most importantly, Yuzuki wouldn't be able to compete. Was that really okay? Was it okay to waste all her efforts so far?

During lunch break, I was walking around in a daze when, unexpectedly, I ran into Yuzuki. Her almond-shaped eyes were wide open. Her eyes were captivating. Yuzuki said nothing and walked past me.

---In that instant. Casually, she touched my hand with her fingertip. My heart leapt, and I turned around without thinking. Yuzuki didn't look back. Suddenly, I felt a gaze. Sakamoto was standing there with a stunned look on his face. I looked away and hurried off.

---Before I knew it, school was over. I went straight home and packed for the trip. I stuffed a change of clothes and a toothbrush set into my backpack. Even as I prepared, I kept hesitating. My unease grew larger. I knew that Yuzuki needed to ``run away from home,'' but my thoughts entered a labyrinth, caught by an unknown dark feeling. It felt as if I were leading Yuzuki in the wrong direction, trying to ruin her life. Somewhere, I could hear the whispers of the girls from Class 3:

``That guy... the one who's been clinging to Yuzuki like a shadow lately...''

After all, wasn't I just Yuzuki's shadow? Wasn't I, a mere shadow, trying to drag Yuzuki from the spotlight into the darkness of the wings?

Half of me was my father, and the other half was my mother. Half was shadow, half was salt. I felt as if I could never become a full-fledged person. Someday, I might find someone annoying, like my father, and make that person unhappy.

I must be an incomplete person. I didn't have the right to run away with Yuzuki.

These thoughts circled endlessly. It was irrational. But at that time, I didn't have such a calm perspective. Like a black hole warping the surrounding space, the pitch-black hole in my heart warped every thought.

With my backpack before me, I sat on the floor. Time passed. Dusk came, and night fell. Snow fluttered outside the window. Was Yuzuki still waiting for me? Was she freezing alone? I wavered and wavered, and in the end, I didn't go to the meeting place.

Unable to sleep, I remained frozen like a stone statue until morning.

\subsection*{10}

Yuzuki did not blame me. In fact, she never brought up the subject of ``running away'' again. As if it had never existed in the first place.

However, something had decisively changed between Yuzuki and me. It was a subtle, yet coldly precise change---like zero turning into one.

Yuzuki stopped playing the piano in front of me. She also stopped going to school. She shut herself in at home and continued practicing until the competition. Alone. Or while being yelled at by Ranako. My role was to be Yuzuki's pressure-release valve---on weekdays, I would just chat with her for about an hour about trivial things. A ``hole'' for releasing pressure is a role quite fitting for a ``shadow.'' But it was a position I had chosen for myself.

Yuzuki, as expected, won the national championship. And then came the Asian competition---Yuzuki shone at the top of the elementary third and fourth grade division.\footnote{The Chopin International Piano Competition in ASIA is held in multiple age categories; the elementary third and fourth grade division is for children aged approximately 8 to 10.} At the Chopin International Piano Competition in ASIA, a commemorative album featuring the performances of gold prize winners is released. The CD containing Yuzuki's performance came out in September.

I was excited, thinking I would finally hear Yuzuki's performance after so long, and I put on the CD. After the ``runaway'' incident, Yuzuki had never played the piano in front of me even once.

From the speakers, Yuzuki's performance flowed out. It was Chopin's \textit{Barcarolle}.

---I was astonished. The sound that had been so murky was now as clear as crystal. And yet the performance carried a profound sorrow and depth. Hearing this, no one would imagine that the performer was a third-grade girl.

But somehow, I felt that this was not Yuzuki's sound.

There was a boat, and there was a sea---but this time, there was no sky. Yuzuki's performance could not reach heaven.

It was a sound that had stopped praying.

\fullpageimage[]{7.jpg}
% ==============================
%        Chapter 3
% ==============================
\clearpage
\thispagestyle{plain}
\phantomsection
\begin{center}
    \large\bfseries Chapter 3
\end{center}
\addcontentsline{toc}{chapter}{Chapter 3}
\label{sec:ch3}
\vspace{1cm}

\subsection*{1}

We became middle school students. Our elementary school classmates became our middle school classmates.

In the fourth grade of elementary school, Yuzuki won the gold prize, first place, in the Concerto B category of the Chopin International Piano Competition in ASIA---a category with no age restrictions---making her the youngest winner in history. Furthermore, the following year, she won the gold prize, first place, in the Concerto C category, again as the youngest winner in history. At the age of eleven, she made her professional debut as a classical pianist.

``The Miraculous Genius Beautiful Girl Pianist'' became Yuzuki's catchphrase. She was tagged with this label that she ``hated so much it made her want to vomit,'' and she appeared on television more and more. Because she was intelligent and had strong nerves, she was also invited to programs unrelated to piano. She was assigned a manager who handled her schedule and other matters.

Kitajo Takashi was that manager---a handsome man with a high nose and slender silver-rimmed glasses. He treated Yuzuki like a princess. He would open every door for her, and he never missed using breath freshener before speaking to her so that he could address her with a refreshing breath. He always wore a confident smile, as if to say that he himself was the prince by her side.

He always carried a heavy-looking camera around his neck and took photos whenever he had the chance. It was his dream to someday turn photography into some form of career.

I have a photo he took. In it, Yuzuki and I are standing side by side. I'm making a sour face. I disliked Kitajo. The way he looked at Yuzuki through his camera, with that entranced, disgusting expression, made me feel sick.

I knew very well that what he felt was a mixture of adoration and love. I knew it because countless boys had directed similar gazes toward Yuzuki.

As Yuzuki grew, she became astonishingly beautiful and charming. Even a boring classroom, with her in it, became as brilliant as if adorned with white flowers. Boys from other classes and upper grades crowded around her door just to see her. Sometimes students from other schools would slip in. Yuzuki would narrow her almond-shaped eyes and shoot them a look of absolute zero. She was still not one to seek attention or show off.


\subsection*{2}

At Sakuranoshita Middle School, students were required to join some kind of club activity.

I had no choice but to join the baseball team, because Shimizu was there.

Shimizu looked easygoing, but he had astonishing talent for baseball. In elementary school, he had been the cleanup batter and pitcher, and apparently he performed remarkably both at bat and on the mound. As a first-year middle school student, he already stood nearly 180 cm tall, bigger than anyone, including the third-years. Before long, he pushed aside the third-years and became a regular.---With such conspicuousness, one might expect jealousy, but on the contrary, he was adored by everyone. Shimizu had a mysterious charm, and everyone loved him.

He was always smiling, cheerful, innocent, kind, and a little bit of a fool. Yet in a pinch, he was reliable. He loved baseball and always practiced with excitement. Just watching him was fun, and seeing him perform well felt great. No one could help but feel affection for him. The number of first-year club members was nearly double that of previous years, and I think Shimizu's charm had a lot to do with it. As long as Shimizu was around, there was a sense of security that no one would ever be left out. When Shimizu practiced happily, even I felt happy and never found practice difficult. Shimizu was like the sun. For some reason, he took a liking to me, who was like a shadow, and thanks to that, I never struggled to make friends. After practice, we would all go home laughing loudly together. When Shimizu got his first failing grade and faced detention, all the first-year baseball players went to his house to study together. I was able to live days so happy they felt almost too good for someone like me.

On the other hand, I had grown distant from Yuzuki. We were in different classes, and she was extraordinarily busy. Between concerts and television appearances, her schedule alone was dizzyingly packed, and she seemed to practice the piano fiercely without missing a day. She probably had no time to meet me.

Even without that, my relationship with Yuzuki was delicate. Since the ``runaway'' incident, she had become someone close yet distant. Even though we laughed together side by side, I felt her heart was far away. Though I could reach out and touch her, she was infinitely distant. Just when I thought I had touched her, she would vanish. As her name suggested, like the moon reflected on a lake's surface.

But I was careless. I naively believed that just as I thought of Yuzuki as special, she also thought of me as special.

One July evening, as I was walking home laughing with the baseball team as usual, Aida suddenly exclaimed. I followed his gaze and was stunned.

Yuzuki and the baseball team captain were walking side by side.

Both of them were slightly looking down, exuding a strange, bittersweet atmosphere. Aida let out a sound as if something had pierced his solar plexus---``Uwaaa... aah...''---and said:

``Is Igarashi going out with the captain...?''

Come to think of it, Aida was the one whose bullying had started Yuzuki's troubles. I hadn't imagined he still liked her. He clutched his head and cried ``Uwaaaa...'' with tears even forming in his eyes.

``Are they already kissing...? Like, doing perverted stuff...?''

Aida, who was popular among girls, showed unexpectedly pure feelings. Looking back, it's amusing, but at the time I couldn't laugh at all. I felt as if I had been hit in the head with a bat.

I liked Yuzuki more than I had realized.


\subsection*{3}

The baseball team captain's name was Roppongi Satoshi.

I began to observe him whenever I had the chance. It was reconnaissance of the enemy. I grew anxious, because the more I looked, the better a man he seemed. He was tall, handsome, good at baseball, and had excellent grades.

``Next month's goal is to grow up and become 'Senbon-gi' (Thousand Trees)!''

He would make such jokes and make everyone laugh---a true entertainer.

``I'll power up 167 times!''

With the cleverness to casually round up to the nearest decimal.

Everyone laughed heartily, while Aida showed pure jealousy.

``'Senbon-gi' my foot. I'll turn you into 'Sanbon-gi' (Three Trees)!...''

He said such foolish things, and I was just as foolish, mentally demoting Roppongi-senpai to ``Moto-senpai'' (Former Senior). It was deforestation. Desperate to somehow defeat Roppongi-senpai, I threw myself into baseball. I went to Round One\footnote{Round One is a Japanese chain of sports and entertainment complexes, featuring batting cages, bowling, arcade games, and other facilities.} at the newly built East Shopping Center in Koriyama Station and honed my skills on the batting machines. My hands became covered in blisters.

During that time, I witnessed Yuzuki and Roppongi-senpai walking home together several times. Each time, Aida would make strange sounds like ``Ah-han'' or ``Uh-hun'' and grieve.

After two weeks of intensive training, the opportunity to face Roppongi-senpai finally came. Even though it was just practice, I was unusually serious. I swung the bat with all my might, fouled twice, and the count was two balls and two strikes. I sharply hit a pitch that came in slightly outside---a single. I felt like pumping my fist, but I played it cool. I was in adolescence.

The next batter was Aida. From the moment he stepped into the batter's box, he had the face of an angry demon and swung at a ball out of the zone, making contact. Another single. Aida let out a cry of joy like ``Pao-on!''---like an excited African elephant.

``Not bad, first-years...''

Roppongi-senpai smiled wryly and wiped the sweat from his brow. Then Shimizu stepped into the batter's box. His massive frame made the bat look like an ice cream stick. He swayed his body happily.

Ah, he's going to hit it, I thought.

\textit{Ka---n!} A pleasant sound rang out. The ball rose incredibly high.

``Uwa-hahaha---!'' Shimizu laughed. He had a habit of laughing when he was sure of a home run. His laugh was indescribably pleasant, and among us he was called the ``Great Demon King.''

Overjoyed, I also laughed, ``Uwa-haha!'' Aida laughed too. The three of us, laughing, rounded the bases and stepped on home plate. Roppongi-senpai made a face that said ``Oh dear.''

I lingered in the afterglow of victory, watching the other first-years bat.

---And gradually, I came back to my senses.

Even if I had beaten Roppongi-senpai a little in baseball, so what...?

I looked beside me; Aida had been grinning all along.


\subsection*{4}

At the end of July---the closing ceremony ended in the morning, and the afternoon was completely free.

I wandered slowly around the school building. The sound of the brass band practicing came from everywhere. I was thinking about Yuzuki. What was she doing right now, and where...?

When I went around to the back of the school building, I heard a voice: ``Oh.''

There, Yuzuki was sitting with her back against the school wall.

``Oh.'' Caught off guard by this unexpected encounter, I didn't know what to do.

``L-long time no see, then...''

``Wait---'' Yuzuki grabbed the hem of my shirt. ``Why are you leaving so soon? Let's catch up, it's been a while.''

She made me sit beside her. My heart was pounding like crazy. Even though we had been together so much in the past, it felt like we were meeting for the first time. I glanced at Yuzuki's profile. Had she always been this beautiful?

``How have you been?''

I answered hastily: ``---I've been fine.''

``Are you eating properly?''

Her questions sounded like a mother worried about her son living alone.

``Yuzuki, you seem busy---are you okay?''

Yuzuki smirked mischievously, leaning forward, and said:

``For you, a mere Yakumo, to worry about me---that's ten years too early.''

``What's that supposed to mean...?''

I turned away, pretending to be sulking---but in truth, I wanted to hide my red cheeks, from the sight of her white thighs showing beneath her skirt and her adorable smile.

A brief silence fell. I ended up asking what I had been meaning not to ask.

``Yuzuki, are you going out with Roppongi-senpai?''

She chuckled, ``Hehe.'' With a triumphant, impish smile, she said:

``Does it bother you?''

``...Not really.''

``Liar. You're bothered, aren't you?''

I closed my mouth. In front of my eyes, the clear black shadows cast by the summer sunlight swayed lazily.

Yuzuki, however, murmured with satisfaction:

``Tomorrow, summer vacation starts.''

Then, peering at me again, she said:

``Hey, are you free tomorrow? I have a favor to ask.''

``A favor?''

``Come over to my place tomorrow.''


\subsection*{5}

The next day, I visited Yuzuki's house. With my heart pounding, I rang the doorbell. The door opened immediately. Yuzuki appeared like a white flower. She was wearing a white dress.

``Welcome---'' She smiled and let me in. It had been a full year since I'd come to the Igarashi house. I felt a nostalgic scent in the air.

As I was lost in sentiment, I suddenly ran into an unexpected person in the living room and nearly let out a gasp. The other person seemed to feel the same. Yuzuki's manager, Kitajo, who was lounging as if he owned the place, twitched his nose and gave a stiff smile when he saw me.

``Hey, long time no see, Yakumo-kun.''

Yuzuki was smiling calmly. Seeing her face, I intuited: she probably had no choice but to be alone with Kitajo today. She didn't want that, so she had summoned me.

\emph{You little...} I thought, but I put on a fake smile and sat down at the table.

``Since we're all here, let's take a commemorative photo.''

I didn't know what was so commemorative, but Kitajo was already taking pictures with the camera hanging from his neck. No doubt he would later cut out the part with me and hang Yuzuki's photo on the wall.

After chatting about various things, we ate the shortcake Kitajo had bought. Of course, there were only two cakes, so Yuzuki gave me half.

``Here, Yakumo-kun, you can have the strawberry too---''

At that moment, Kitajo's nose twitched again.

``Then I'll give my strawberry to Yuzuki-chan.''

``Ah, I don't really feel like eating strawberries today, so it's okay.''

``What a coincidence. I feel the same, so don't hold back.''

``Oh, then I'll give it to Yakumo-kun.''

My unbalanced, thin cake now had two strawberries on it. Seeing Kitajo's resentful face, I had a hard time suppressing my laughter. He looked as if he wanted to say ``Give me back my strawberry!''

Then Kitajo suddenly stood up, went to the next room, and put on a CD. Yuzuki's performance began to play. At that moment, Yuzuki shot me a look. Kitajo wore what he probably thought was a clever smile and sat down leisurely. Then, without warning, he began enthusiastically explaining to me what was so wonderful about Yuzuki's performance.

Yuzuki's face turned red. Kitajo, seeing her reaction from the corner of his eye, grew even more excited. Yuzuki's face turned red to her ears. Kitajo probably thought she was embarrassed, but in reality, she was furious.

Thinking Yuzuki might actually lash out at Kitajo, I quickly stood up.

I pretended to go to the bathroom and took another route to the room with the stereo. Then, at the moment the track changed, I secretly swapped the CD and returned to the living room with an innocent look. During my absence, Kitajo had apparently said something more, and Yuzuki's anger had reached its peak.

Just then, the music began to play---

Yuzuki's expression relaxed, and she looked at me. I gave her a slight smile from the corner of my mouth.

Kitajo, completely unaware that the performer had changed, continued praising in the same tone.

Yuzuki laughed with a snort and quickly looked down to hide it.

Kitajo seemed to misunderstand something and praised even more extravagantly.

``This is a sound only Yuzuki-chan can produce!''

I barely held back my laughter and said:

``You're absolutely right.''

Yuzuki snorted again and covered it by pretending to sneeze.

Whenever Kitajo said something ridiculous, Yuzuki and I were enveloped in a strange, sweet atmosphere, like partners in crime.---Suddenly, under the table, my right little finger and Yuzuki's left little finger touched. I didn't know if it was an accident or intentional.

For some reason, I didn't want to confirm. The ambiguity was fine. We didn't separate, nor did we interlock; we just kept our little fingers touching.

I silently thanked the female pianist who had played such wonderful music in Yuzuki's stead.

Thank you, Martha Argerich.\footnote{Martha Argerich (b. 1941) is an Argentine classical pianist, widely regarded as one of the greatest pianists of the 20th and 21st centuries, known for her virtuosity and passionate interpretations.}


\subsection*{6}

Since that incident, the distance between Yuzuki and me had somehow grown closer again. Sometimes I visited her house, and sometimes she visited mine. Still, she never played the piano in front of me.

In August, we went together to the ``Uname Festival''\footnote{The Uname Festival (うねめ祭り) is a summer festival held in Koriyama City, Fukushima Prefecture, featuring traditional dancing, parades, and fireworks.} held in front of Koriyama Station.

As I waited in front of the Igarashi house, the front door opened and Yuzuki appeared.

I was captivated. Yuzuki was wearing a yukata---a blue yukata with a pattern of bellflowers,\footnote{Bellflower (姫射干, \textit{Belamcanda chinensis}) is a flowering plant with orange-red spotted petals, often used in Japanese fabric patterns.} and the sash was a vivid red. On her tied-up black hair, cute white flower hairpins dotted about.

It almost felt as if the darkness around us brightened the moment she appeared.

Yuzuki saw me and smiled.

``Well? Does it suit me?''

She turned around to show me. Her geta\footnote{Geta are traditional Japanese wooden sandals, often worn with yukata during summer festivals.} made a cool sound, and a pleasant fragrance drifted, making it hard for me to look at her properly.

Just then, Sosuke appeared from the entrance. I was startled. The last time I had faced him was when I broke the double window. ``I'll drive you to the station,'' he said, walking past me toward the garage. Yuzuki walked on as if nothing was amiss. Though confused, I followed.

---The white coupe began to move. The townscape flowed slowly by.

I didn't know what to say, so I stayed quiet. Yuzuki was as still as a hina doll.\footnote{Hina dolls are decorative dolls displayed during Japan's Hinamatsuri (Girls' Day) festival. The comparison suggests quiet, formal stillness.}

Suddenly, Sosuke spoke.

``It's been a while, Yakumo-kun. I always hear about you from Yuzuki.''

``...It's been a while.''

I glanced at Yuzuki. She was looking out the window.

``I hear you joined the baseball club.''

We talked about club activities. I briefly mentioned Roppongi-senpai, but Yuzuki's expression didn't change.

``Sosuke-san, what club were you in during middle school?''

``I didn't do any club activities. I played the piano. I went to music college, and there I met Ranako. ...I ended up giving up, though. So Yuzuki's talent comes from Ranako. Her looks too. I wonder what part of me Yuzuki resembles at all.''

His tone was somewhat lonely, as if saying not a single part. I couldn't help but say:

``Isn't it her kindness?''

Sosuke looked surprised. Our eyes met through the rearview mirror. Something seemed to flicker deep in his eyes---a pure emotion like a child's. But he quickly looked away and said in a bittersweet voice:

``...I'm not such a kind person...''

``Dad is kind.''

Yuzuki said it, still looking out the window. Sosuke kept his gaze straight ahead, opening his mouth and closing it again. Even though they were in the same car, they seemed like two people on opposite lanes.

``Fine''---the word rose from the depths of my memory. Like a shell buried in mud, a foreign object that never disappeared.


\subsection*{7}

The station square was overflowing with people.

Yuzuki and I strolled through the food stalls, eating cotton candy, candied apples, and kebab. I, who lived alone and managed my own finances, was a bit put off by the festival's inflated prices, but Yuzuki happily kept buying things and laughing adorably.

Festival music played without pause. The people flowing like a gentle river stepped and turned, unknowingly syncing with the rhythm of the drums. Hairpins swayed; goldfish spun. It seemed chaotic yet gently rhythmic. Bright rhythms swelled like waves here and there, then softly dissolved.

I felt as if paint had been splattered inside my head. It always happens in crowds. Excess information and emotions flood in, making me dizzy. It felt like being car sick in an unfamiliar land.

A line of dancers passed by---

Suddenly, I heard a voice: ``Ah.'' I came to my senses and looked toward the sound.

It was Roppongi-senpai. He was with two other third-years. They seemed to have come to the festival too.

But something was off. Senpai's eyes were wide open as he looked back and forth between Yuzuki and me.

``Ah, Igarashi---'' Roppongi-senpai said. The five of us ended up stopping in the middle of the crowd. Passersby steered around us with annoyed looks. Senpai scratched his cheek with his index finger.

``Igarashi, didn't you have something to do today...?''

I looked at Yuzuki. She remained calm.

``This is what I had to do. Yakumo-kun invited me first.''

I had a feeling that it was Yuzuki, not me, who had first suggested going to the festival. Senpai still seemed to be catching up. ``Oh, I see...''

``Well then, see you. Let's go, Yakumo-kun.''

Yuzuki's geta clattered as she walked away briskly. I bowed to the still-stunned seniors and chased after her.

``Yuzuki---'' Finally catching up, I asked: ``Weren't you going out with Roppongi-senpai...?''

``We weren't going out,'' Yuzuki said in a somewhat stiff voice. ``He kept asking, so I just walked home with him. I thought it would be rude to turn him down flat. He's a nice person.''

``...But isn't it more cruel to lead him on like that?''

Yuzuki spun around, her voice a little irritated.

``Then would you rather I had gone out with Roppongi-san?''

``I didn't say that.''

A reddish-blue color tinged the corner of Yuzuki's eyes. She turned around again and walked on.


\subsection*{8}

We took a bus bound for Miharu and Funabiki, got off at Mizuana about fifteen minutes later, and walked another five minutes to reach the Fukuyama Yume Fireworks venue.\footnote{The Fukuyama Yume Fireworks is a summer fireworks display held along the Abukuma River in Koriyama City.} The riverbed of the Abukuma River was packed with people. Since all the good spots were taken, we walked along the river and spread a plastic sheet a little further away. There was still some time before the launch. We still carried the awkward atmosphere from the Roppongi incident. Suddenly, Yuzuki spoke.

``...Hey, do you know the origin of the Uname Festival?''

Then she told me the ``Uneme Legend.''\footnote{The Uneme Legend (采女伝説) is a tragic love story from the Heian period, associated with Koriyama City and the famous Uneme legend of Nara's Sarusawa Pond. It tells of a woman named Haruhime who was offered to the Emperor but escaped and returned home, only to find her lover had died.}

---About thirteen hundred years ago. The village of Asaka in Mutsu Province, the predecessor of Koriyama City, had suffered continuous cold damage and could not even offer tribute to the imperial court. The villagers pleaded their plight to the imperial inspector, Katsuraki-no-Okimi,\footnote{Katsuraki-no-Okimi (葛城王) is a historical figure, a prince and official of the Nara period.} who had come from the capital of Nara, and begged to be exempted from tribute. But the prince would not listen. That night, a banquet was held in his honor, but the prince was displeased, feeling the hospitality was lacking.

Then, an elegant and beautiful young woman---Haruhime---appeared. She held water in her right hand and sake in her left. She tapped the prince's knee and recited a poem:

*Asaka-yama kage sae miyuru yama-no-ino no asaki kokoro o waga omowanaku ni*

Her poem meant: ``My hospitality is not as shallow as the well that reflects even the shadow of Mount Asaka; I do not hold such shallow feelings toward you.'' The prince must have seen the ``mountain shadow'' and ``shallow heart'' in the water of her right hand, and the ``true heart'' in the sake of her left. He gladly drank the ``true heart'' she offered. His mood restored, the prince agreed to exempt the village from tribute for three years, on the condition that Haruhime be offered to the Emperor as an uneme (court lady).

Haruhime had a betrothed lover named Jiro, but she tearfully parted from him.

In the capital, Haruhime received the Emperor's favor. But her longing for Jiro only grew, and one moonlit autumn evening, she slipped away from a banquet, ran to the shore of Sarusawa Pond, hung her clothes on a willow branch, and fled home, feigning a drowning.

It must have been a grueling journey. Exhausted in body and spirit, Haruhime finally arrived home, only to be met with a sad reality. Jiro, in his grief over losing Haruhime, had already thrown himself into the clear waters of Yamanoi Well.

On a snowy night, Haruhime, too, finally sank into the same clear waters of Yamanoi Well.

When spring came and the snow melted, around the clear waters of Yamanoi Well bloomed countless pale purple flowers of unknown name. Those flowers, as if the crystallization of the two lovers' love reborn, came to be called ``Asaka no Hanakatsumi''\footnote{``Asaka no Hanakatsumi'' (安積の花かつみ) is a local name for a pale purple flower that blooms in the Koriyama area, associated with the Uneme Legend.}---

This legend has another version from the Nara perspective, in which Haruhime drowned herself in Sarusawa Pond out of grief over losing the Emperor's favor. Yuzuki told the Koriyama version, a tragic love story, and even that was abridged.

``That's a sad story,'' I said.

Yuzuki sighed.

``What do you think is the lesson we can learn from this story?''

``A lesson---?'' I thought for a moment and said: ``Love is beautiful?''

``How banal,'' Yuzuki said, a little exasperated. Then, after a pause: ``The lesson is, 'Women must use all their wisdom.' Women are weak, so without using all their wits and guile, they cannot fulfill their wishes. ...Not even one single love.''

``Is that so?''

``You wouldn't understand, Yakumo-kun. You're still a child---''

Her words echoed with a certain loneliness.

Just then, a large firework burst in the night sky. The surface of the Abukuma River reflected the light like a mirror.

It was a beautiful summer night.

\fullpageimage[]{8.jpg}
% ==============================
%        Chapter 4
% ==============================
\clearpage
\thispagestyle{plain}
\phantomsection
\begin{center}
    \large\bfseries Chapter 4
\end{center}
\addcontentsline{toc}{chapter}{Chapter 4}
\label{sec:ch4}
\vspace{1cm}

\subsection*{1}

A period of peace continued for a while. I became a second-year middle school student, and Roppongi-senpai graduated. In the end, after that Uname Festival, Yuzuki and Roppongi-senpai were never seen together again.

Since I had been so eager to beat Senpai, my batting improved and I became a regular. What I was even better at was stealing bases. Somehow, I could read the pitcher's carelessness---or rather, the blind spots in their awareness. Shimizu, of course, and Aida also became power hitters and made the regular lineup.

My relationship with Yuzuki continued at a certain distance---close but not too close. It felt like walking along a long path divided by tall fences, chatting side by side without holding hands. She would hide somewhere only when playing the piano, and after that, we'd start walking together again...

What should I call this kind of relationship?

---But those lukewarm days eventually came to an end. I don't know how to describe the temperature of the days that followed. It was March, just before I became a third-year middle school student.

The room shook.

Everything fell from every shelf. Vases shattered. Books flew through the air. The lights went out. Dressers slid across the floor. The entire apartment building creaked and groaned. Screams came from somewhere...

It was March 11. The Great East Japan Earthquake.\footnote{The Great East Japan Earthquake (東日本大震災) struck on March 11, 2011, with a magnitude of 9.0, triggering a massive tsunami and the Fukushima Daiichi nuclear disaster.}

It wasn't just the room---all of Japan was shaking.

When the shaking finally subsided, I crawled out from under the bed covered in dust, sat stunned for a moment, and quickly sent a message to Yuzuki to check on her. At the time, I was using what's called a ``garakei'' (feature phone).\footnote{Garakei (ガラケー), short for ``Galapagos mobile phone,'' refers to Japanese feature phones that were popular before smartphones became widespread.}

A large crack ran through the wall---like a black river flowing down the center of the west wall. I couldn't take my eyes off that crack. It felt as if forgotten dark emotions were seeping out through it.

A message arrived. But it wasn't from Yuzuki. It was from my father.

``Are you okay?''

I remembered that I had a father.

I closed my phone once, thought for a moment, and replied:

``I'm fine.''

That was the end of the exchange.

Finally, a message came from Yuzuki. She had been in Tokyo for a concert, and she was worried about me. I replied again:

``I'm fine.''

It was 3:00 PM. I went outside. New cracks had appeared all over the apartment stairs. People were running around in confusion on the streets. The sky, which had been a quiet overcast until just moments ago, had turned into a blizzard right after the earthquake.

I remembered the summer of third grade, when my mother turned to salt. The giant cumulonimbus clouds beyond Mount Adatara, like a bomb had dropped, like the world had ended---that hot summer...

I might have come to the other side of that mountain. It was cold, gray, and snowing.


\subsection*{2}

The damage in Koriyama was relatively light compared to the coastal areas. In Sakuranoshita, where I lived, internet and phone connections were temporarily lost, but that was about it. Even within the same city, the damage varied by district---some areas lost infrastructure entirely, while others could continue living almost as before. Due to disrupted distribution, food disappeared from store shelves. Hospitals became evacuation centers, overflowing with people.

On the other hand, the scenes shown on television were truly devastating.

The massive tsunami struck mainly the coastal areas of Iwate, Miyagi, and Fukushima, causing enormous damage. Furthermore, the tsunami struck the Tokyo Electric Power Company's Fukushima Daiichi Nuclear Power Plant,\footnote{The Fukushima Daiichi Nuclear Power Plant (福島第一原子力発電所) suffered meltdowns in reactors 1, 2, and 3 following the 2011 tsunami, leading to the release of radioactive materials and a major nuclear disaster.} which lost power and cooling capabilities, leading to meltdowns in reactors 1, 2, and 3, releasing large amounts of radioactive material.

I stayed in touch with friends and Yuzuki while watching the news nonstop.

When the internet became available again, I watched videos uploaded by ordinary people. Brown, turbid seawater surged in, sweeping away buildings and cars. The cameraperson's shocked voice, people's screams...

I ran to the toilet---and vomited.

The sour vomit of melted frozen food burned my throat. Tears overflowed uncontrollably. My mind was in chaos; I was dizzy and couldn't stand.

At that time, many people across Japan apparently fell ill from watching footage of the disaster zone. It's called ``compassion fatigue''---exhaustion from empathizing too deeply with the suffering and pain of others.

In my case, that special phantom-limb-like pain was added to it. The land left devastated by the tsunami, reduced to rubble, was a vast emptiness. A void of unprecedented scale... After throwing up everything in my stomach, I felt as if my body had become transparent. I had no strength anywhere, and I just sat there blankly in front of the vomit-stained toilet.


\subsection*{3}

Three days after the earthquake, the doorbell rang.

I had been bedridden, just like when my mother died, and groggily opened the door.

A shadow stood there.

The last time I'd seen him was that summer in third grade. The day I told him to go to hell... Over five years had passed. My father looked exactly the same, still wearing black clothes.

``You look well, kiddo.''

He said this with a grin. My feelings were complicated. It was like rediscovering an old, ugly scar I'd almost forgotten...

I neither rejected him nor invited him in; I just stood there.

``Mind if I come in?''

The shadow entered without waiting. He set down the bags of groceries he was carrying, then went back down the apartment stairs to get more. I returned to my bed, closed my eyes, and listened to his footsteps.

---For the next three days, I lived with the shadow.

I slept and woke repeatedly, day and night, spending most of my time in bed. Whenever I woke up, I could sense my father's presence. He was always clicking away at a keyboard, writing his novels. He had seen that watching the news worsened my condition, so he unplugged the TV antenna. Instead, he played movies nonstop. He liked Disney movies. Without much conversation, we sat side by side eating cup noodles and watching Winnie the Pooh and Aladdin. It was a strange time---like being inside a strange dream from long ago.

My father was still a shadow---but somehow, I felt a sense of closeness and reassurance. And then, before I knew it, the shadow had disappeared.


\subsection*{4}

Even when summer came, my health was still poor. I began to miss school frequently. Shimizu was very worried, but I didn't know how to explain the cause of my condition.

The radioactive materials released from the Fukushima Daiichi Nuclear Plant had become a national issue. Every day, I saw dark news: farmers who lost their crops taking their own lives, children who evacuated to other prefectures being bullied there. The term ``reputation damage'' became well-known.\footnote{``Reputation damage'' (風評被害) refers to economic and social harm caused by unfounded rumors or fear, particularly affecting products from Fukushima after the nuclear disaster.} Agricultural, fishery, and livestock products—even those properly tested and proven safe—wouldn't sell just because they were from Fukushima. Rumors flew everywhere, and it became impossible to tell what was true and what was false.

I started picking flowers again, just like when my mother lost her limbs. But this time, the emptiness was too vast. No matter how many flowers I had, it wasn't enough.

At a loss, I escaped into the world of video games. An online game. There were enemies that dropped ``flowers'' when defeated, and I collected them endlessly. \textit{Sukatan-bogubogu}---the enemy made a sound and dropped a ``flower.'' \textit{Sukatan-bogubogu-chararira, sukatan-bogubogu-chararira, sukatan-tata-bobobo-bogubogu-wellington-chararira}---like beating a distorted, cute drum, it was a rhythmic, endless, monotonous task. But strangely, the pain subsided.

It didn't have to be real flowers. Because it was a conceptual act.

Meanwhile, Yuzuki had been irritable ever since the earthquake. There were too many things in the world that displeased her.

We started sneaking out of our houses at midnight to take walks. The midnight streets were quiet, as if nothing was wrong. The playground equipment looked like large animals sleeping peacefully. But in reality, radioactive materials were floating in the air, spreading contamination like a slow poison... It was very strange and frustrating.

---Why had Yuzuki been so angry that day?

I think it was because some celebrity had said that people still living in Fukushima were negligent.

``Leave now.''

``Even knowing the danger of radiation, you still live in Fukushima?''

``Don't give your children thyroid cancer because of your selfishness.''

``Shut up, you idiot---that kind of feeling,'' Yuzuki cursed with unusual vulgarity. ``At your age, you can't even make a calm analysis; you're just driven by emotion and saying things that only increase the damage. Where do you get off talking like that? Everyone is enduring things that aren't their fault, so why are people in safe places getting upset and saying whatever they want? Everyone carries unresolved feelings, so why can't they have even that much imagination...''

Before I knew it, Yuzuki was crying. She cried a lot during this time. She fought with Ranako every day, and her emotions were always on edge.

I looked up at the summer night sky. The moon was bright, the stars beautiful. I sighed and said:

``...Well, I guess that's just how it is. There are plenty of people without imagination in this world. I don't even know if I have any myself.''

``Yakumo-kun worries whether his lack of imagination might hurt someone. That shows you do have it.''

We sat on the bank of the Abukuma River. I wondered if there would be fireworks this year.

``I wonder if it's okay for me to keep playing the piano---'' Yuzuki said, looking very uneasy. ``Playing the piano doesn't help anyone, does it? It doesn't fill anyone's stomach, doesn't save anyone's life. It can't save anyone...''

I couldn't easily say, ``That's not true.'' Yuzuki was genuinely troubled. I was genuinely troubled too, and in the end, I couldn't find the right words.

``...There must be people who are saved by Yuzuki's piano.''

``...I hope so.''

Yuzuki had once mentioned that because she became unusually popular, some people came to her concerts mistaking them for idol concerts or something. That probably weakened her heart too.

A heavy silence fell. The murmur of the river sounded strangely lonely. Suddenly, Yuzuki said:

``Hey, Yakumo-kun. If I were gone, what would you do?''

``Eh---?''

I looked at Yuzuki involuntarily. She was staring straight at me.

``My mother says Fukushima is dangerous, so after graduation I should study abroad in Italy...''

``So that's what you were fighting about...''

Yuzuki nodded.

``What do you think?''

Her eyes seemed almost pleading. I looked into them for a while, then said:

``...I think you should go abroad.''

Yuzuki looked surprised. She shook her head slightly.

``...Why?''

``...In the end, nobody knows what will happen with the nuclear plant. If you can evacuate, you should.''

``...Don't say the same thing as my mother...''

Yuzuki suddenly stood up. Then she walked into the night river.

With a splash, she disappeared beneath the surface.

I was stunned and stood up involuntarily.

Yuzuki emerged almost immediately. Her wet hair was blacker than the night sky.

My heart pounded with anxiety. In the end, I too was afraid of radiation. The Abukuma River looked like a river of poison. No matter how beautiful the scenery, I had unconsciously been conditioned to see it as contaminated.

``Radiation can go to hell---'' Yuzuki said with strong eyes. ``I was born in Fukushima. I grew up eating the food of this land and drinking its water. I can't leave that easily. Even though I'm so sad, so frustrated, if I went abroad, these feelings wouldn't reach anyone. They'd just say I ran away, that I abandoned Fukushima!''

And then she cried. I was somewhat stunned by Yuzuki's passion. It was an emotion I lacked. I didn't have that much love for Fukushima. I had never thought of Fukushima as my homeland. I had never felt love for Fukushima.

That certainly wasn't Fukushima's fault---it was mine.

I thought about the difference between ``lightness'' and ``light-heartedness.'' They seem similar, but they are completely different.

Yuzuki's soul was light-hearted, but not light. My soul was not light-hearted, but it was light. There was no gravity pulling my soul. Not to my school. Not to my family. Not to my homeland.

I was just there because I had nowhere else to go. Like hair swirling around a drain, I was just being sucked in by emptiness, drowning, struggling, getting tangled up.

I don't know why I turned out that way. The precious things Yuzuki had received from Fukushima, I had been dropping and losing along the way.

Realizing that in such a short time, I felt---ashamed.

My soul, light as a helium balloon, was unbearably shameful.

As if making excuses to myself, I said:

``...Yuzuki, it's not about running away or abandoning Fukushima. You're going to Italy to study. To learn from good teachers, to play more beautiful piano. To be able to convey your feelings through beautiful performances, not through worthless words.''

Yuzuki fell silent again. Then, in a voice like a bubble rising from the bottom of the water, she said:

``...And you, Yakumo-kun?''

``Me?''

``Yakumo-kun can't manage without me. Like a goldfish scooped up at a festival, you'd die right away. You're already groaning and useless. There's no way you'd be okay if I left.''

She was worried about me, I thought. And then I considered: if I told her I couldn't manage without her, would she stay in Fukushima?

---I didn't have the right to do that. My weightless soul should not hold back Yuzuki's heavy soul. The Earth and the Moon are together because they both have mass.

``...I'll be fine. I'm not a kid anymore.''

Yuzuki looked down and bit her lower lip. Then she wiped her tears with her arm. She climbed up the bank and said:

``...I understand. Goodbye.''

She passed by me and walked away with heavy steps...

Her figure was so heartbreaking that I couldn't help but call out.

``Yuzuki---'' She stopped. But I still couldn't find the words, and said, ``Take care...''

Yuzuki half-turned toward me and then walked away again with heavy steps.

After her white figure disappeared, I also started walking.

Yuzuki's wet footprints dotted the road here and there. I went home by a different path.


\subsection*{5}

In November, Yuzuki's new CD was released. The cover photo enraged her.

That summer night---Yuzuki's figure as she entered the Abukuma River, tears streaming down her face, had been turned into the CD cover just as it was. The photo was beautifully edited. Yuzuki glowed faintly white like the moon in the night sky, and the river's surface reflected the starry sky like a mirror.

But Yuzuki's raw emotion was vividly sealed within it. The title was ``SADNESS.'' The subtitle was ``Prayer for a Wounded Homeland.''

The CD sold explosively. On top of the already excellent quality of her performance, the fact that Yuzuki was from the disaster area sparked sensational interest. The CD was released with that calculation from the start. It was obvious who had taken the photo.

Kitajo Takashi---that day, he had apparently been staying at the Igarashi house for work. When Yuzuki confronted him, he claimed it was a coincidence.

``I couldn't sleep and went for a walk, and I saw Yuzuki-chan enter the river. The scene was so beautiful that I couldn't help but press the shutter.''

It was doubtful whether this was true. Yuzuki protested fiercely.

``Who gave you permission to use a secretly taken photo like this for the cover!''

``I did---'' Ranako said. Then, to Yuzuki, ``It's a good opportunity, isn't it? Thanks to that photo, it's selling worldwide.''

Yuzuki must have shaken her head in disbelief.

``You monster!''

She screamed and ran out of the house, running through the rain without an umbrella, soaking wet, and burst into my apartment. Yuzuki and I had been in a cold war since that summer night. The fact that she fled here suggested that an even worse war had broken out elsewhere.

Unsure of what to do with the sobbing Yuzuki, I handed her a pure white towel first.

``This... it's new... hic... it doesn't absorb water!''

She said this between sobs, so I switched it for an old bath towel---I strangely remember that detail. Yuzuki took a shower and changed into my jersey. She still kept crying, apologizing over and over. She was sorry to the other disaster victims. This would naturally be called self-promotion.

``I... live in Koriyama... and during the earthquake I was in Tokyo... I haven't suffered at all... just because I live in the disaster area... I act like I'm the most wounded person in the world... when there are so many people who lost family... They'd be angry seeing this CD... saying I'm exploiting people's sadness... what if someone gets hurt because of this...''

Yuzuki cried like a child. I also felt like crying along with her, and secretly checked Amazon reviews and personal blogs for feedback on the CD. No matter how good something is, if enough people evaluate it, there will always be a certain amount of abuse.

However, the reviews for ``SADNESS'' were strangely all positive.

``...What do you know? Everyone understands that Yuzuki is a hardworking, kind girl.''

Yuzuki cried even harder. I didn't know what to do, so I took out the CD of Tanaka Kiyoko from the shelf and played it. Beautiful piano tones flowed out.

Yuzuki stopped crying for a moment and listened. Then:

``Waaaaah...! Kiyoko-sensei...! Kiyoko-sensei...!''

She began to cry even more intensely. She was crying no matter what I did.

There was, of course, no teacher-student relationship between Tanaka Kiyoko and Yuzuki. Tanaka Kiyoko passed away on February 26, 1996. Yuzuki was born the following year, on March 3, almost as if passing each other.\footnote{Tanaka Kiyoko (田中希代子, 1932–1990) was a Japanese pianist and the first Japanese prize-winner at the 5th International Chopin Piano Competition in 1955.}

Nevertheless, through the recorded piano tones, Yuzuki offered such respect and love, calling her ``sensei,'' treating her as a model, and even finding a pillar for her heart in her---

That seemed so beautiful to me, and it made my heart warm.

Listening to the piano, Yuzuki eventually cried herself to sleep. A single droplet remained on her cheek, like transparent ice beginning to melt. A tear as beautiful as the tones played by Tanaka Kiyoko-sensei.


\fullpageimage[]{9.jpg}
% ==============================
%        Chapter 5
% ==============================
\clearpage
\thispagestyle{plain}
\phantomsection
\begin{center}
    \large\bfseries Chapter 5
\end{center}
\addcontentsline{toc}{chapter}{Chapter 5}
\label{sec:ch5}
\vspace{1cm}

\subsection*{1}

We became high school students.

Despite my health problems, I managed to get into a decent college-preparatory high school in the city.

Yuzuki went to study at the Milan Conservatory in Italy.\footnote{The Milan Conservatory (Conservatorio di Milano) is one of the oldest and most prestigious music schools in Italy, founded in 1807.}

For the first time since meeting Yuzuki, I spent days without her.

To put it plainly, just as Yuzuki had said, my high school life was as lifeless and hopeless as a goldfish scooped up at a festival. Since it was a college-prep school, I had to study for entrance exams and worry about my mock exam scores, but I couldn't feel any reality in those days. The pain of the earthquake was overwhelmingly more real. What's the point of getting into a good university? I thought. What Yuzuki had carefully picked up and I had shamefully dropped wasn't grades. I didn't want to become rich, nor did I want to be respected, nor did I want to do meaningful work.

I just wanted to become a decent person.

I wanted to fill something that was irreparably lacking.

I began to read voraciously. Just as I had collected flowers in elementary school and ``flowers'' in middle school, in high school I began to collect ``stories.''

The trigger was seeing an article online about the ``Contributor Awards for the Great East Japan Earthquake.'' Stories of people who risked danger to save others, or who actually sacrificed themselves, were compiled in short articles. Despite the matter-of-fact writing, I cried reading them. The brave figures of people who displayed their humanity in times of crisis were beautifully drawn in my mind. I felt as if I, along with the people they had saved, had been saved as well.---I thought that beautiful stories could save me. That they could save something within me that couldn't be saved.

Just as the flowers didn't have to be real, the stories didn't have to be nonfiction. Ridiculous, absurd fiction was fine. To carve something beautiful out of thin air, a chisel called ``lies'' was necessary. What mattered was that it was made with sincerity. That it was written with blood and soul. I wanted the real thing, even if it was rough. I didn't want well-made fakes. I hated novels that were like skillfully manufactured products. Even if they weren't perfect, I wanted something that conveyed passion. That was similar to the condition for a good parent for a child.

One day, on a whim, I picked up one of my father's novels. It was frustrating, but I thought it was a good novel. There was something there that good novels have and bad novels never do. There was a certain urgency, as if he was writing at the cost of his lifespan. But I felt that the lifespan he was cutting short might have been my mother's, which made my feelings complicated.

Yuzuki had gone to Italy and never contacted me. Since I had said I would be fine without her, contacting her felt like admitting defeat. Still, I secretly kept track of her activities. I remember a video uploaded to YouTube in May. It was a concert video. Yuzuki was wearing a simple black dress.

In mourning---

That's how she looked, standing in the dimness at the back of the stage as if she might disappear into it.

The graceful performance began. Frédéric Chopin's Nocturne No. 2 in E-flat major, Op. 9 No. 2\footnote{Chopin's Nocturne No. 2 is one of his most famous and beloved works, known for its lyrical melody and expressive depth.}---

I noticed that the quality of Yuzuki's performance had changed. There was an urgent resonance in the sound. I could almost see the starry night sky in my mind. It reminded me of Tanaka Kiyoko-sensei's piano.

Yuzuki had begun to play again as if in prayer---


\subsection*{2}

Before I knew it, summer vacation had arrived.

Summer school was held, but I never attended. I spent my days reading books in a dark room, staring blankly at the summer blue sky through the window.

I heard from Shimizu that he would be going to Koshien,\footnote{Koshien (甲子園) refers to the national high school baseball tournaments held at Hanshin Koshien Stadium in Hyogo Prefecture, a prestigious event in Japanese sports culture.} so on August 11 I turned on the TV.

The summer Koshien was broadcasted in all its brilliance.

Shimizu, as a first-year, was wearing the number 4 for Seiko Gakuin.\footnote{Seiko Gakuin (聖光学院) is a real high school in Fukushima Prefecture, known for its strong baseball program.} What a man, I marveled.

Their opponent was Nihon University's Third High School.\footnote{Nihon University's Third High School (日大三) is a well-known baseball powerhouse from Tokyo.} Standing in the batter's box, Shimizu hit the first pitch. A sharp hit flew. The second baseman dove for the ball and immediately threw to first. Shimizu slid in---

The call was safe. I let out a breath. I had been holding it without realizing.

``Shimizu has gotten so fast.''

I muttered to myself. A huge cheer rose from the stadium.

With both teams scoreless, in the bottom of the 8th---with two outs and a runner on second, Shimizu's turn came. He adjusted his grip carefully, aiming his sharp eyes far beyond the pitcher. Then he swayed his body with an almost playful air. He's going for a home run, I knew.

On the third pitch, Shimizu swung with tremendous force.

A sharp sound rang out. The ball flew high into the air. The announcer's excited voice---

``He got it! High! Will it go!''

I knew it would. In the brief moment he appeared on screen, Shimizu was definitely laughing: ``Wahaha!'' Shimizu was still the Great Demon King, and I was happy.

``It's gone! Home run!''

The dugout greeted Shimizu with smiles, slapping him on the back. Wherever he went, he was loved.

In the top of the 9th, Nihon University Third couldn't score the two runs they needed, and Seiko Gakuin won 2-1.

Seiko Gakuin was overjoyed, and Nihon University Third was in tears---both looked magnificent.

Suddenly, I reflected on myself, sitting alone in a dark room.

What on earth was I doing---?

Thinking that, I felt a little like dying.


\subsection*{3}

The new semester began with interviews.

My homeroom teacher, Sumida-sensei, was a social studies teacher in his mid-forties---tall and lanky, wearing a brown jacket all year round, with a square face. His black-framed glasses were also square, making his face look like a geometric shape. Everyone in the class came down with a condition where they saw all sorts of buildings in the teacher's face.

Sumida-sensei blinked his small eyes behind his square glasses and said:

``Yakumo~, do you have any motivation~?''

I thought he was probably talking about my grades. I thought for a moment and said:

``No.''

He blinked again.

``...But this is a college-prep school, you know~''

``I know it's a problem, but I just can't seem to find any motivation.''

``Are you feeling unwell~?''

``Maybe.''

``...Then go get some counseling~. You're not stupid, you know~''

``I don't want to.''

``...Then it's a three-way conference, you hear~!''

I absolutely didn't want a three-way conference. I decided to get counseling. But since I knew it would be useless, I went straight to the psychiatric department of a hospital. If it could be cured, I wanted to be cured.

The doctor seemed quite intelligent. I swallowed my embarrassment and talked about my special phantom-limb-like condition.

``---And the pain I felt during the earthquake was so great that since then, I haven't been able to feel any reality in normal life. Only the earthquake remains real. Right now it's like I'm living in a mirage. Novels full of lies but incredibly powerful, interesting games, beautiful music, works of art---they all feel much more real to me than daily life.''

The doctor looked stumped and exchanged glances with the nurse behind him. After rubbing his right cheek with his hand, he said:

``It sounds like Tokatonton.\footnote{Tokatonton (トカトントン) is a short story by Osamu Dazai about a man haunted by a phantom hammering sound after Japan's defeat in World War II.}''

``Tokatonton?''

``By Dazai Osamu.---The protagonist is a man who was a soldier. When Japan accepted the Potsdam Declaration and lost the Second World War, the man heard the sound of a hammer, \textit{tokatonton}. Since then, whenever he tries to immerse himself in something, he hears a phantom \textit{tokatonton} from somewhere, and everything immediately seems foolish, so he gives up---'' He searched on his laptop, ``Ah, it's not \textit{tokatonton}, it's \textit{tokatonton}--wait, I got it wrong.''

The doctor showed me his laptop, and I read \textit{Tokatonton} from Aozora Bunko.\footnote{Aozora Bunko (青空文庫) is a Japanese digital library of public domain literature, including works by Osamu Dazai.}

It was a short story in the form of a letter sent to the author by a man tormented by the \textit{tokatonton} sound that even shattered nothingness. ``Please tell me: what is this sound? And what should I do to escape from it?''---

The author responded with these words:

Matthew 10:28: ``And fear not them which kill the body, but are not able to kill the soul: but rather fear him which is able to destroy both soul and body in hell.''

``If you can feel a flash of lightning in these words of Jesus, your phantom sounds should stop. \hfill Osamu Dazai''

``What does that mean?''

I asked the doctor. He rubbed his right cheek again and said:

``Honestly, I'm not sure.''

I thought for a moment and said:

``I think this man was suffering because the illusion of Japan as a nation had collapsed with the defeat. Isn't that like the death of the god he believed in? Isn't it a bit nonsensical to make someone who's troubled by that listen to the words of the enemy's god?''

``Huh? Are you criticizing Dazai Osamu?''

``Huh? Oh, is this meta-commentary that's aware of all that?''

``Huh?''

``Huh?''

``...I'm not really sure, but aren't you overthinking it?''

``No, I'm actually suffering. Even you, sensei, would think desperately about any difficult problem if your family had a serious illness, wouldn't you?''

Sumida-sensei fell silent. His face gradually turned pale.

``I think you thought of \textit{Tokatonton} because of the similarity between defeat and the earthquake. The protagonist was in a state of collapse after his illusion of the Japanese state was shattered. He couldn't adapt to postwar daily life. On the other hand, I never believed in anything from the start. There were no gods or Buddhas to begin with. The everyday world was destroyed by the tsunami, and I just returned to everyday life---''

Sumida-sensei raised his hand as if to say ``enough.''

``...I get it, I get it. It seems your head has gotten a bit strange from the shock of the earthquake. I'll prescribe some medicine, so take it properly.''

``It's not that I went crazy from the shock of the earthquake. I've always been a little strange.''

``...Agh, I can't stand it, I can't stand it!'' Sumida-sensei suddenly raised his voice. ``I never should have become a psychiatrist!'' She hugged her knees on the round stool, rubbing her right cheek vigorously, and spun around crying. ``You people are all, all, all crazy! Disgusting! Don't come to my hospital, you crazies! No entry! Only normal people should come to me!''

Startled, the nurse standing behind her gently massaged Sumida-sensei's back with a practiced hand, and said to me apologetically:

``I'm sorry. She's been off ever since one of her patients committed suicide.''

Somewhat dumbfounded, I said:

``...Shouldn't she see a doctor?''

``She is a doctor.''

I realized that was true. I left the examination room and picked up the antidepressants at the hospital pharmacy.

On my way out, worried about Sumida-sensei, I went around to the back of the hospital and looked up at the examination room window.

---I was startled. Sumida-sensei was looking down at me from the window. She looked somehow dejected, like a shadow cut away from reality with a butter knife, a lonely figure.

``Take care of yourself,'' she murmured.

I was deeply moved by that sight for some reason.

``Thank you.''

I bowed and left the hospital. Sumida-sensei was probably watching my back the whole time.


\subsection*{4}

The antidepressants worked surprisingly well. The pain I had always felt when looking at anything seemed to become somewhat faint. The normal world I had seen before my senses twisted flickered before my eyes like a mirage. For some reason, I felt very nostalgic. The medicine changed my perception and made it hard to read books, so during that time I listened to music. Since it became harder to notice flaws, it sounded more beautiful than usual.

---I wondered whether those who have ears to hear or those who don't are happier.

But in the end, the antidepressants didn't save me. When the medicine ran out, all that remained was the loneliness of having finished a delicious bottle of ramune.\footnote{Ramune (ラムネ) is a Japanese carbonated soft drink sold in a distinctive glass bottle with a marble stopper.} I never went back to the hospital.

Around that time, Yuzuki seemed to have become devoted to Chopin. I listened to polonaises and researched Chopin.

Chopin was a Polish composer of the Romantic era. Gifted from a young age, he composed the Polonaise in G minor at age seven. Plagued by tuberculosis throughout his life, he lost his sister to the same disease when he was seventeen.

On November 2, 1830, at twenty years old---already a successful performer and composer, Chopin decided to leave his homeland due to worsening domestic conditions. At the time, Poland was partitioned among Russia, Austria, and Prussia, and independence movements were on the rise.

``I feel like I'm setting out to die.''

It is said he wrote such words in a letter to a friend. He must have had a bad premonition. Whether true or not, Chopin supposedly set out for Vienna, Austria, carrying a ring exchanged with Konstancja Gładkowska and a silver cup filled with soil from his homeland.

On November 29, 1830, the November Uprising broke out.\footnote{The November Uprising (1830–1831) was an armed rebellion in Poland against the Russian Empire.} Armed citizens drove the Russian army north of Warsaw. The patriotic Chopin wanted to join the revolution, but his friend Tytus is said to have advised him: ``You should serve your homeland through music,'' so he remained in Vienna.

However, in Vienna, anti-Polish sentiment rose in response to the November Uprising, and Chopin was treated coldly. In the end, Chopin left Vienna without achieving much.

Then, in Stuttgart, Germany, he learned that the Warsaw revolutionary forces had been crushed by the Russian army. He must have spent very painful days. Worrying about the safety of his family and friends in his homeland, tormented by loneliness... Heart-wrenching sorrow must have struck him.

As his ominous premonition foretold, Chopin never set foot on his homeland again.

I felt I understood why Yuzuki had become devoted to Chopin.

The sadness of longing for one's homeland but being unable to return, the anger of having one's homeland wounded with no outlet---she probably saw herself in Chopin, who could only pour such feelings into his music.

On December 25, 1831, Chopin confided his feelings in a letter to Tytus:

``I appear cheerful on the surface, especially among my 'comrades' (by which I mean Poles). But inside, I am always tormented by something. Forebodings, anxieties, dreams---or sleeplessness, melancholy, indifference---a desire for life, and the next moment a desire for death. A kind of pleasant peace, a numbing vagueness, but sometimes clear memories revive and make me anxious. Sour, bitter, salty---my feelings are terribly mixed and deeply confused.''

Reading this passage, I felt as if my own inner life had been described. I was certain Yuzuki must have similar feelings. I heard that in Polish there is a word ``ZAL.'' It is said to be a uniquely Polish emotion---``forlorn resignation,'' ``source of deep resentment,'' ``fierce protest,'' ``sorrow at losing what should be there''... I interpreted it as a massive sense of loss and the accompanying hatred, sadness, and the helplessness of being stunned into stillness.

Didn't we who suffered from the disaster feel ``ZAL''?

Thinking so, I felt that the feelings of the protagonist of \textit{Tokatonton} were also close to ``ZAL.'' He too, while still in his homeland, lost it due to defeat.

Then why did ``ZAL'' make Chopin's and Yuzuki's piano beautiful?

It is said that Schumann described Chopin's works as ``cannons hidden among flowers.'' This likely means that beneath the surface of gorgeous, elegant music lie passion, sorrow, and a rebellious spirit.

We are captivated by the pitiful flowers, but I think what Chopin really wanted to express was the cannons. Even if he presented the cannons as they were, people wouldn't accept them. So he concealed them with flowers and handed them over as bouquets. There, a touching human beauty is born.

I recalled the countless flowers offered to the dead. The gesture of offering beautiful flowers before memorial tablets. And Chopin's gesture of wrapping fearsome weapons in beautiful flowers.

Weren't those gestures of prayer? Not as firm as praying to God, but gentle gestures. Isn't it exactly such quiet gestures of prayer that make the piano sing beautifully...?


\subsection*{5}

Was I also praying?

Was I turning pages of novels as if turning pages of the Bible?

Was I continuing to throw bouquets, hoping that someday the dark hole would be filled?

But only monks can live by praying; as just an ordinary high school boy, I continued to charge straight down the path to being a dropout. I just couldn't study. When I studied, I felt like I was dying.

``You're an idiot,'' said my classmate Sekigahara.\footnote{The name Sekigahara (関ヶ原) is a reference to the famous Battle of Sekigahara in 1600, but here it is simply a character's name.}

``At this rate, you'll be left behind by Fukushima.''

``Huh? What do you mean?''

It was a classroom at dusk. The textbook from the last class was still on my desk. I had been pretending to listen and reading a book, and before I knew it, class was over and evening had come.

``People who stay in Fukushima are losers. Fukushima is already finished. Actually, it was finished from the start. It's full of idiots with no ambition, who are fine staying idiots, and who are enthusiastic about dragging others down. They're arrogant. They put mayonnaise on hiyashi chuka.\footnote{Hiyashi chuka (冷やし中華) is a chilled Chinese-style noodle dish popular in Japan during summer. Some Japanese enjoy adding mayonnaise to it, though it's a matter of regional and personal preference.}''

``You're arrogant and put mayonnaise on hiyashi chuka too.''

``...That's my cute point, so overlook it. I'm saying Fukushima is a land of losers.''

``The concept of 'losers' didn't exist in the first place.''

``You always say things are embarrassing. The embarrassing ones are the losers.''

``I'm only ever embarrassed of myself.''

``Try projecting that feeling onto others. Fukushima is a cesspool of embarrassing people.''

I gave it a try.

``...Sorry, I don't mean to be rude, but you're the most embarrassing one of all.''

``...Tch.'' Sekigahara clicked her tongue and stood up. As a parting shot, she said, ``You're gonna end up a pathetic loser too, you idiot!''

And she stormed off. She had approached me saying ``I feel like you're different from the other idiots,'' only to turn on me so quickly.

Later, Sekigahara tried to leave me and join another group, but failed. She was rather disliked. In the end, she came back to me and ate hiyashi chuka with mayonnaise for lunch.

That was Sekigahara's cute point.


\subsection*{6}

I became a second-year high school student.

Until middle school, I had felt a sense of progress each year. But after entering high school, I couldn't shake the feeling of marking time in the same place. Only piles of books grew in the corner of my room. Nothing changed; only the seasons passed.

---Summer came. Shimizu was supposed to go to Koshien again. But it didn't happen.

He had been in an accident. While on an early morning run in light rain, at an intersection with poor visibility, he was hit by a car traveling at 80 kilometers per hour.

As soon as I received the news, I dropped everything and headed to the hospital in Fukushima City where Shimizu had been taken.

At the hospital, three boys from our middle school baseball team who had been contacted were already there, along with eight of Shimizu's teammates. Aida was also there. He had gotten a bit flashy since high school. More people kept arriving, eventually reaching twenty-eight. In front of the operating room, we prayed for Shimizu's safety. It was as if we had returned to the club locker room, nervous before a big game.

---Eventually, the ``surgery in progress'' light went out, and the surgeon emerged.

``Doctor...! Is Shimizu...?''

Aida asked. The surgeon removed his mask and smiled.

``He's going to be fine.''

We sighed in relief and smiled with reassurance.

---But that smile died when we saw Shimizu.

Shimizu had lost his left leg below the knee. The stump was wrapped in bandages, like a cocoon. I remembered last summer, in the game against Nihon University Third, when Shimizu slid into first and was safe. Shimizu had been fast too. The legs he used to hit and run were gone...

The emptiness left by Shimizu's lost leg brought me terrible pain. It was so great, so sad, that I couldn't help crying. Then everyone else packed into the hospital room started crying too. Apart from his left leg, he was covered in injuries like a mummy, and Shimizu, who had been dazed, hurriedly said:

``Everyone, don't cry, it's okay.''

He said it with a strangely bright smile.

``I'm going to kill him---'' Aida, next to me, suddenly said, his face pale. ``I'm going to kill the bastard who was driving recklessly and took Shimizu's leg.''

He tried to push through the crowd and leave the room. Chaos erupted.

``Whoa! He's serious! Stop Aida!''

The muscular members of Seiko Gakuin's baseball team grabbed Aida, and he fell to the floor. More guys piled on top, pinning him down. With his cheek pressed against the linoleum floor, his face looking like a Hyottoko mask,\footnote{Hyottoko is a traditional Japanese comical character with a distorted face, often depicted in festivals and masks.} Aida cried:

``Ugh... damn it... damn it...!''

He cried with so much frustration that his tears and snot formed a puddle on the floor.

``I'm really okay, everyone, you really don't have to cry,'' said Shimizu, the only one smiling until the end.


\subsection*{7}

Shimizu was transferred to a hospital near his parents' house. Visitors to his room never ceased. Colorful flowers and fruits were always on the table. Shimizu always had a round smile, like a prince from a happy kingdom. Among the visitors were girls. Kobayashi Koyomi was one of them. She had moved to Soma after elementary school, and I hadn't seen her since. She had apparently become friends with Shimizu and they had been keeping in touch.

She was quiet and unassuming, but had grown into a cute girl.

One day, when I went to visit, I happened upon her peeling an apple for him. She cut the apple into rabbit-shaped pieces. Her careful knife work alone suggested a meticulous nature and kindness. When she finished, she smiled and handed them to Shimizu. He would examine them happily, then eat them in one bite. Over and over. It was like little rabbits hopping back to their burrows.

Both Shimizu and Kobayashi were adorable and heartwarming. I walked back the way I came, feeling happy.

I visited Shimizu more than anyone else. Kobayashi lived in Soma, so she couldn't come that often, and everyone was busy with studies, clubs, or romance. I, on the other hand, was overwhelmingly free. As free as a seal lounging on an ice floe.

Shimizu knew this too, and often messaged me, ``Yacchan, come.'' I replied, ``Okay.'' Checking now, I've never refused even once. I thought about bringing flowers to fill the void of Shimizu's lost leg, but bringing them every time would be a bother, so instead I stopped by a used bookstore and bought one volume of \textit{JoJo's Bizarre Adventure} at a time.\footnote{\textit{JoJo's Bizarre Adventure} (ジョジョの奇妙な冒険) is a long-running Japanese manga series by Hirohiko Araki, known for its unique art style and creative storytelling.} As expected, Shimizu got hooked and started shouting ``Scarlet Overdrive!'' and such. What an idiot. I did it too. Since we were idiots, it was a lot of fun.

Sometimes Shimizu had phantom limb pain. His supposedly missing leg hurt. Whenever that happened, Shimizu would groan and grit his teeth, hugging his left knee. I too felt pain in that void at his knee, and stroked his broad back.

Seiko Gakuin went to Koshien without Shimizu. All of us from the middle school baseball team gathered and cheered from the hospital room TV. Aida brought a ``Fukimodoshi''---a toy paper tube that extends when blown---and it was annoyingly loud.\footnote{Fukimodoshi (吹き戻し) is a traditional Japanese toy consisting of a paper tube attached to a whistle; when blown, the tube extends and then curls back.}

Seiko Gakuin won their first-round game against Aichi Kogyo Meiden 4-3.

Six days later---they faced Fukui Shogyo.\footnote{Fukui Shogyo (福井商) is a Fukui Prefecture high school baseball team.}

We gathered again, and Aida brought that annoying toy again, and on top of that, he seemed to be having a messy breakup with his girlfriend, so he was really loud.

Fukui Shogyo scored 1 point in the bottom of the 1st inning. Both teams remained scoreless for a while, and the tension dragged on until the top of the 6th, when Seiko Gakuin finally tied the score. We cheered too loudly and were warned by the nurse. Then in the bottom of the 8th, Fukui Shogyo scored another run. We held our heads. Aida was dumped by his girlfriend.

In the top of the 9th, Seiko Gakuin couldn't score the run they needed, and they lost.

I thought, if Shimizu had been there, maybe they could have won.

When the defeated players crying on screen appeared, a strange sound came: ``Uoooh.''

Looking, Shimizu was sobbing. The face that had been smiling all along was now terribly contorted. The emotions he had hidden deep inside must have overflowed uncontrollably. His frustration, sadness, and regret were painfully conveyed.

We cried along with him. Everyone had wanted to see Shimizu perform.


\subsection*{8}

At the beginning of November, when the winter chill crept in, Shimizu was finally discharged from the hospital.

We had accumulated up to volume 63 of \textit{JoJo's Bizarre Adventure}, through the end of Part 5.

Shimizu moved to a rehabilitation center for training. He would wear a temporary prosthetic leg and train to return to society.

``Prosthetics are expensive, right?''

When I asked, Shimizu said it wasn't so bad.

``I got a Level 4 physical disability certificate, so apparently the temporary leg costs 30\% and the permanent one costs 10\%. Plus, I got about 50 million yen in compensation.''

I had no idea if 50 million was a lot or a little.

I watched Shimizu grit his teeth and train with his temporary leg from outside the training room window. He had lost so much muscle mass that even simple movements looked painful. I wondered if Shimizu would ever run again.

---Around the time high school went on winter break, gossip about Yuzuki appeared. A male student at the Milan Conservatory posted a two-shot with Yuzuki on Facebook. The background was some square, and the man was a tall, blonde, blue-eyed handsome guy. He seemed to be a promising pianist. I marveled at the existence of such CGI-like beautiful men... but I was also anxious. Was he her boyfriend?

Since Yuzuki herself didn't disclose any personal information, it seemed like an unexpected leak. It became a topic on some parts of the internet, but no further information surfaced.

I wondered whether to message Yuzuki. I hadn't contacted her even once since she went to Italy. I typed ``Long time no see''---and deleted it. It felt clingy and pathetic. I gave up and listened to Yuzuki's new performance uploaded to a video site.

Her playing was becoming more and more beautiful. I resented myself for not changing at all.


\subsection*{9}

``You can't go on like this,'' Shimizu said. I looked up from the book I was reading.

It was mid-February, in the evening hospital room, nearing night. Since the accident, Shimizu's nose had become a little crooked, but his eyes were straight.

``Yacchan---you can't keep doing this forever.''

His unusually serious tone surprised me. Shimizu was always smiling. I'd never seen him look so serious and composed.

``...Huh? Doing what?''

Shimizu paused for a moment, then said:

``I had my accident in July, lost a leg, and was discharged in November. I've been training for three months since then, and I can already walk well with the temporary leg. It's like I lost a leg and grew a new one. That's how much time has passed.''

``Grew a new one...''

I looked at the row of \textit{JoJo's Bizarre Adventure} volumes lined up by the bed. I had caught up to the latest volume, and the prices of the later used volumes had gone up, so it wasn't easy. That much time had passed—enough to read the entire twenty-plus years of work by Hirohiko Araki. Shimizu said:

``When I get the permanent leg and go back to school, I'm going to rejoin the baseball team. I'll practice desperately and go to Koshien. And I'll hit a home run.---What about you, Yacchan?''

I thought he said something amazing. I was taken aback.

``Me... nothing...''

``Yacchan, I think you're much more amazing than you think you are. I think you can do something. But I think you'll really go down the drain if you keep going like this.''

I was at a loss. For me, Shimizu was the embodiment of magnanimity. Even he said I couldn't go on like this. It seemed the situation was that bad.

``...But now, I don't even know what to do.''

Shimizu said what seemed to have been on his mind:

``Why don't you write a novel, Yacchan?''

``A novel...?'' I was surprised. ``Huh? Why a novel?''

``Conversely, I don't understand why you wouldn't write. You're always reading novels, and your father is a novelist, right? Don't you have talent?''

``...No, I don't, I don't.'' I waved my hand in front of my face. ``I'm a reader. I've never thought about writing.''

``I think you do.''

``Huh? Why?''

``Don't you remember? The reason I was able to go to Koshien was because of your composition.''

That was apparently from first grade of elementary school. Shimizu was in youth baseball, but he was just big; he didn't actually like sports, and was forced to do it by his parents. Apparently he felt depressed every time there was practice. Then one day, during an open class, I read a composition.

It was titled ``White Clouds.''

Shimizu stood at the plate, made a sharp \textit{kan} sound, and ran powerfully. Against the cloudless blue sky, a white baseball flew high. ``It looked like a cloud, like a free bird, and it felt so good I felt like laughing out loud, 'Wahaha.' I thought Shimizu-kun, who can do such amazing things, is really great.''

``---After that, I started loving baseball. And then I came to like myself playing baseball. It's all thanks to you, Yacchan.''

I didn't remember that at all. When Shimizu said it, I vaguely recalled it. When I finished reading that composition, Shimizu, who sat diagonally in front of me, turned around and smiled. The teacher praised me, and my mother, who had come for the open class, must have praised me too. She said, ``You're really Ryunosuke's child, aren't you?'' I remember feeling very happy about those words.

Come to think of it, Shimizu had been my friend ever since that day.

For some reason, tears welled up. Shimizu, also blinking back tears, said:

``I was so happy with that composition that I took it from you. It's still framed in my room. I read it every day.''

``I see...''

I had never imagined that Shimizu's ``Wahaha'' originated from my composition.

``So, the person I am now was built by a single sheet of composition paper you wrote. I've ended up like this, but I'm still going to keep trying. I'm going to prove that you're amazing. So, just trust me a little—write a novel. Consider yourself fooled...''

And Shimizu smiled. It was the same childish smile he'd had all along since that day.


\subsection*{10}

I began to write a novel.

I opened Word on my Windows PC and typed something. But it wouldn't become a story no matter what I did. It was like a memory-lost inchworm trying to learn how to walk from scratch—a clumsy trial and error.

I had devoured so many stories, so why could I produce nothing? Why was my head like an empty jar? I thought.

For a whole month, I couldn't come up with any story. But once I got past that struggle, ideas started bubbling up disgustingly fast, like bubbles bursting from the bottom of a boiling pot. I just had to wait until the invisible water reached the right temperature.

Then, I suffered from how bad my writing was. Since I'd only ever sought out and read first-rate works, my eye had grown too sharp. The gaze that had mercilessly dissected others' novels now turned on my own. The curses I had spat at boring novels now fell upon my own. As they say, curse others and you'll dig two graves; curses always come back around. To return the curse, I had to write with my life on the line. But the more serious I got, the more unbearable my own clumsiness became. Maybe I was concentrating too hard—I couldn't stand my poor writing and vomited often. I wondered what was wrong with my body. I couldn't understand the reason for a creature that writes and then vomits from its own bad writing. Could there be a more idiotic set-up?

In April, I threw down my pen. Then, on a whim, I decided to go to Seiko Gakuin.

Saturday, 7:30 AM. I took a train from Koriyama Station, got off at Fukushima Station, and then moved to Date Station. It was about an hour and a half. By 9 AM, I was already at Seiko Gakuin.

I secretly peeked at the field; practice had already begun. I searched for Shimizu.---There he was. Shimizu was apart from the group, silently moving his body in a corner of the field. There was still a noticeable awkwardness in his movements. His left leg prosthetic looked somewhat thinner than his right leg.

Swing, swing—Shimizu swung the bat. His body's axis was off, and his form had changed from before. I could see he was trying to correct it. He gritted his teeth. His face was drenched in sweat, as if he'd been caught in a rain shower.

Shimizu was serious. He really meant to go to Koshien---

Thinking that, my chest felt warm. Now wasn't the time to be doing this—I thought, and immediately left Seiko Gakuin and returned to Koriyama.

And then I began writing my clumsy novel again.


\subsection*{11}

In May, I finished my first novel.

No matter what kind of work it is, completion brings joy. I was over the moon alone in my room. I even began to feel it was a masterpiece. What to do with it---

After much deliberation, I decided to post it online, in a community known for harsh criticism.

The result was devastating. I think the main feedback was ``I don't get it.'' The flow of \textit{kishōtenketsu}\footnote{Kishōtenketsu (起承転結) is a traditional four-part structure used in Japanese and Chinese narratives: introduction, development, twist, and conclusion.} was unclear, and they didn't know how to emotionally engage. Some parts of the prose had moments that made them go ``oh,'' though.

I was made to realize that my sensibilities were a little strange. Both my reading and writing of stories deviated from others'. At this rate, no one would read my novels, and nothing would reach anyone.

Overwhelmed by the setback, I wrapped myself in my futon during the day. I realized the gap between Yuzuki and me. Her music reaches hundreds of millions, but my novel doesn't reach even one person.

---When I woke up, my head felt like it was splitting. I had slept for over twenty hours.

I crawled out of bed and started watching movies. I watched a series of masterpieces, and carefully read people's evaluations of them. What makes a work interesting? How do people appreciate it, and how do they feel? What do they see and what do they miss?---I studied these things relentlessly. It was like simulating an ordinary person. I was sad to think that I, who was lacking, had to try to act like a normal person.

But looking back on it now, that process was also a time of prayer. With cannons alone, nothing reaches anyone. They had to be hidden among bouquets. I was just gathering the flowers for that purpose. Gathering is my life.


\subsection*{12}

Before I knew it, summer had come.

While my classmates were busy with summer school, I watched movies, read novels, wrote scraps of writing, and threw them away. I still hadn't written a single line of a story. I had a sense of what I wanted to write, but I always stumbled at the opening. I couldn't help but wonder: Is this really right? I'd become paranoid, feeling as if some invisible critic's voice came from somewhere.

In the Fukushima High School Baseball Championship, Seiko Gakuin defeated Nihon University's Tohoku High School and secured a spot at Koshien.\footnote{Nihon University's Tohoku High School (日大東北) is a high school in Fukushima Prefecture, known as a regional rival to Seiko Gakuin.} Shimizu must have imposed grueling training on himself and overcome it. He made the bench. But he wasn't given a chance to play. The other players were also excellent, so the hurdle for Shimizu, who had only just become one-legged, must have been unbelievably high. From the stands, I secretly watched Shimizu, carrying complex feelings, cheering on his teammates from the bench.

And on August 19---Seiko Gakuin faced Saku Chosei at Koshien Stadium.\footnote{Saku Chosei (佐久長聖) is a Nagano Prefecture high school baseball team.}

I went with my middle school friends to cheer at Koshien Stadium. It was a clear, sweltering day. The temperature was over 30 degrees, and the sun was strong; Aida kept slathering on sunscreen. Seiko Gakuin won 4-2. We were overjoyed, but also conflicted.

Shimizu never once stepped into the batter's box.

We stayed overnight in Hyogo and fooled around. It felt kind of like a school trip.

``What are you up to, Yakumo?'' ``I'm writing a novel.'' ``Wow, cool! What kind?'' ``I haven't written a single line yet.'' ``Huh? Then you haven't written anything...'' ``I haven't written it, but I'm writing it...''

Everyone looked puzzled. I didn't know how to explain.

August 21---Seiko Gakuin faced Omi.\footnote{Omi (近江) is a Shiga Prefecture high school baseball team.}

It was also terribly hot that day. The game was a tense one. Both teams were scoreless until the bottom of the 5th, and in the top of the 6th, Omi scored first. 7th---8th---the score remained 1-0. In the top of the 9th, Seiko Gakuin held Omi scoreless without trouble. And then, the fateful bottom of the 9th---Seiko Gakuin scored 2 runs for a walk-off victory. The joy was all the greater for the fierce contest.

But again, Shimizu didn't get to bat.


\subsection*{13}

August 22---Seiko Gakuin faced Nihon Bunri.\footnote{Nihon Bunri (日本文理) is a Niigata Prefecture high school baseball team.}

It was also clear that day, but felt a bit more comfortable than the day before.

However, the game was tough. Nihon Bunri scored 1 run in each of the 1st and 2nd innings, while Seiko Gakuin struggled to score. Both teams were scoreless through the bottom of the 6th. Then in the top of the 7th, Nihon Bunri scored 1 run. For Seiko Gakuin, who was chasing, this was a painful point. Both teams were scoreless in the 8th. Then in the top of the 9th, the nail was driven in. Nihon Bunri scored 2 runs all at once. 0-5. It was a desperate situation. Bottom of the 9th, Seiko Gakuin was down to their last chance.---But the first batter flied out, and the second batter grounded out to short. In an instant, two outs.

We in the stands held our heads. It was over.

Shimizu's last Koshien was about to end.

---That was when it happened.

``Pinch hitter---number---13---Shimizu Kentaro---''

Shimizu's name was announced. At the very last moment, the manager gave him a chance.

Shimizu came out of the dugout.

``Shimizu...!'' ``It's Shimizu...!'' ``Shimizuuu!''

We shouted excitedly. Even from a distance, Shimizu's left leg looked somewhat thin. The fact that Shimizu was on the bench with a prosthetic leg was somewhat famous, and people around were talking about it.

Shimizu stepped into the batter's box. He looked up at the sky and took a deep breath. He must have been full of emotion. It was for this very moment that he had been tirelessly working.

Shimizu took his stance. He swayed his body with apparent joy.

My eyes were already welling up.

Shimizu still loved baseball. And now---he was going for a home run.

The pitcher threw. \textit{Ka-n}, a sharp sound; a foul ball rose high. A murmur spread. It was in response to the brilliance of Shimizu's swing.

Second pitch. The ball was sucked into the low-inside corner. Shimizu swung, but there was a certain unnaturalness in his motion. Maybe the low-inside corner was difficult with a prosthetic leg.

Third pitch. Aida and the others shouted: ``Shimizu! Hit it!'' I couldn't help gasping. The catcher set his glove in the same spot as before. I shouted too:

``Hit it!''

Shimizu swung. It was a magnificent swing, as if he had read the course from the start. He caught the low-inside pitch with the sweet spot.

A sharp sound rang out.

The baseball soared high into the air.

Against the cloudless blue sky, like a small cloud.

A roar of applause. Aida and the others shouted:

``Go! Go!''

I knew it would go in. Because Shimizu was smiling.

``Wahahaha!''

The Great Demon King, laughing, made a joyous tour of the Koshien bases.

The ball landed in the stands. Home run! A huge cheer enveloped Shimizu. Generous applause was given to the batter with the prosthetic leg. Shimizu pumped his fist and shouted: ``How's that!'' ``How's that!'' And then, laughing heartily again, he rounded second. Laughter mixed with crying. Shimizu, face contorted, ran while laughing and crying loudly. We couldn't hold back our tears either. Aida and the others stood up, thrust their fists toward the sky, and shouted through their sobs:

``Shimizuuu! Shimizuuu!''

I remained seated, covering my face and crying out loud.

I just couldn't stand up---

Thank you, I kept saying in my heart.

Thank you, Shimizu.

I will write a novel---


\fullpageimage[]{10.jpg}
% ==============================
%        Chapter 6
% ==============================
\clearpage
\thispagestyle{plain}
\phantomsection
\begin{center}
    \large\bfseries Chapter 6
\end{center}
\addcontentsline{toc}{chapter}{Chapter 6}
\label{sec:ch6}
\vspace{1cm}

\subsection*{1}

After I started writing novels, my view of \textit{Tokatonton} changed.

It wasn't a letter exchanged between a reader and a novelist, but between one novelist and another. In the middle of the letter, there is a passage describing the man's frustration when trying to write a novel.

``Should I end it with that kind of splendid, sorrowful conclusion, like the final chapter of \textit{Eugene Onegin}? Or should I go with a despairing ending like Gogol's \textit{The Marriage}? I was terribly excited, wondering, and as I looked up at the bare light bulb hanging from the high ceiling of the bathhouse, I heard it from afar: \textit{tokatonton}, the sound of that hammer.''

And at the end of the letter: ``Before I had even written half of this letter, I could already hear \textit{tokatonton} loudly. The tediousness of writing such a letter. Still, I endured and wrote this much at least. And because it was so tedious, I gave up and wrote nothing but lies. There is no woman named Hanue, and I didn't see any demonstrations. Most of the other things seem to be lies too. But only \textit{tokatonton} seems not to be a lie. I send this to you as it is, without rereading it.''

I feel that beneath this passage, the desire ``I want to write a novel'' is faintly visible. By writing lies, the letter becomes a novel. And by taking the form of a letter, he cleverly gets a professional novelist to read his novel. The novelist who saw through this replies:

``Matthew 10:28: 'And fear not them which kill the body, but are not able to kill the soul: but rather fear him which is able to destroy both soul and body in hell.'''

For someone who wants to write but cannot, these words strike like a thunderbolt. One must be conscious of the reader's eyes, but not overly so.

With Chopin and Dazai Osamu as my teachers, I began writing novels again. Then I submitted one to a publisher's new writer award. It was November.

New writer awards receive an enormous number of submissions, so the judging side also has a hard time. It's normal for it to take half a year for results, but the award I entered was the type that is judged in several rounds throughout the year, so the results came in two months.

Unfortunately, I was rejected——

My work had the highest evaluation, but no winner was selected.

However, I was able to meet and talk directly with the editor who had judged it.


\subsection*{2}

I took an overnight bus to Tokyo.

I was too nervous to sleep. I felt that something was beginning to move. Tokyo is much closer than Hyogo, where Koshien Stadium is located. Yet I felt like I was going much farther than I had then. The editor I was to meet was apparently quite famous; he had a Wikipedia page. I looked at the information on my smartphone and felt intimidated all over again. The list of works he had edited included many famous titles. He had also handled several writers who had won prestigious awards...

I bought e-books and read his works all night long.

---I don't like Tokyo. There are too many people, and it makes me feel sick. The information overload, like during the Uname Festival, continues endlessly. For some reason, I see too many things and hear too many sounds.

The meeting place was a coffee shop a little ways from Shinjuku Station.

At 2 PM, I entered the coffee shop and looked for the editor, Furuta-san.

The shop was three times larger than I'd imagined. The window-side seats were bright and open, while the windowless back was lit by warm-colored lamps. The floor was covered in crimson carpet, and the chairs and tables were made of ebony. It was a stylish place that somehow reminded me of the movie \textit{The Godfather}.

As I looked around, someone sitting in the back raised a hand lightly.

He was a strange man. A rather sturdy build, wearing a bright red leather jacket. A round face with square black-rimmed glasses. The discord with the background was striking. It was as if a Marvel Comics character had wandered into the set of \textit{The Godfather}. He seemed to be about forty.

\emph{This person?} I thought. He was nothing like what I had imagined. Wondering if the real editor might be hiding somewhere, I approached him while glancing around.

``Hello, I'm Furuta.''

He was indeed Furuta. We shook hands and sat down at the table. He said abruptly:

``You've got an aura!''

``Huh? ...Oh, is that so...?''

``I can sense talent! Yeah!''

Furuta seemed perfectly serious, but I felt as if my neck had suddenly been dislocated. \emph{Ah, he's a strange person}, I thought. Small talk began.

In any case, Furuta talked a lot. He was quick-witted, firing off topics one after another, jumping here and there, yet somehow never losing the thread. His tone had a unique quality; occasionally he'd throw in phrases like ``you know'' or ``I wonder,'' so-called \textit{onee} speech, which gave rhythm and drew you into his talk. He didn't seem to do it consciously. He was simply born with a gift for speaking. Before I knew it, four hours had passed, and I snapped back to reality. We hadn't talked about my work at all. At this rate, I'd have come all the way to Tokyo just for small talk.

``...Um... what did you think of my novel...?''

Furuta was momentarily dumbfounded.

``Ah! Right, right! That's what today was about!''

\emph{Oh my}, I thought. But Furuta suddenly became cool and composed.

``In a word---your novel is 'too intense.'''

``'Too intense'---?''

``It was interesting, though---the story is too difficult and too complex. I thought ordinary people wouldn't be able to follow it. It doesn't feel like it would sell.''

Furuta dropped a sugar cube into his cup, then fished it out and licked it like candy. The brain needs sugar to function. He continued:

``The writing is very good. The detailed metaphorical expressions are also skillful. There's something truly impressive about it.---But it's too rich in flavor. It feels like just one sip is enough. You're too serious, or rather, you're putting too much force into it.''

I thought, I don't remember any of this. I had tried so hard to simulate an ordinary person's sensibilities, yet there was still such a gap. I felt dumbfounded.

``...But I want to write seriously. Ordinary novels don't satisfy me. Unless it's something written with blood, I can't consider it good.''

Furuta popped a second sugar cube into his mouth.

``Works written seriously must be read seriously by readers. In this day and age, everyone is already overwhelmed by the real world and exhausted, so I think such things don't sell. It means people like you are a minority.''

...Was that really true? I fell silent. Perhaps I still didn't have enough flowers. Maybe I needed to cover the cannons with many more flowers---

``Here's my contact---'' Furuta took out a business card from his wallet. ``If you write a new work, please send it to me by email. I'll read it. I'm looking forward to it.''

Furuta quickly took the bill, paid, and left. I remained seated, thinking about what he had said.---Then suddenly, Furuta came back.

``Oh, I forgot to mention, you should change your pen name! It's uncool!''

And he left again. \emph{He came back just to say that?} I thought.

I changed my pen name.


\subsection*{3}

Exam season arrived. I didn't take any exams. I quietly wrote novels by myself.

My homeroom teacher at the time was persistent in trying to make me take exams, so I spent my days avoiding him. I understood his feelings well enough. I still remember passing Sumida-sensei, my first-year homeroom teacher, in the hallway. The tall, bamboo-like teacher stopped in the middle of the corridor.

``Yakumo~, are you still not feeling well~?''

I thought for a moment.

``...I'm better than before.''

He smiled.

``That's good to hear~.''

Even people who only brush past you in life can still worry about you.

I also told my father that I wouldn't be taking exams. We had this exchange via text:

``I see. What are you doing now?''

``Writing a novel.''

After a pause:

``I see. Good luck.''

Even after three years, our conversation ended in just three lines.

I smoothly graduated from high school.


\subsection*{4}

My father gave me a MacBook. The text was clearer than on Windows, making it much easier to write. My days of writing novels began in earnest.

I have almost no memories of this period. The desk and the MacBook were nearly my entire world.

Every day, I chewed up books, movies, and internet information and typed it all out. I rarely went out, so my face grew pale. Going to the nearby supermarket felt like an expedition. I'd never been particular about food, and I gradually lost weight.

---In two months, I finished a work. I thought I'd thinned it out quite a bit. I immediately sent it to Furuta by email.

Just one hour later, a reply came:

``Intense!''

...Huh? That's it? Having a novel I'd spent two months of my life on dismissed in a single word, I lost my composure and made a meaningless expedition to the supermarket. I searched meaninglessly for a radish with an obscene shape. The label on the bottled tea I picked up read ``Intense.''

Next, I wrote a work in one month. Double speed. The concentration should be halved.

I sent it to Furuta. A reply came two hours later:

``Still intense!''

I couldn't believe it. I felt like it had only a quarter of the intensity of the work that had led me to meet Furuta, yet it was still intense... I even suspected that even though I'd halved the concentration, Furuta had spent twice as long reading it, creating the illusion that nothing had changed.

I thought maybe the problem was the kind of things I read, so I read various ``light'' works. They weren't to my taste at all, but I could easily analyze what made them interesting.

I incorporated that lightness and spent two months writing another new work.

How about this, Furuta!

``That's not what I meant... I want you to write without lowering your 'deviation value.' In fact, isn't this more boring than before?''

...What the hell am I supposed to do! What's a ``deviation value''! What do you want me to write! I hated Furuta. I even wrote him an email. He ignored it.

I was bedridden for about a month---


\subsection*{5}

With no results to show for it, autumn came. I didn't even notice when summer began or ended. The air in my room was always kept at a constant temperature by the air conditioner.

I kept in touch with Shimizu intermittently. He had moved to Soma City and become a fisherman. He'd moved near Kobayashi Koyomi. Because of reputation damage, Fukushima seafood wasn't selling, and he was struggling.

``Even if it's the same bonito caught in the same waters, the port where it's landed determines its origin, so they go out of their way to land it at Kesennuma in Miyagi. Just because it's from Fukushima, it won't sell. It's ridiculous.''

That was quite harsh language for Shimizu. He must have had his own anger.

At the end of September, I received a contact from an unexpected person.

It was Yuzuki.

``Get a passport. Get one right now.''

After more than three years of silence, this came out of nowhere. Between my father, Furuta, and Yuzuki, it seemed I was limited to sending only a few bytes of data.

But I was also in a state of mind where I didn't want to write anything other than novels, so I just replied, ``OK.'' And I actually went out to fill out the paperwork and get a passport.

When I stepped outside after so long, I felt dizzy. Even though I was in my hometown, I felt like a foreigner. I'd forgotten how to behave normally, and I worried that I might seem suspicious.

I submitted the documents at the passport counter of the Koriyama Joint Government Building. Faced with the female official, I was at a loss. I had forgotten how to make my voice come out.


\subsection*{6}

In mid-October, Yuzuki contacted me again.

``Get ready for a trip!''

``OK.''

I was obedient. I wasn't in a position to argue anyway. I was just an idle person with nothing to do but write novels. But five days later, the message at 6 AM on October 13 took me by surprise.

``Come to Warsaw! My performance is on the 16th!''

I was stunned.---No, it probably meant I was slower than Yuzuki had imagined. After all, many people were paying attention to her activities...

Yuzuki was competing in the Frédéric Chopin International Piano Competition.\footnote{The International Chopin Piano Competition is one of the most prestigious piano competitions in the world, held every five years in Warsaw, Poland.}

The competition was introduced on the website of the Tokyo International Arts Association as follows:

``Held every five years in Warsaw, the capital of Poland and Chopin's homeland, for about three weeks around October 17, the anniversary of Chopin's death. It is the oldest of the currently continuing international music competitions, and its prize winners are world-renowned musicians. The required pieces are exclusively Chopin's works, including études, sonatas, fantasies, waltzes, nocturnes, concertos, and more. It is said to be one of the world's most authoritative competitions, considered one of the three major international competitions along with the Queen Elisabeth Competition and the Tchaikovsky International Competition. It is a gateway for pianists aiming for the world stage.''

``Amazing,'' I murmured involuntarily. I remembered the day I first met Yuzuki, when she made me listen to Maurizio Pollini's CD. He won the sixth competition in 1960 at the age of eighteen. The pianists I'd naturally encountered through my time with Yuzuki---Ashkenazy, Argerich, and even Tanaka Kiyoko-sensei---were all prize winners of the same competition.

And Yuzuki was challenging the 17th competition at the same age as Pollini when he won.

It was as if fate had willed it---

I looked it up and it seemed Yuzuki had sent me ``Come to Warsaw!'' right after the announcement of the second stage results. If I left immediately, I'd make it in time for the third stage.

I immediately booked a flight, took a local train from Koriyama Station for four hours to Tokyo, went to Hamamatsucho, and took the monorail to Haneda Airport. Lost and dizzy from the crowds, I managed to board the plane at 9 PM.

``I'm on the plane,'' I messaged Yuzuki.

``Good boy. Thank you. What time do you arrive?''

``Around 6 PM your time. Warsaw airport.''

``OK. I'll pick you up.''

The plane first flew to Kansai International Airport, then transferred and flew for ten and a half hours to Helsinki-Vantaa Airport in Finland. After another transfer, I finally arrived at Warsaw Chopin Airport.


\subsection*{7}

As soon as I got off the plane, I couldn't help murmuring, ``It's cold.''

The average temperature in Fukushima in October is 15°C, while Warsaw's average is around 8°C. I hadn't taken the climate into account and came wearing the same clothes.

Feeling somewhat pathetic, I headed for the gate.

I would see Yuzuki for the first time in three and a half years---just thinking about it made my heart race.

I passed through the gate. A blond man walking in front of me embraced a red-haired woman, apparently his girlfriend waiting for him, and kissed her. I searched for Yuzuki---

It was as if an electric current had shot through me.

I found Yuzuki.

I was captivated in an instant.

She had become so beautiful.

Her long black hair still flowed smoothly, and her lightly made-up face was more beautiful than ever. Her facial lines had sharpened, her almond-shaped eyes were more captivating, and her cherry-colored lips looked more alluring. Earrings sparkled in her ears. She wore a striped knit top, a white coat, slim black pants, and heels. A vivid young green scarf was wrapped around her neck.

She struck me as a very sophisticated woman. I, on the other hand, could only be described as wretched.

At first, Yuzuki smiled when she saw me, but as I got closer, her eyes widened.

``Whoa! You came in that thin clothes? And your complexion is way too pale!''

It was a first greeting that didn't feel like we hadn't seen each other in three and a half years. I felt like the whole trip was ruined.

---We walked through the airport side by side, chatting about various things. In an open, glass-walled area, there was a grand piano. A man with a thick beard and a bodybuilder's build was playing Chopin's Waltz No. 13. The sound was unexpectedly bright and sweet.

``Ah---'' I realized. ``I only brought Japanese yen. I need to exchange money.''

The Polish currency is the zloty. One zloty is about 31 yen. Yuzuki said:

``The exchange rate at the airport is bad, so it's better to go to a money changer in the city! I'll cover you for now, and you can pay me back later all at once.''

I checked later and the airport exchange rate was 1 zloty for 40 yen.

We went into a clothing store in the airport, and Yuzuki picked out a coat that would suit me. She said that dark colors would make my paleness stand out, so gray would be good. Seeing her pay with a credit card, I felt ashamed of my own lack of social competence.

We left the airport and took a taxi. I finally noticed that Yuzuki was wearing perfume.

``How was my playing?''

Yuzuki asked. I tilted my head.

``Playing? Which playing?''

``Huh! Didn't you watch the YouTube videos?''

Apparently, during the 17th Chopin International Piano Competition, videos of the participants were streamed in real-time, and viewers shared their impressions on Facebook. While I was on a desert island, the world's communication methods had advanced.

``I didn't watch the archives either. Actually, I didn't even know you were in the competition.''

``...So you came right away just because you got a message, without knowing anything?''

``That's about it.''

Yuzuki looked dumbfounded.

``What kind of life do you have to live to end up like that?''

I briefly explained my life after high school. Not that it took much explaining; it had always been simple. I thought she might scold me for my hopeless lifestyle.---But Yuzuki just sighed.

``I always thought you'd become a novelist someday, Yakumo-kun.''

``Huh? Really?''

``Actually, I thought you could only become a novelist.''

``Why?''

``Anyone can tell after spending a little time with you.''

I didn't really understand, but I thought perhaps it was true. I'd started writing novels because Shimizu told me to. He might have thought the same as Yuzuki.

Through the taxi window, the streets of Warsaw, already approaching night, passed by. The airport area was quiet, but as we neared the city center, it gradually grew livelier.


\subsection*{8}

Poland, located between Germany and Belarus, at the center of Europe, is called ``the heart of Europe.'' And at the heart of that lies the capital, Warsaw.

The Vistula River, Poland's longest river at 1,047 kilometers, originates in the southern Beskid Mountains and flows into the Baltic Sea, splitting Warsaw in two.

The left bank was where we did everything; we never once crossed to the right.

The taxi stopped in front of a hotel in the area between the Warsaw Philharmonic, the venue for the Chopin International Piano Competition, and the Chopin Music Academy. The exterior looked almost like ones in Japan---but the interior had detailed, stylish touches.

It cost about 200 zloty per night. The room next to Yuzuki's happened to be empty, so I checked in there. When I put my luggage in the room, it was already dark outside.

Yuzuki came to my room and sat on the bed. A soft, pleasant fragrance drifted. I sank into the blue leather sofa.---Instantly, exhaustion and drowsiness washed over me. I'd traveled a total of twenty-seven hours from Fukushima. I turned back my watch, which showed three o'clock, by eight hours, setting it to 7 PM.

``...By the way, what about your parents? Aren't they here to support you?''

Yuzuki looked down slightly.

``They seem to be here, but I haven't seen them even once. If I see them, my playing gets thrown off.''

It seemed Yuzuki and Ranako's conflict was still ongoing.

We had dinner together at the hotel restaurant. We ate pierogi, a Polish traditional dish like dumplings, borscht, bread called ``frytki,'' and many other dishes whose names I didn't know... They were all delicious, but the feeling of being on a plane hadn't worn off, so I couldn't eat much. Yuzuki, on the other hand, didn't lose her appetite from nerves and ate normally.

As we talked about each other's lives since we'd last met, time flew by. I wanted to keep talking forever, but we cut it short early to prepare for the next day, and I went to bed in my room.

The scent of Yuzuki lingered on the bed. It was the scent of freesia.


\subsection*{9}

October 15---We had breakfast at the hotel restaurant. I drank orange juice and glanced sideways at Yuzuki eating bread with honey. Her makeup seemed a bit lighter than yesterday. She was so naturally beautiful that it seemed she didn't even need makeup.

``Yuzuki, you're performing tomorrow, right? Will you practice on a piano at the Chopin Academy?''

``I'll do that in the morning. Let's have lunch together, and then we'll go to the Old Town.''

Even though I was in Poland, I holed up in my hotel room writing novels and researching the Warsaw Old Town again.

When Yuzuki returned, we ate at a nearby restaurant and then walked north along Krakowskie Przedmieście. It was a wide, beautiful cobblestone street. Buildings in Renaissance and Neoclassical styles lined both sides. The streetlights were shaped like lily-of-the-valley flowers. We saw a statue of Copernicus in front of the Polish Academy of Sciences and a statue of Poniatowski in front of the Presidential Palace. I wanted to see the famous Chopin statue, but it seemed to be in the opposite direction.

After walking about twenty minutes, we arrived at our destination.

The Warsaw Old Town---a UNESCO World Heritage site, this place is a new old town.

This place once disappeared from the world and was revived.

---On September 1, 1939, Germany and its ally, independent Slovakia, invaded Poland. On the 3rd, Poland's allies Britain and France declared war on Germany, beginning the Second World War. On the 17th, the Soviet Union, which had signed a non-aggression pact with Germany, also joined the invasion, and by October 6, all of Poland was occupied by Germany and the Soviet Union.

Then, on June 22, 1944---when the Soviet army launched Operation Bagration,\footnote{Operation Bagration (1944) was the Soviet Union's major offensive against Nazi Germany on the Eastern Front.} the largest counteroffensive against German forces in Belarus, the German army began to retreat repeatedly. By this time, the Soviet Union had become Germany's enemy. As they pushed back the Germans and the occupied territories expanded across eastern Poland, the Soviet Union called for an uprising by the Polish resistance.

And at exactly 5 PM on August 1---the Polish Home Army, including women and children, rose up in armed revolt against the German forces. This was the Warsaw Uprising.\footnote{The Warsaw Uprising (1944) was a major World War II operation by the Polish resistance Home Army to liberate Warsaw from Nazi occupation. It was suppressed by the Germans after 63 days.}

However, the Soviet forces that had encouraged the uprising betrayed them and did not support the citizens of Warsaw. In an overwhelmingly disadvantageous situation, the citizens of Warsaw continued their desperate resistance, but after sixty-three days, they were finally defeated. The number of people massacred is estimated to have been between 180,000 and 250,000.

Afterward, Warsaw was systematically destroyed by punitive German attacks.

It was a dark history, almost smelling of death---

After the war, the Soviet Union, which had replaced Nazi Germany in ruling Poland, planned to rebuild Warsaw into a Soviet-style city based on socialism. The citizens, upon learning this, united in opposition. With the belief that ``a cityscape destroyed with intent and purpose must be restored with intent and purpose,'' and the principle that ``the restoration of what was lost is a responsibility to the future,'' they chose to rebuild the Old Town exactly as it had been before its destruction.

It is said that all citizens participated in this enormous undertaking, and the expenses were covered by their donations.

The basis for the reconstruction was the realistic landscape paintings of Bernardo Bellotto,\footnote{Bernardo Bellotto (1721–1780) was an Italian painter known for his detailed cityscapes of Warsaw, which were crucial in the post-war reconstruction.} an 18th-century court painter, and over 35,000 detailed drawings of buildings and other structures left by students of the Warsaw Polytechnic Institute.---In 1939, teachers who foresaw the Nazi occupation had collaborated with students to draw and preserve these records of Warsaw's cityscape, risking their lives to protect them.

Based on these, Warsaw's streets were reproduced ``with fidelity down to every crack in the walls.''

And in 1980, the Warsaw Old Town became the first World Heritage site to be recognized for ``the efforts of people in reconstructing and maintaining from destruction.''


\subsection*{10}

Together with Yuzuki, I walked through the beautifully reconstructed medieval streets. Buildings in various warm hues, from Gothic to Neoclassical, lined the streets.

As I walked, I felt a profound sense of awe. What wonderful people the Poles were, I thought. Polish history is a turbulent one. They had been invaded many times, lost countless lives, lost their homeland, and each time they rose again like a phoenix.

I found bullet holes here and there in the streets. During the reconstruction, as much as possible of the original bricks and fragments were reused, leaving traces of the war.

I stopped involuntarily and touched one of the traces. At that moment---I felt I heard something.

The sound of pencils running across paper. The sound of drawing 35,000 blueprints. Thinking of their feelings made my heart feel tight. The sadness of knowing their homeland was about to be destroyed, the frustration, the deep love for Warsaw. And the sound of turning those 35,000 blueprints. The sound of millions, tens of millions of bricks being painstakingly laid... It was as if, like countless bouquets, they gently covered the terrible gunfire and the sad faces of the dead.

That is why, when I saw the blank space of the bullet holes, I didn't feel pain. It was a love so overflowing it made me feel ashamed, like when I saw Yuzuki enter the Abukuma River after the earthquake.

On the cobblestones of Warsaw lingers the melancholy of a dark past. But even more than that, the present of the people is beautiful and bright. That is why Chopin's music, with its ``cannons hidden among flowers,'' rings most beautifully in the sky of this city---

I gazed at the bullet holes and, standing there, found myself crying without realizing it. Through writing novels, I had learned that ordinary people don't cry in such places, and that those who do are considered strange. But I couldn't help crying.

``...Yakumo-kun, you haven't changed.''

Yuzuki said that, and gently came closer to me.


\subsection*{11}

October 16, 10 AM. Warsaw Philharmonic.\footnote{The Warsaw Philharmonic is the main concert hall of the National Philharmonic in Warsaw, a major venue for classical music in Poland.}

Yuzuki was the first performer of the day. I sat in the ground-floor audience seats. The second-floor seats should have been lined with the world's greatest pianists as judges (including Martha Argerich!), but from my seat, I couldn't see them due to the angle.

The program and other details were announced in both English and Polish.

Yuzuki appeared on the wooden-ceilinged stage---wearing a brilliant red dress. The contrast between her glossy black hair and white skin was beautiful. The audience greeted her with thunderous applause. Having already passed the first and second stages, they already favored Yuzuki.

Yuzuki smiled, made an elegant bow, and sat down at the piano.

The piano was a YAMAHA CFX. Contestants could choose from four brands---Yamaha, Kawai, Steinway, and Fazioli---and there was competition among manufacturers and tuners.

The hall fell silent---

Yuzuki took a deep breath.

Tension filled the air.

Yuzuki's fingers began to move---

\emph{Three Mazurkas, Op. 59}---Polonaises and mazurkas are Polish folk dances, but while the polonaise spread mainly among the nobility, the mazurka was beloved by commoners. It is said that Chopin wrote mazurkas like a diary, almost daily.

Op. 59 was composed between 1844 and 1845. At that time, in addition to his usual poor health, Chopin was suffering from conflicts with Maurice, the son of his lover George Sand.\footnote{George Sand (1804–1876) was a French novelist and Chopin's romantic partner for nearly a decade.}

The first movement begins with a melancholy melody, the second is full of vitality, and the third expresses intense emotion like anger at the start before transitioning to a sweet melody and ending on a positive note---

Since it was four years before his death, there is an undercurrent of death in it.

It had been nearly ten years since I'd heard Yuzuki perform live.

The moment I heard the sound, goosebumps ran down my back.

She had improved dramatically---

While inheriting the beauty, purity, mellowness, and transparency of tone from Tanaka Kiyoko-sensei, Yuzuki's own rich emotion and sense of color added depth. The tones seemed to paint a vivid life. Just as life changes, Op. 59 constantly changed color. Yuzuki expressed it with remarkable richness. There was even an aromatic flavor like fine wine. This must be what Yuzuki had called maturity. I thought of Yuzuki's life in Italy. She must have spent very fulfilling days. The fruit of each accumulated day had ripened and appeared in her tone. It also seemed to revive some nostalgic memories. As if the seasons were rewinding, and cherry blossom petals that had fallen to the ground were floating back up into the sky. And yet, something indescribable, precious, like a blue pearl, seemed to be forming in my heart.

\emph{ZAL}---Yuzuki had expressed it brilliantly.

---When the nearly sixty-minute performance ended, thunderous applause erupted.

Yuzuki smiled. She rose from her chair, bowed, and disappeared into the wings.

Her performance was overwhelming.


\subsection*{12}

October 17---This day is the anniversary of Chopin's death.

The announcement came in the morning, confirming Yuzuki's place in the final. After the announcement, I watched from a distance as Yuzuki was surrounded by the press. It felt like we lived in different worlds.

I believed Yuzuki would undoubtedly win. She would stand shoulder to shoulder with the top pianists she had always aimed for. And many people would hear her play.---That was a very good thing.

We attended the Chopin Requiem Mozart\footnote{The Chopin Requiem Mozart is a memorial concert held in Warsaw on the anniversary of Chopin's death, featuring Mozart's Requiem.} at the Holy Cross Church, a memorial concert for Chopin. At 8 PM, flowers were offered at the pillar housing Chopin's heart. Chopin had never returned to his homeland in life; after his death, only his heart was secretly brought back by his sister. During World War II, Chopin's heart was evacuated and protected by the people. The Holy Cross Church was also destroyed during the Warsaw Uprising and rebuilt.

The pillar is inscribed with Matthew 6:21: ``For where your treasure is, there your heart will be also.'' ---``Do not store up for yourselves treasures on earth, where moth and rust destroy, and where thieves break in and steal. But store up for yourselves treasures in heaven, where moth and rust do not destroy, and where thieves do not break in and steal. For where your treasure is, there your heart will be also.''

Chopin's heart was in the heaven of his homeland. On his deathbed, he is said to have expressed his wish: ``At my funeral, I want Mozart's Requiem to be played.''

Mozart's Requiem resounded through the high ceiling of the Holy Cross Church---

We listened to the beautiful performance from not very good seats. Yuzuki had politely declined better seats, saying that as Japanese, we couldn't possibly sit in better seats than the Polish people.

The sorrowful and solemn singing and music moved us deeply.

``On that tear-filled day,
when all people rise again from the ashes,
to be judged for their sins,
O God, give us mercy.
Merciful Jesus, Lord,
grant them eternal rest. Amen.''

The orchestra performing before the altar was brightly lit by chandeliers and lights. The audience behind them was shrouded in dim light. The faces of people with mournfully lowered eyes, sometimes even showing tears, were like a painting by Rembrandt.

Yuzuki also lowered her long eyelashes and quietly offered a silent prayer.


\subsection*{13}

From October 18, the finals of the competition began.

I listened to performances sometimes with Yuzuki, sometimes alone. The pieces were either Piano Concerto No. 1, Op. 11 or Piano Concerto No. 2, Op. 21.\footnote{Chopin's Piano Concertos Nos. 1 and 2 are among the most famous works in the piano repertoire, both composed in his youth.} They were performed with the Warsaw National Philharmonic Orchestra, conducted by Jacek Kaspszyk.\footnote{Jacek Kaspszyk (b. 1952) is a Polish conductor who has worked with major orchestras across Europe.} Since almost all finalists chose No. 1, I heard the same piece many times. But they were all wonderful, as expected of those who had advanced; I was completely absorbed in each of the roughly forty-minute performances.

I actually happened to encounter Martha Argerich inside the venue and was able to get her autograph.

``Thank you.''

I said it, putting various feelings into those words. She smiled.

Yuzuki spent most of her time alone, seemingly building her focus.

And finally, on October 20, at 7:50 PM---on the final day, the last performer was Yuzuki.

She appeared wearing a red dress. Loud applause greeted her. Yuzuki shook hands with the concertmaster and the principal second violinist, bowed to the audience, and sat down at the piano.

And then---she took a deep breath.

She met the conductor's eyes and nodded.

The performance began---

Yuzuki had chosen Piano Concerto No. 1, Op. 11---the exposition is played by the orchestra, and the piano doesn't enter for some time. Yuzuki quietly lowered her eyes, gazing at the keyboard. I was the one who was about to fall apart from tension. After about four and a half minutes, the piano finally sounded its first note.

A brilliant fortissimo chord---

From there, it was Yuzuki's world. A rich bass and high notes like polished glass beads beautifully depicted Chopin's poetic sentiment. I was captivated by Yuzuki's expression, her finger movements, her pedaling, her every gesture---so much so that I almost felt the orchestra had receded into the background, leaving only Yuzuki.

A sweet melody flowed.

The piano was singing---

I felt as if Yuzuki and the piano were one inseparable existence.

A beautiful dream seemed to be playing music.

Fate was sounding---

Like when I first met Yuzuki. No, even more strongly, I felt it. For some reason, tears welled in my eyes. It was that beautiful a performance. Without a doubt, Yuzuki would win. I was certain of it.

---It was about four minutes before the end of the first movement.

A dramatic, beautiful passage was being played, like running down a steep slope, like something built collapsing, like dominoes falling.

The piano sound stopped---

A murmur arose. A scream was heard. The orchestra's performance stopped as if stumbling.

I opened my closed eyes---

Yuzuki was staring at her own hand, her almond-shaped eyes wide open.

Her left ring finger---

It had snapped between the first and second joints and was lying on the keyboard.

Yuzuki's face showed that she didn't understand what had happened. Then, her face quickly turned pale as a dead person's. Her beautiful face distorted as if cracking.

Yuzuki let out a terrible scream. The conductor moved. He put the finger that had fallen on the keyboard into his breast pocket, took Yuzuki by the shoulder, and led her into the wings. Great confusion and commotion arose. Some people covered their mouths and cried. I forgot even to blink, sitting there stunned and paralyzed.

In that instant, I had seen the cross-section of Yuzuki's ring finger.

---It was pure white.

I understood what that meant.

Yuzuki, too, had contracted salt disease---

\fullpageimage[]{11.jpg}
% ==============================
%        Chapter 7
% ==============================
\clearpage
\thispagestyle{plain}
\phantomsection
\begin{center}
    \large\bfseries Chapter 7
\end{center}
\addcontentsline{toc}{chapter}{Chapter 7}
\label{sec:ch7}
\vspace{1cm}

\subsection*{1}

Yuzuki was taken by ambulance to a hospital in Warsaw.

I rode in the ambulance with Yuzuki and two people who seemed to be event staff. While Yuzuki was being examined, I stood around aimlessly, staring blankly at a statue of Christ in the hospital. Then Ranako and Sosuke appeared. They must have followed the ambulance.

We exchanged bewildered looks and fell silent. We had all lost our words.

---Eventually, Yuzuki returned. She was still crying violently.

``Yakumo-kun...! Yakumo-kun...!''

I embraced her, still wearing her dress with a coat over it. Over her shoulder, I saw Ranako and Sosuke. They wore expressions of hurt and devastation, indescribable. Yuzuki completely ignored her parents and walked alone toward the exit.---Their conflict must have deepened while I wasn't there. A little confused, I spoke to her parents.

``...I think it's salt disease.''

Ranako covered her mouth. Sosuke gaped. I ran after Yuzuki.


\subsection*{2}

Yuzuki sat on the hotel bed and cried endlessly.

I sat beside her, stroking her back.

Neither of us spoke. It was too sad, too sudden. Yuzuki could no longer play the piano, and it was certain she would turn to salt and die within about a year. I was filled with disbelief. It all seemed like a cruel joke. Yuzuki cried until 3 AM, then suddenly collapsed as if a string had snapped. I panicked, but her breathing was normal. I laid her on the bed in her red dress, and slumped onto the sofa.

I checked online; the news of Yuzuki's accident was causing a huge stir. Her finger falling off had been broadcast live, and many people had witnessed it. Words resembling ``salt disease'' in various languages were flying across communities. This rare disease, barely known to anyone, had become explosively famous thanks to Yuzuki.

Of course, the official video archives cut Yuzuki's performance, but regular viewers had clipped just the scene of her finger falling off and shared it online. Seeing it filled me with both pain and intense anger.

I looked at Yuzuki's finger, wrapped in cloth and left on the table. The salinification was progressing rapidly; about seventy percent had already turned into fine salt. The parts that remained as flesh looked grotesquely strange. Frightened, I quickly rewrapped it. I thought of my mother.

I was worried Yuzuki might attempt suicide, so I fell into a light sleep on the sofa.


\subsection*{3}

The next day, the results of the Chopin International Competition were announced.

The winner was showered with praise and would begin a brilliant career as a pianist---

We didn't watch any of it. We just stayed in the dim room. We only learned much later that Yuzuki had received rave reviews from the judges and audience and had been awarded a special prize.

Yuzuki was quiet as death. She neither grieved nor laughed. She blinked at regular intervals, occasionally glancing at her left hand.

We took the evening flight to Italy together. A direct flight of two hours. Unable to bear the silence, I tried talking to her several times, but Yuzuki only gave vague, absent-minded replies.

From Milan Malpensa Airport, we took a taxi to Yuzuki's apartment.

It was a beautiful white-walled apartment with an impressive, stylish green railing on the balcony.

It was surprisingly large for a single person, with even a room containing a grand piano. Yuzuki took off her coat and went to that room. On the piano sat that familiar Anpanman doll from long ago. Yuzuki gazed at it quietly, then looked down at the piano keyboard and burst into tears as if a dam had broken. She cried for three days and nights---

At a loss, I still managed to move for Yuzuki's sake. I shopped in a foreign city, confused, cooked using recipes I found online, and brought the food to her crying room.

``Sorry, I don't want to eat...''

For two full days, Yuzuki ate nothing. Even after that much time, she shed tears of thick sadness, like dissolved blue paint. I feared the room would become a deep blue sea, swallow the piano, and swallow Yuzuki too.

I thought of the void left by Yuzuki's lost finger and felt terrible pain. To fill that void, I even ate all of Yuzuki's leftover food, forcing it down until I felt sick.

On the morning of the third day, Yuzuki ate a few slices of orange. Dark circles had formed under her eyes, as if stained by the color of grief. She only ate a few pieces.

``...Thank you, Yakumo-kun...'' she said, and began crying again.

Late at night---I woke suddenly and heard a faint \textit{pon, pon} sound of a piano.

Yuzuki was striking the keyboard. Half-asleep, I imagined an African elephant. Its mother had died somehow, lying on the ground, separated from the herd. The baby elephant, perhaps unable to understand death, nudged the area around its mother's backside with its trunk and kicked the carcass with its front feet---\textit{ton, ton}---as if saying, ``Wake up.''

The piano sounds were that sad.

As my mind cleared and Yuzuki's sobbing pressed against my ears, I felt a sense of hatred mixed into the piano notes. It was hatred born from love. Yuzuki had devoted her entire life to the piano without a single regret, finding happiness in it. That love had suddenly been betrayed and turned into hatred. I could even sense a desire to smash the keyboard. But because she had always played with such care and devotion, she couldn't bring herself to destroy it entirely, and the keyboard responded with sorrowful notes. Those notes sounded like Yuzuki's own voice.

``Don't leave me. Don't leave me alone.''

That was what she seemed to be saying, wordlessly.


\subsection*{4}

On the morning of the fourth day, I woke to a pleasant aroma and found dishes set on the dining table.

I looked toward the kitchen; Yuzuki was making the familiar rhythmic \textit{sutoton-toto} sound with a knife. Her beautiful black hair was tied in a ponytail, and she wore a green apron as she added another dish to the table. I caught a glimpse of her slender, white neck.

Half-dazed, I watched her work.

``Good morning, Yakumo-kun.''

Yuzuki saw me and smiled sweetly. The dark circles under her eyes had disappeared.

``Ah... yeah... good morning...''

Though confused, I followed Yuzuki's lead and sat at the table. Golden-brown toast, consommé soup, caprese salad, bacon and eggs, orange juice---the table was so bright and vivid, it was hard to believe it was the same table that had been drowning in blue sadness just days before.

``Let's eat---''

We ate the simple yet delicious food Yuzuki had made. She ate heartily too.

She asked me what I'd been doing while she was crying all that time. A man with no life experience, I'd made quite a few mistakes. Yuzuki's eyes widened and she laughed, ``Ahaha,'' gulping down her orange juice. After the meal, she cleared the dishes with a \textit{sukoton-pappa} sound and even washed them with a \textit{shassha-shassha-pirupiru-kachanko-tonton}.

``Yakumo-kun, you want to be a novelist, so you need to write proper novels.''

With that, she pushed me into the corner room and made me tap away at the keyboard, while she vacuumed with a weird \textit{pururu-woooon} sound. I heard a \textit{gorigori} scraping sound, then a delicious aroma wafted in, the door opened with a click, and Yuzuki appeared saying, ``Good work,'' setting a coffee cup beside me. She had been grinding beans. I took a sip. It was wonderful. I started tapping again, and a \textit{ping} announced an email from Furuta.

``Yakumo-kun, you've been quiet lately. How are you doing?''

I recalled my recent whirlwind life, then shook it all off.

``I'm writing novels. As usual.''

``I think you should try writing a harmless romantic comedy once.''

As always, it came out of nowhere. A bit irritated, I typed:

``What do you mean?''

``Clear your head, write what you love with abandon, and maybe you'll find a way.''

``That's as vague as a fortune teller. Besides, I don't like romantic comedies at all. In fact, I can't think of something I like right away.''

``What? I have lots of things I like!''

``What do you like?''

``What I like most is boobs! Boobs!!''

Had this guy really lived twice as long as I had...?


\subsection*{5}

In the evening, Yuzuki and I watched a movie together: \textit{Cinema Paradiso}.

The next paragraph contains spoilers, so if you haven't seen it, please do so before reading on. It's a masterpiece. Personally, I recommend the theatrical version over the director's cut.

---The final scene. The protagonist, Toto, watches a montage of love scenes cut from old films, laughing and crying with nostalgia. This beautiful scene moved both Yuzuki and me to tears. The movie ended and the credits rolled. The dark room was lit by the pale light of the TV. We happened to meet each other's eyes in that light.

The dark circles had returned under Yuzuki's eyes. She had been covering them with makeup.

Noticing me looking, Yuzuki said:

``You know, Tanaka Kiyoko-sensei also contracted collagen disease in her thirties and could no longer play the piano. She must have been terribly frustrated and suffered pain like being cut by a blade. But even so, she went on to train many students and lived with fortitude until the end. I decided to follow her example. I stopped grieving and decided to live.''

I was once again amazed by Yuzuki's strength. She had already begun to fill the terrifying silence left by the loss of the piano she loved so much with the sounds of everyday life.

Overwhelmed with respect, I fell silent. Then Yuzuki spoke:

``...Hey, Yakumo-kun, have you ever kissed anyone?''

Startled, I looked at Yuzuki. She was staring fixedly at the featureless menu screen after the credits had finished. The kiss scene at the end must have prompted her question.

``...Uh... no.''

``...Me neither.''

We met each other's eyes again. Yuzuki's face seemed a little tense. I asked:

``...Do you want to?''

``...Not especially...''

``...Then we don't have to, right?''

Yuzuki's left eyebrow twitched. She said:

``...I suppose. I'm tired, so I'm going to bed. Good night.''

She stood up quickly, went to the bathroom, brushed her teeth with a \textit{shassha}, and climbed into bed without even taking a bath.


\subsection*{6}

We continued our strange cohabitation in Italy.

In early December, a film director came. Daniel Miller\footnote{Daniel Miller is a fictional character, likely a pastiche of a stereotypical eccentric filmmaker.}---he had an appearance as interesting as a manga character. Short and stocky, with messy chestnut hair that connected his sideburns to a magnificent beard. His square glasses had red frames. He wore a gray stylish jacket over a Superman T-shirt.

His assistant, perhaps a secretary, was taller than the director, a long-legged Latin beauty in heels. Her black hair was cut in a strict bob, her nose was high like a witch's, and her abnormally long lashes, combined with eyeshadow, made her look a bit like Cleopatra.

While Yuzuki conversed fluently in English, I prepared tea---just dunking a green tea bag you can find in any Japanese supermarket.

Miller drank it with apparent relish, saying, ``I love Japanese tea!'' Cleopatra only gave a thin smile and nodded in agreement.

``Is he your boyfriend?'' Miller asked Yuzuki, gesturing toward me.

``No, he's a friend.''

She emphasized the word ``FRIEND'' unnecessarily, it seemed.

At this point, Yuzuki had also lost her right pinky and middle fingers, and her left thumb. The salt disease was steadily progressing. So holding the teacup required some effort.

Miller spoke extremely fast and had a slight lisp, so I couldn't follow the conversation at all. I just sat across from Cleopatra, who looked like an Egyptian mural, sipping Japanese tea while Yuzuki and the director talked intently.

About two hours later, the director and his party left.

``Sayonara~!'' Miller waved with a cheerful smile. Cleopatra hadn't spoken a single word the entire time.

``---So, what was that all about?''

I asked Yuzuki after they left.

``He wants to make a movie about me.''

``What, that's amazing!'' I said, surprised. ``So, what did you say?''

``I said I'd think about it.''

We watched several of Miller's films together. As his Superman T-shirt suggested, his basic aesthetic was American comics, but he threw in Japanese subculture and disguised it with a faux-Spielberg flavor. Every single one was like that.

``I hate to say it...'' I scratched my cheek. ``It feels third-rate...''

``Pretending to be original while actually being unoriginal. I think he lacks aesthetics and philosophy.''

``You're pretty harsh.''

``Because that man said, 'Film is my soul!' with nothing but smooth talk.''

``Yuzuki, are you angry about something?''

``I could see right through his scheme...''

Yuzuki sighed, folded her damaged hands like a puzzle, and said:

``...He just wants to use me as a stepping stone. He wants to film my suffering from the salt disease, my overcoming it, and then---my death. If you were crying beside me, Yakumo-kun, it would be perfect. A perfect tearjerker. Since I'm famous from the Chopin competition, the film would get a lot of attention. It might make the third-rate director Miller famous overnight...''

``I see...''

``And most people wouldn't notice Miller's ulterior motives. People drawn to the theater by the tagline 'It'll make you cry' would cry exactly as Miller intends, feel refreshed, sleep soundly, and forget everything the next day. They'd casually consume my death.---I don't want that. I wasn't born for that. I didn't laugh and cry for that. I'm not so naive as to die just to turn a third-rate film into a second-rate one---''

Amidst the peaceful days, Yuzuki's passion suddenly came into view.

---She didn't want to be consumed.

I thought that was the steadfast feeling within Yuzuki. I understood that feeling all too well. The earthquake and the CD cover incident had given birth to this rebellious spirit. The earthquake had been consumed by some people. There were volunteers and donors, but also those who used it for profit and attention. And Yuzuki had been made a tool in that, through the CD cover.

So Yuzuki couldn't forgive the easy consumption of others' misfortune and death, nor did she want to be consumed.

``That man would make people with the same disease sad, and make films that appear gentle on the surface but hurt people deeply inside. No matter how much it sold, it would be a shameful failure---''

``...That's harsh. But I kind of wanted to see the film...''

I said carefully, and Yuzuki's stern expression softened.

``Then you shoot it, Yakumo-kun---''


\subsection*{7}

I decided to leave a video record using Yuzuki's video camera.

I filmed Yuzuki in everyday moments. Whenever I approached, she would smile and wave at me. She was adorable. Yuzuki looked so good on camera from every angle that I almost believed I had good skills myself.

It was strange how just filming Yuzuki's profile as she walked with a smile through the beautiful streets of Milan could become like a work of art. There are people who can capture hearts just by walking, without having to painstakingly write words like I do.

In the footage, Yuzuki seemed to grow more beautiful. Even more beautiful than before she contracted the disease. Maybe it was the beauty of death. Like a sparkler that blazes brightest just before it falls.

Yuzuki couldn't pick up the bright red cherries she bought at the market. When she managed to lift the fruit with her remaining fingers, she popped it into her mouth and smiled a little shyly, as if to cover up her clumsiness. The contrast between the white crystallized stumps and the red cherries was so vivid that they looked like blood, and it made me uneasy.

``How did it feel filming, Yakumo-kun? Did it make you want to become a director?''

Yuzuki tilted her head slightly and asked. I thought for a moment.

``It was a lot more fun than I expected.---But I still want to write novels. I feel like what I want to express lies in words.''

``I see...'' Yuzuki's face became serious. ``What kind of novels do you want to write?''

I thought for a while.

``I want to save---''

I was about to say ``Yuzuki,'' but I couldn't. I wanted to write a novel that could save Yuzuki from despair. I wanted to write a novel that could heal the sadness of Yuzuki losing the piano.---But I knew that was impossible. Yuzuki's despair and sadness were too deep, and I was too inexperienced. So I said, ``someone.''

``I want to write a novel that can help even a little someone who can only be saved by stories.''

``...That's naive. But I think it's very you, Yakumo-kun.''

Yuzuki said that and smiled gently.

``Then someday, write about me too. I want to help you help people, Yakumo-kun.''

``Okay, I will. I'll write about you.''

I said it carelessly, without thinking.


\subsection*{8}

I knew that Yuzuki cried alone in the middle of the night.

I heard her quiet sobs from her bedroom. I couldn't do anything about it. So I wrote a novel, making Yuzuki happy within the story. I knew it was meaningless, but I wrote it like a prayer: please, let Yuzuki be saved.

One night, I woke up and found the atmosphere different from usual. I didn't hear Yuzuki crying. I sensed a presence in the living room.

I slipped out of bed quietly and opened the door gently.

I saw Yuzuki's back.

She was sitting at a desk by the large window facing the street, illuminated by the orange light of a table lamp, her silhouette somehow nostalgic.

She started and turned toward me.

``Oh, Yakumo-kun, were you awake?''

As I approached, I saw a microscope on the desk. A very old one, made of faded brass, with a simple shape like a tube on a stand.

``Wow, what's this?''

``I found it in an antique shop. It felt right... so I bought it.''

Yuzuki shifted her weight and made room on her chair. When I sat down, our skin touched. Her body temperature was higher than I'd expected. I stared at the microscope's beautiful craftsmanship, then looked through the lens.

I saw translucent, square grains.

``Were you looking at salt crystals---?''

I was surprised and looked at Yuzuki. She wore a faint smile.

``...I felt like I had to see them.''

She turned back to the microscope. She brushed her hair behind her ear and opened her almond-shaped eyes wide.

Looking at her beautiful profile, I felt a chill.

Yuzuki was staring straight at her own death---

``It's beautiful---'' she said. ``Scary, but so beautiful... It's strange that I can become such a simple, beautiful crystal, when I'm so complex and ugly...''

The battlefield of Warsaw came to mind. Many people had been killed, reduced to ash. The ash formed a gray desert. The destination of all people is so simple from the start.

So simple, so quiet, even beautiful.

``But I think it's a lonely beauty.''

``But deep down, maybe everyone wants to dissolve into that lonely beauty. Into the stillness of a clear surface, without anger, hatred, sadness, or pain... Like drifting into a gentle nocturne, dozing off...''

Outside the window was midnight Milan. A deep darkness, like the bottom of the sea---

I wondered if Yuzuki was beginning to accept death. Whether she was dissolving into the night, becoming beautiful salt, little by little...

``I wish the afterlife were more lively. I wish all kinds of people, with all kinds of beauty and ugliness, could be happy there. I wish there were all kinds of salvation for all kinds of sorrow. And I wish they could all laugh together.''

``...That would be wonderful.''

Yuzuki smiled softly. Then she picked up the bottle filled with her own salted remains and tilted it. White salt trickled out, forming a small hill.

With her remaining right index finger, Yuzuki broke the hill and drew the shape of a flower. A sand painting---

Yuzuki smiled, and I smiled back. Then I smudged the flower and drew a whale.

``Oh, cute! You have an artistic touch, Yakumo-kun.''

``A flying whale.''

``When I die, I want that whale to come pick me up.''

``First class.''

``Hehe.''

Yuzuki laughed and leaned against me. I felt her warmth, her breath, the slight sway of her body. My chin brushed her soft hair, and I caught a faint pleasant scent. I gently held her slender shoulder. She didn't resist at all.

A bittersweet, sorrowful time passed---

``When I die---'' Yuzuki murmured. ``I want to die in the place I was born.''

Her words felt like a cold blade piercing my chest.

``...Yeah.''

``Let's go home, Yakumo-kun. Back to our homeland---''


\subsection*{9}

In February, we moved to Koriyama City, Fukushima Prefecture.

Perhaps because we'd moved from the Mediterranean coast, or because Yuzuki's death was approaching, Fukushima felt somehow dark. Standing on her homeland, Yuzuki's expression was surprisingly calm.

We lived in an apartment spacious enough for two. I handled most of the tedious paperwork and moving tasks. When we finished, I adjusted the angle of a wall clock with an orange cross-section design according to Yuzuki's instructions, and we laughed together.

Yuzuki no longer had any fingers left.

I began doing all the cooking. I didn't want Yuzuki to eat bad food, so I studied and spent time preparing elaborate dishes. Yuzuki ate them with great relish.

``Good~ Yakumo-kun, you have talent~'' Yuzuki beamed. ``Here, aah~''

She opened her mouth. I used chopsticks or a fork to feed her. It seemed like a lot of trouble, but Yuzuki looked so happy.

Yuzuki devised various ways to live without fingers. For example, she would wrap a rubber band around the back of her hand and use it as a substitute for grip. She would wedge a comb between her fingers to brush her hair, or a toothbrush to brush her teeth... But apparently it was still a bother, and she preferred me to do things for her. She even made me let her rest her head on my lap when brushing her teeth.

Washing her hair was also my job. Whenever I heard her call ``Yakumo-kun~,'' I took off my socks, rolled up my pant legs, and went into the bathroom. Yuzuki, wrapped tightly in a bath towel, would be sitting on a stool waiting. She wore rubber gloves to keep the salt disease's crystallized stumps from getting wet. I also sat on a stool and washed Yuzuki's hair---

``Is anywhere itchy, ma'am?''

``Hehe.''

I joked around, pretending to be a hairdresser, but honestly, I was nervous inside. Yuzuki had gained a little weight, and her white shoulders were smooth and full. The bath towel was a little tight around her body. When the hot water splashed on her hair, it became the color of a wet crow's wing, her nape even whiter, and the towel clung to her body, accentuating the curve of her waist.

I pushed Furuta's words out of my head and reached from behind to wash her hair.

As I washed her hair slowly, I felt a pang of sadness. Yuzuki's back looked like a fragile white shadow, and I remembered her as a child.

The sad figure from the harsh lessons I'd spied through the double window.

I thought that Yuzuki might be returning to her younger self's heart. In the face of death, perhaps she was returning to that time and trying to reclaim the love she was never given then.

Thinking that, Yuzuki's head and body suddenly seemed small and lonely.

With that melancholy feeling, I left the bathroom. Then I heard her call again: ``Yakumo-kun~.''

I opened the closed door wondering what it was, and was surprised.

Yuzuki had lifted one white leg out of the bathtub.

``A service for washing my hair.''

And she winked sexily. It was the first time I'd seen a Japanese person wink so well. Confused, I said:

``...Ah... yeah... thanks.''

Yuzuki's face turned bright red.

``---What kind of reaction is that! Idiot! Pervert! Get out!''

And despite having no fingers, she splashed water amazingly well, hitting me squarely in the face.

It was unreasonable. I still don't know what the right answer would have been.


\subsection*{10}

When I asked if she would see her parents, Yuzuki got very angry.

``I won't go. I've cut ties with them.''

``...Maybe just once? They must be worried...''

``No. I decided I won't see them again.''

``I understand why you resent them, but Ranako-san was strict out of love---''

``What do you know?'' Yuzuki cut me off sharply. ``---Besides, do you have the right to interfere in my family, Yakumo-kun?''

...No, I didn't. Not at all. Maybe living together with Yuzuki had made me feel like we were married. We weren't even lovers.

Yuzuki still seemed angry and said:

``Yakumo-kun, are you poisoned by dramaturgy?''

By ``dramaturgy,'' I think she meant the rules of storytelling.

``Don't you have that novelist's patterned thinking? Like parents who fight at the beginning of the story reconcile at the end for a happy ending?''

I hadn't thought of that. Maybe it was true. Maybe I'd unconsciously absorbed that way of thinking. It wasn't impossible. Yuzuki said:

``Stop it. People's feelings aren't that simple. People who know-it-all say things like 'you have to value your parents' or 'you have to forgive family' as if it's a template, but that's just the opinion of people whose parents were decent. It's the opinion of people from privileged environments who lack imagination.---Thanks to them, I did become a pianist. But for everything else, I hate them. I don't think I'll ever forgive them. It's not that simple. So don't try to process me with such a crude framework, like Miller's films. Don't turn me into a story so easily, with such insensitive hands.''

I couldn't say anything more. I'd thought of myself as sensitive, but Yuzuki had so vividly exposed my insensitivity. The person who had cried at the bullet holes in Warsaw's Old Town was the same person who had tried to intrude so carelessly into Yuzuki's delicate family issues---while that was strangely absurd, it was also entirely natural. Human imagination is, after all, limited. The same words that comforted someone yesterday can hurt someone today. Such things are common. That's why we must imagine: that we lack imagination, that no one has perfect imagination---

``Turning into a story, huh...''

Those words struck my chest strangely. I remembered Yuzuki saying she didn't want to be ``consumed'' when Miller came. At that time, she was also refusing to be ``turned into a story.''

What does ``turning into a story'' mean? And what does ``a story'' mean?

I thought about it. And why, on the other hand, did Yuzuki want to become a story in my novel?


\subsection*{11}

In early March, a package arrived from Poland.

It was a long, thin cardboard box, like one that might contain a violin.

Inside the tightly packed cushioning, there was a smaller box. When I opened it, a beautiful silver arm lay on a black sponge base.

``What is this---?''

Yuzuki opened and read the letter that came with it. Meanwhile, I picked up the silver arm and looked at it. The arm had delicate craftsmanship and appeared like a work of art at first glance... The intricate internal structure was intentionally exposed.

``It looks like a prosthetic arm. A mechanical prosthetic hand.''

Yuzuki said, handing me the letter. It was written in Japanese:

``I have asked a Japanese friend to write this for me.

My name is Emil Kamiński, a Pole.

I apologize for suddenly sending you this package. I really wanted to meet you in person, but due to my busy schedule, I must extend my greetings this way. Please forgive me.

(middle omitted)

---I am currently involved in the development of mechanical prosthetic hands at a company in Poland called Copernicus Technologies. What I have sent you is a new model under development called 'AGATERAM.' It has not yet been announced, and there is no information about it online.

The innovative mechanism, in addition to the conventional method using surface electromyography, uses AI deep learning on ultrasound data and repeated fine adjustments to achieve a custom-made fit...

(middle omitted)

Thus, we have created a device that is as light as a feather, requires far less training time than existing prosthetics, and allows free movement. It is also fully waterproof and suitable for all situations...

(middle omitted)

My request is actually about my daughter. Her name is Miaha. She is six years old.

Miaha was born without arms. She has a congenital absence of both forearms.

My daughter and I watched your performance at the 17th International Chopin Piano Competition at the Warsaw Philharmonic. She was deeply moved, even shedding tears, and admired you. She said she wanted to play the piano too.

But Miaha has no arms, and I knew it was impossible. Since then, she has been depressed, barely eating, and only watching your performance videos. It's as if she has finally realized that she has no arms...

As a developer of mechanical prosthetics, I boasted that I would one day make her arms, but she wouldn't believe me.

Miaha is currently in zero grade (I will add a note as the proxy writer: unlike in Japan, Poland does not have kindergarten; before entering first grade at age seven in September, one must complete zero grade). She also suffers from the difficulty of being in a group for the first time with a disability, often missing school and becoming withdrawn.

I told her, ``There are many hard things in life, Miaha, and you have to overcome them too.'' But she answered, ``Dad has arms, so he doesn't understand,'' and I'm ashamed to say I was at a loss.

(middle omitted)

I know this is an impertinent request, but would it be possible for Igarashi Yuzuki-sama to wear AGATERAM and play the piano? Even a little would be enough. Seeing that, I believe my daughter will have hope for the future and find the courage to face this harsh reality---

(remaining omitted).''

The letter was verbose. It was clearly written with a lot of hesitation, but also with passion.

``Yuzuki, what are you going to do---?''

``What can I do---right now it doesn't fit, so there's nothing...''

AGATERAM was made assuming a defect up to the middle of the forearm. Yuzuki still had her hands up to the back of her hand, so it would probably be a few months before she could wear it. She might have to spend what little time she had left practicing with the prosthetic...

``When it fits...?''

``I want to cooperate as much as I can.'' She answered immediately. ``But I don't know how well this arm will work...''

We researched mechanical prosthetic hands. The mainstream type works by picking up myoelectric signals from the skin surface and converting them into electrical signals for movement. They haven't become widespread in Japan. This is because most mechanical prosthetics distributed in Japan are German-made and cost as much as 1.5 million yen. A government subsidy is available with a doctor's certificate proving proficiency, but there are only about thirty training facilities nationwide where one can obtain that certification, and only three where children can train. Moreover, training takes about two to three years, requiring a training prosthetic, which is not subsidized and must be paid in full. It's no wonder they haven't spread.

In comparison, AGATERAM, produced using 3D printing technology, can be made at a much lower cost, requires almost no training, and allows natural control. That would be truly revolutionary. It could surely become a hope for many people.

We watched several videos of mechanical prosthetics online. Just recreating complex hand movements mechanically is amazing. I couldn't help but respect the engineers. However, none of the prosthetics seemed to have reached the stage where they could play the piano.

Of course, there were no videos of AGATERAM. It was a company secret. So in the end, we had no idea how this beautiful arm would move.


\subsection*{12}

March 11---

Yuzuki offered a silent prayer. Her figure, slightly bowed and eyes closed, was so beautiful it surprised me.


\subsection*{13}

April---Yuzuki completely lost everything below her wrists.

She began doing many things with her feet. She was quite dexterous, using her toes to operate her smartphone smoothly, which surprised me. Yet, I never once saw her make an indelicate movement. When she was about to make one, she must have skillfully avoided my gaze.

We often went for walks to that flower field. The wildflower field that Yuzuki had shown me in third grade, where I used to pick flowers for my mother. Passing through it, Yuzuki's family home was just nearby. We lived within sight of them.

Our living area was surrounded by a field of dandelions. When a strong wind blew, it became a vivid yellow sea. Golden waves rose, gently dissolved in the distant flower field, brushed the tops of the trees, and danced into the sky. It was like the playful whim of a wind tired of its endless journey.

Yuzuki sat down in the dandelion field and took off her shoes and socks. Her bare white feet on the vivid yellow flowers. And the crystallized stumps, whiter than her skin... Yuzuki was beginning to lose her toes. She might lose the ability to walk at any time... Perhaps that fear was why she started going on walks more often.

Yuzuki touched the stumps of her toes with her fingertips and said:

``It feels a little tingly.''

``Tingly? Like phantom pain?''

``Probably. Will it get worse? What was your mother like?''

``...She also suffered from phantom pain.''

When a gentle breeze blew, the dandelion flowers softly brushed her bare feet. It looked like a mother gently rubbing a child's aching belly. I wondered if that touch could heal Yuzuki's phantom pain.

``...Yuzuki, where have you been going lately?''

Yuzuki had started going out alone these days, sneakily, with a guilty look.

Gazing at the swaying dandelion flowers, she said quietly:

``...Before I can't move at all, I'm thinking of entering a hospice.''

Unlike hospitals focused on treatment, a hospice is a facility for relieving the suffering of patients nearing death. I felt a jolt, as if I could hear the approaching footsteps of Yuzuki's death.

``I don't want you to take care of me, so I'd like to go early.''

``Don't want me to take care of you?''

``Ah, it's not that I don't like you, Yakumo-kun...''

Yuzuki didn't say any more.


\subsection*{14}

May---Yuzuki had lost most of her toes.

She walked unsteadily, her steps uncertain, as if unable to balance properly. When she stumbled and was helped, she would smile with a little embarrassment, as if to cover it up.

Still, she kept walking. I think she must have treasured the feeling of her own feet pressing into the earth. Each calm step was a deeply spun farewell.

Before I knew it, it looked as if snow had fallen. The vivid yellow dandelion field had turned into a pure white snowfield. When the wind blew, the snowfield rustled and rippled---

All the dandelion flowers had turned into fluffy seeds at once.

``Wow---''

Yuzuki's eyes sparkled and she let out a breath that seemed almost white.

We walked, leaving footprints on the smooth, pristine snow. With each step, the fluff drifted softly into the air. When we sat down quietly, it rustled. Each round fluff, packed with seeds, made a sound like a tiny sacred bell, full of some kind of premonition. Afraid that speaking loudly would shatter the peaceful scene, we leaned close to each other's ears and whispered. Yuzuki's breath tickled.

``Hey, since we're here, let's tell each other secrets.''

Her voice was sweet.

``Okay.''

``You first, Yakumo-kun.''

``I secretly think this---'' I paused. ``'The Law of Perpetual Cultivation' has a really nice ring to it.''

Yuzuki laughed, but frowned.

``What's that? Yakumo-kun, you're the worst.---Okay, my turn.''

Yuzuki brought her lips to my ear. The tip of her nose coldly touched my ear. I heard a faint breath.

---Suddenly, the wind blew. A rough wind. I closed my eyes as my hair was blown about.

The wind softened. I opened my eyes, and a dreamlike scene spread before me.

The dandelion fluff that had covered the field like snow was taking flight all at once. Silently, violently, dancing in the wind, crossing the field, brushing the tops of the trees, rising into the May blue sky---

It was so white it filled my entire vision, a breathtaking beauty.

``It's scary---'' Yuzuki said. Her voice trembled slightly. ``My grandmother, who has passed away, once told me that if dandelion fluff gets in your ear, you won't be able to hear anymore. Yakumo-kun, cover my ears for me---''

She seemed genuinely frightened. It was strange that she could be so brave in the face of approaching death, yet so afraid of superstition, but it also felt very human.

I covered Yuzuki's ears with both hands.---Her large almond-shaped eyes blinked. Her breath trembled slightly. I felt like I was protecting a fragile life within my hands.

Suddenly, Yuzuki seemed unbearably endearing, lovable.

We looked at each other.

From there, everything felt natural, as if under a spell.

We shared our first kiss.


\subsection*{15}

For a while after our first kiss, Yuzuki would turn bright red just from meeting my eyes, and scurry off to hide somewhere. I was just as red-faced, and the two of us seemed to be endlessly playing a game of hide-and-seek without a seeker.

---One night, I was sitting on the living room sofa when Yuzuki suddenly sat down beside me. I was surprised by her sudden move. It felt like the day of reckoning had finally arrived.

We sat motionless, side by side, staring at a TV we didn't want to watch. A cheetah mother and child were on the savannah. A voice narrated: ``Oh, it's dangerous! An angry African buffalo!'' and ``---Whew. Close call, they're safe.'' If the camera had been on us, the narration might have been: ``Look! An awkward couple who kissed before dating!''

``Yakumo-kun---'' Yuzuki, face bright red, stared at the TV and asked: ``When did you start liking me?''

I couldn't help but startle. It was a fastball that could even strike out Shimizu, and she used polite language.---\textit{Bang!} A gunshot came from the TV.

``In the past, many African elephants were poached for their ivory---''

``...Did I ever say I liked you?''

``...Hmph. Even though you don't like me, you did that...?''

``The impala is being driven to the water's edge, with no escape---''

I gave in and said it. I think my face was bright red.

``...I liked you.''

``...F-from when?''

``...From the first time we met... I've always liked you...''

``...I-I did too...'' Yuzuki swallowed. ``...From the first time we met... I've always liked you...''

Being told she liked me, I was happy and embarrassed, not knowing what to do. On TV, the cheetah mother and child were making a bloody meal of the impala.

``...Um, Yuzuki-san, what should we do...?''

``...How about we try holding hands, Yakumo-san...?''

``...Um... Yuzuki-san, you don't have hands...''

``...Oh, right...''

Was it okay to laugh? Yuzuki said:

``...Then, how about... a kiss instead...?''

It was a mystery why she chose such formal words, but I had no time to worry about it. We faced each other. Yuzuki blushed, pursed her lips shyly, and blinked repeatedly. Then she softly closed her eyes.

Yuzuki's face was beautiful. I noticed again how long her eyelashes were.

My heart pounded loudly. Slowly, I moved my face closer and closed my eyes.

---A soft sensation.

Like a snowflake melting silently on my lips.

We pulled our faces away twice as fast as we'd brought them together. Yuzuki's face was as red as a cherry. She looked like she might even cry from embarrassment. And then:

``...Thank you very much. I look forward to working with you in the future.''

She said something strange. I was probably bright red too.

``Ah, yes, likewise. I look forward to working with you too.''

``...Good night, then.''

``Good night.''

Yuzuki left quickly. I was confused by the strange exchange, but listened to my still-pounding heart.

The eldest cheetah, now grown, gazed at the evening savannah with a dignified expression.


\subsection*{16}

That's how we clumsily became lovers. Yuzuki gradually revealed an unexpected side of herself. She was a kiss monster.

``Hey, Yakumo-kun, let's kiss.''

That became Yuzuki's catchphrase. Even after just a quick peck, her cheeks would turn pink, she'd giggle with satisfaction, and then go off somewhere. Like a drive-by kisser.

It seemed Yuzuki said ``kiss'' because saying ``kiss'' was too embarrassing. The order of embarrassment seemed to be: kiss, French kiss, snog. A mysterious ecology.

So when I said ``Let's kiss,'' she'd blush and shake her head, but when I said ``Let's kiss,'' she'd easily agree. Feeling a bit mischievous, I'd focus on \"French kiss,\" and she'd gradually get annoyed, finally getting really angry.

If Yuzuki came over and skillfully tucked her hair behind her ear with her armless stumps, she usually wanted a kiss.---So I'd think, \"She's probably about to ask,\" and my heart would pound, which was bad for it. I never got used to kissing. I liked Yuzuki too much; strangely, it made it harder to kiss her. So for a while, whenever Yuzuki approached, I'd flinch like a nervous herbivore. A sad ecology.


\subsection*{17}

While our sweet life continued, Yuzuki's salt disease steadily progressed.

By the end of May, her arms had turned to salt and crumbled away up to about the middle of her forearms. The void left by her lost arms brought me sharp, terrible pain. The progression rate of the disease varies. Thinking about how much longer Yuzuki could live made my heart feel tight.

``Hey, should we try AGATERAM soon?''

I'd completely forgotten. From the back of the closet, I took out that now-nostalgic box. Opening it, the two silver arms sat there, as beautiful as ever.

I took Yuzuki's soft arm and attached the hard one---

``It fits perfectly!''

Yuzuki raised her arms and looked at them curiously. They fit perfectly because they were custom-made for her, probably with measurements taken from performance footage.

Using the smartphone app that came with it, I activated AGATERAM.

Suddenly, a ``Pon!'' sound came from AGATERAM, and Yuzuki flinched.

``---Wow! It's already moving! It's moving!''

The silver fingers were wriggling. The smoothness and silence of the movement were astonishing. It was like watching a spider move. I performed a few more initial setup steps through the app, fine-tuning it specifically for Yuzuki. After completing all the sequences, the hand could move quite freely. Yuzuki stared at the opening and closing fingers with an expression like she was looking at an unparalleled jewel.

``This is amazing... It's so beautiful...---Here!''

Yuzuki suddenly pinched my nose. I laughed and said:

``Ow, ow, ow...''

``Ahaha, break off~''

The reason my eyes got a little wet was because I was happy that Yuzuki had working hands again.

``With this, maybe I can play the piano---!''

We hurried to a music store. On the way, we bought black leather gloves that could hide AGATERAM. It didn't seem right to walk around showing off a company secret.

We stood before the piano on display. Yuzuki positioned her hands over the keyboard.

And she took a deep breath---

The piano she had once lost. If she couldn't regain it, she would be sad. Even if she did regain it, knowing she would lose it again would be sad. Either way, sadness would follow.

But Yuzuki began to play, slowly---

A beautiful sound rang out. I shuddered involuntarily.

It was the continuation of the piece she hadn't been able to finish at the Chopin International Piano Competition.

Of course, the sound wasn't as beautiful as then, and the tempo was much slower. But the sound AGATERAM produced certainly filled something that had been lost then.

Yuzuki cried as she played.---When she finished, she buried her face in my chest and cried, muffling her sobs. I stroked her back gently until she stopped.


\subsection*{18}

In early July---we set foot on Polish soil again, in Warsaw. Yuzuki had wanted to go, so we arranged to meet Emil and his daughter there.

We toured places we hadn't been able to visit last time: the Chopin statue in Łazienki Park, Chopin's birthplace in Żelazowa Wola, the Pleyel piano in the Chopin Museum... We even extended our trip to the UNESCO-listed Wieliczka Salt Mine in Kraków.\footnote{The Wieliczka Salt Mine is a historic salt mine near Kraków, Poland, dating back to the 13th century. It is a UNESCO World Heritage Site.}

---And then, the day to perform for Miaha finally came.

In the morning light streaming through the hotel window, Yuzuki adjusted AGATERAM. The salt disease progressed by the hour, so she stuffed padding to make it fit snugly. The tension was like a gunslinger meticulously maintaining his weapon before a duel.

For this day, Yuzuki had spent a considerable amount of her remaining life. She bought an electric piano, had it brought to the apartment, and practiced about five hours a day.

``Half-hearted practice wouldn't give Miaha-chan hope---''

To avoid disturbing the neighbors, Yuzuki wore headphones, so I couldn't hear the sound. I tapped away at the keyboard in another room. While doing so, I thought that humans are essentially solitary beings. Like two boats traveling side by side on a canal. Even if they seem like one boat for a time, they are fated to drift apart on the vast sea...

---July in Warsaw was cool.

That day, the high was around 20 degrees, very pleasant. The sky was clear and blue, as if the piano's sound could rise lightly into it. The streets of Warsaw seemed to be dozing, a calm and peaceful day. We rented a room at the Chopin Music Academy. It was originally a day off, but the lack of people might have been the academy's consideration.

Mr. Emil was a slender gentleman with a long nose and kind, drooping eyes. His thin silver-rimmed glasses made him look every inch the creator of AGATERAM. His brown suit had a somewhat faded hue, and his tall stature somehow reminded me of ``the big old clock.''

``Nice to meet you. Thank you for today. I am grateful.''

Emil greeted us in Japanese. He must have practiced hard. Then, with a slightly hunched back, he shook our hands. He was already smiling, as if about to cry with emotion.

From behind his long legs, a small girl peeked out. She had been hiding out of shyness.

She was as cute as a doll. Wavy blonde hair, sky-blue eyes, a wide round forehead. She must have dressed up for the occasion, wearing a blue dress with white lace and a ribbon. Her arms ended at the elbows.---This was Miaha.

Yuzuki spoke to her in Polish. Miaha beamed, swayed shyly, and exchanged a few words. It was a heartwarming sight.

Emil set up a video camera on a tripod. The lens pointed at the piano. Clear light from the window fell diagonally across the keyboard. Since they would send us the footage later, I didn't need to bring anything. Yuzuki approached the grand piano, looking a little shy, and bowed. She wore a pure white satin dress. The simple beauty of AGATERAM stood out.

We applauded. Miaha also clapped her short arms happily.

Yuzuki sat down and saw the Anpanman plush doll on the piano. She had been so anxious about today's performance that she had brought it all the way here.

Then she lowered her eyelids and took a slow, deep breath.

At that moment, I gasped. Just like in her final performance at the Chopin International Competition, Yuzuki looked as if she had become part of the piano, a beautiful instrument.

A powerful forte, imbued with her soul, launched into the silence---

AGATERAM began to play. Spinning notes smoothly, as if they were real arms.

The black-key accompaniment, like a swaying boat.

Chopin's \textit{Barcarolle}---

In a rush, childhood memories came back to me. Memories from the day after I met Yuzuki---

The memory of that beautiful girl, who unabashedly compared her own performance to Maurizio Pollini's, saying things like ``I want more profundity and maturity'' or ``I want to fall in love quickly,'' so far beyond her years. The memory of a girl obsessed with the piano---

Those memories blossomed vividly, as if in an instant.

I was spellbound, captivated by the performance.

What a clear, pure tone it was---a transparent, angelic sound, as if a heavenly instrument were playing itself, impossible to believe it was played by a human.

I could see a gondola carrying lovers gliding through the canals of Venice.

The sea holding darkness, the sky connected to heaven, appeared in vivid illusions.

AGATERAM and the notes themselves were dazzlingly beautiful, like light shimmering on water.

Romantic, and somehow deeply nostalgic. So nostalgic it brought tears.

And somehow, sad. So sad, and so dear.

It was a poignant sound, cherishing and regretting the days that remained.

The boat did not stop, swaying and moving along with the flow of time---

Before I knew it, the streets of Venice had become the streets of Warsaw.

The destroyed Old Town---

The sound of prayer, polished by time, quietly showed the streets of the afterlife through the mist.

Gentle rain of prayer, funeral bouquets, falling upon the wounded city.

So that one day this city might heal.

So that it might quietly find salvation.

I was shedding tears without realizing it.

Yuzuki sang with the piano, so joyfully, so lovingly.

A pure hymn for the fragile human soul---

Emil covered his mouth with both hands, crying as if about to collapse. Tears formed small seas on his glasses, dripping down.

Miaha's eyes sparkled. They were eyes that could gaze straight at something beautiful and pure, eyes only a child possesses, clear as the blue sky.


\subsection*{19}

When the performance ended, we gave unreserved applause.

Yuzuki took the Anpanman doll and walked over to Miaha.

Miaha, with sparkling eyes, said:

``That was a beautiful performance!''

I somehow understood what she said. Though every word was in a foreign language, it seeped into my heart and made sense. Yuzuki smiled gently.

``Thank you. These hands are beautiful too, aren't they?''

``Yes, very beautiful!''

``Miaha's father made them for me.''

Miaha looked up at Emil, still crying, then turned back.

``Really?''

``Really. Your father will make hands for you too, Miaha.''

Miaha looked up again. Emil wiped his tears.

``Yes, leave it to me!''

He puffed out his chest. Miaha smiled, her face blooming like a flower.

``Will Miaha be able to play the piano like the big sister?''

Yuzuki smiled like the sun.

``You will! If it ever gets hard, look at this doll and remember me. It will give you courage.''

And she handed over the Anpanman plush. Miaha, with a joyful, loving look, hugged it tightly with her short arms, as if holding her own heart.

``Thank you---!''

Anpanman's red cape looked like the color of life itself.


\subsection*{20}

We headed to Chopin Airport to catch the evening flight.

The setting sun, streaming through the full wall of windows, dyed the airport interior red.

People passing by looked like silhouettes.

It looked like the final scene of a movie, I thought. A movie I'd seen in an empty theater long ago, whose title and plot I could no longer remember. The empty theater was stained with a lonely sunset color. I had been crying for some reason, and I walked home crying too. A memory I could no longer recall...

Thinking that somewhere within this red light was little Miaha, with her sky-blue eyes, felt strange.

The \textit{Tristesse} Étude was quietly playing---

I recognized the person playing. It was the same burly, bearded man I'd seen playing the Waltz No. 13 at Chopin Airport before.

I felt a strange familiarity toward him. But we were strangers, and would probably never meet again. Thinking that made me feel a little sad, and fond.

I looked back just a little after passing by him.

He noticed me and smiled. I smiled back.

On the plane, after buckling our seatbelts and settling in, Yuzuki said:

``Hey, Yakumo-kun, how was my performance?''

``It was wonderful. Beyond words.''

Yuzuki smiled a little mischievously.

``I really did fall in love three times, you know. All because of someone.''

``Huh? Me---?''

``Place your hand on your heart and think about it.''

I placed my hand on my heart and thought. But I couldn't figure it out at the time.

Yuzuki made a cute, frowning face---then smiled softly, looking a little sad.

``When everything is over and I think about it, it's all been like a cheesy movie. Even more than Miller's films...''

``Now that you mention it, maybe it is.''

``But you know, I was honestly happy. That straightforwardness. Miaha's directness. It made me so happy and blessed. It felt like I was saved...''

The plane took off. Looking down at the shrinking city outside the window, Yuzuki murmured:

``Goodbye, Warsaw.''

A tear traced down her cheek.

It wasn't a blue tear of sadness, like dissolved paint.

It was a transparent tear, as if the colors of the sunset were passing straight through it.


\fullpageimage[]{12.jpg}
% ==============================
%        Chapter 8
% ==============================
\clearpage
\thispagestyle{plain}
\phantomsection
\begin{center}
    \large\bfseries Chapter 8
\end{center}
\addcontentsline{toc}{chapter}{Chapter 8}
\label{sec:ch8}
\vspace{1cm}

\subsection*{1}

Perhaps after performing the \textit{Barcarolle}, something had changed within Yuzuki.

She became lighter than before, like a white bird soaring across the blue sky. It seemed contradictory that Yuzuki, who could no longer move without a wheelchair after losing her feet, had become lighter. I must not have understood the true nature of freedom.

True freedom is not about going to see a distant blue sky, but about keeping the most beautiful blue sky within your heart. Miaha had become Yuzuki's blue sky.

That left me feeling like the one left behind. So I asked:

``Is there anything you wish for...?''

It was a small chain to keep Yuzuki tethered to this world.

``Hmm, let me think about it...''

Yuzuki thought for about three days. Then, just when I'd started to forget, she suddenly said:

``Yakumo-kun, I want to wear a wedding dress.''

``Huh?---Oh, a wish. Where can we rent a wedding dress?''

``...That's not what I meant.''

``Huh?''

``...And you call yourself a writer? I'll give you thirty minutes, so think carefully!''

Yuzuki huffed and rolled her wheelchair away to the next room.

I gaped like an idiot and actually thought for ten whole minutes.

Then I rushed out of the room and went to a flower shop.

To buy a bouquet for a proposal---


\subsection*{2}

Yuzuki insisted on greeting my father, so I had to contact him.

After agonizing for about three hours, I finally sent this message:

``It's been a while. I'm getting married.''

About three hours later, a reply came:

``Congratulations. I'm surprised.''

The lack of information was so glaring it actually gave me a headache. We continued the exchange in bits and pieces.

``My future wife wants to greet you.''

``I see. When?''

``As soon as possible. She's in a wheelchair, so I'd like you to come here.''

``I'll come tomorrow. 3 PM.''

He was surprisingly prompt. Maybe novelists are flexible in these situations.

Worried about this belated three-way meeting, I couldn't sleep at all that night.

The next day---at the promised 3 PM, my father's beat-up Mercedes stopped in front of the apartment.

The doorbell rang. I opened the door, and a shadow stood there.

---It was my father, whom I hadn't seen in five and a half years.

He seemed much smaller than in my memories. I had grown taller. He was still wearing black, as slim as a shadow. He hadn't aged so much as to seem old, and with his eyepatch, he'd gained a certain ruggedness, like Date Masamune,\footnote{Date Masamune (伊達政宗, 1567–1636) was a famous feudal lord of the Sengoku period, known for his distinctive crescent-moon helmet crest and for being one-eyed.} which annoyed me for some reason.

``Thank you for coming. I'm Igarashi Yuzuki---''

Yuzuki smiled and held out her hands. My father looked alternately at Yuzuki and AGATERAM. I couldn't read his expression at all. He smiled and shook her hand.

We sat at the living room table. It was a full three-way meeting. As I sweated nervously, unsure of what to do, my father spoke:

``I never imagined that the person Yakumo would marry would be the pianist Igarashi Yuzuki. I knew you went to the same school, but... I was so surprised I nearly fell over---'' That was surprised? I thought. ``Your performance at the Chopin International Competition was truly wonderful. I was moved.''

``Thank you, Father-in-law.''

Yuzuki smiled. \emph{Father-in-law}, huh---I thought.

My father scratched his cheek, looking a little embarrassed.

``You can call me something more casual, if you like.''

``Then, Father.''

My father nodded---but he still looked a little embarrassed. After that, Yuzuki and my father chatted easily. They talked at length about various music, composers, pianists, and even the folk music of different countries. I thought I had a fairly broad knowledge, but I couldn't follow them at all.

``Even among Eskimos, there are groups that can sing in harmony and those that can't, apparently---''

``Yes, so I've heard. Between the whale-eating groups and the caribou-eating groups---''

It was terrifying how casually they spoke of it. My father occasionally cracked clever jokes that made Yuzuki laugh. Had he always been this talkative? I wondered. It must have been because of the stark contrast with his terse, curt emails. When I thought back to my childhood, my father had always chattered on about all sorts of things.

In a short time, Yuzuki and my father exchanged a vast amount of information and seemed to have hit it off completely. Then, as he was leaving at 6 PM, he said:

``If you're going to have a ceremony, I can provide financial support. I can cover the whole amount.''

``You have that much money?'' I interjected.

``I have plenty. I earn it, and I don't spend it.''

``Plenty, huh...''

It was quite a shocking fact, but upon reflection, it made perfect sense.

``Thank you,'' Yuzuki said.

``Thank you, Yuzuki-san. For marrying Yakumo. I'm truly happy.''

The moment my father said that, Yuzuki lowered her face and covered her mouth with her left hand.---Her eyes glistened with tears. It was only at this moment that I realized Yuzuki had been carrying a burden of guilt.

She must have thought that someone about to die getting married was ridiculous. And then, receiving such kind words, she couldn't help but cry.

In that moment---I felt I could forgive my father. Even though he was lacking in some human qualities, had made my mother unhappy, and had left me on my own. For some reason, I felt I could forgive him---

After my father left, Yuzuki said:

``You spoke ill of him, but I thought he was a wonderful father.''

I nodded quietly and said:

``He has both good and bad sides.''

Yuzuki nodded. After a pause, I said:

``I think we should greet your parents next. What do you think?''

In the past, she would have refused immediately, but Yuzuki's expression stiffened as she hesitated.

After playing the \textit{Barcarolle}, something had indeed changed within Yuzuki.


\subsection*{3}

``Hey, let's not go after all---?'' Whenever Yuzuki said that, I stopped properly.

``Are you sure you want to stop?''

Then Yuzuki would fall silent on her wheelchair, like a doll, and I would start walking again. Our apartment and Yuzuki's family home were less than a twenty-minute walk apart. Yet we had been wandering the neighborhood for nearly two hours. It was early August, morning. The sunlight had the transparent brilliance typical of summer. I had a feeling it would get hot today. Above the New Town streets where beautiful summer flowers bloomed, white cumulonimbus clouds loomed.

``...Ya-Yakumo-kun, let's take a detour. Let's go this way...''

I obediently changed course, fully understanding how difficult it must be for her.

``Oh, I wonder if Melody is still around?''

I said. Melody was the golden retriever who loved Yuzuki so much that whenever they met, she'd wet herself with joy.

``Now that you mention it. I used to visit her often up until middle school, but how is she now?''

We headed to the house where Melody was kept.

``---Ah, there she is, there she is,'' I said, seeing her from a distance. But Yuzuki shook her head.

``No, she looks very similar, but that's not Melody.''

Indeed, Yuzuki was right. The dog was too young.

``Rhythm''---a nameplate on the doghouse read.

``Melody... died...?''

Yuzuki turned to look up at me with a very sad expression.

``...Maybe. But Rhythm could be Melody's child, couldn't she?''

Yuzuki's eyes widened.

``Maybe! Hey, Rhythm, come here.''

Rhythm joyfully wagged her tail so hard it seemed it might come off, and jumped on Yuzuki. Yuzuki, a little surprised, laughed and stroked Rhythm.

``Hmm, she does seem to resemble Melody---''

Just then, Rhythm suddenly let out a rhythmic stream of pee. ``---Whoa!'' Yuzuki's skirt was soaked. She looked at me with a stunned expression.

I couldn't help but burst out laughing.

Yuzuki also burst out laughing.

We laughed until our stomachs hurt. I said, still laughing:

``No doubt about it, she's Melody's child!''

``Ahahahaha, no doubt! Ahaha, it hurts, ahahahaha!''

Yuzuki's eyes glistened as she hugged Rhythm tightly.


\subsection*{4}

Seeing us suddenly appear, Ranako and Sosuke were stunned. Yuzuki was in a wheelchair, had a cool silver arm, and was covered in dog pee—so it was understandable.

Rather than being moved, they seemed frozen in place with sadness.

While Yuzuki showered, I sat in the living room drinking tea with Ranako and the others.

The inhabitants of fairy tales never age. The Grimm-themed furniture and decorations from my childhood remained unchanged. On the other hand, Ranako and Sosuke had visibly aged. Their faces had more wrinkles, and their sharpness had diminished, giving them a softer presence. When we met at the hospital last year, perhaps the emergency had kept me from noticing.

---No, it must have been the approach of Yuzuki's death that had aged them so rapidly. They exuded an air of being worn down by nights of worry, unable to see their dying daughter.

Both of them had downcast brows, looking lost and sorrowful.

``Yuzuki said she didn't want to see you, Ranako-san and Sosuke-san---''

I slowly recounted the events leading up to now. They listened with sad faces, nodding.

``Yakumo-kun, help me again!''

Hearing Yuzuki's voice, I went to the bathroom to help her get dressed.

``Whoa, the waist is loose. Mom, you got fat...''

Yuzuki said, wearing Ranako's clothes.

I put her in the wheelchair and pushed her to the living room. I stopped her in front of her parents, across the table, and sat down myself. Finally, the four-way meeting began.

``I'm marrying this person. That's all. ...Well then, goodbye.''

Yuzuki tried to end it quickly and wheel away, but I hurriedly stopped her.

``Wait, wait, Yuzuki---there's more you need to talk about.''

With her wheelchair held in place, she seemed to give up. She looked down and was silent for a while.

``I---'' Yuzuki said, still looking down. ``I still haven't forgiven you two.''

Her voice was cold. Ranako and Sosuke hung their heads like defendants.

What followed was a trial. Yuzuki, with all her feelings, accused her parents of their crimes.

``Mom---just because you ended up as a second-rate pianist, you tried to make me first-rate instead, forcing me to practice and even resorting to violence. I think that's truly despicable. Don't impose your own ego and inferiority complex on others. Don't make someone a sacrifice for your own desires. Because of you, I nearly came to hate the piano I loved. I hated you, Mom; you seemed so ugly. Thinking that half of me is you made me terrified, and I wanted to throw away half my blood...''

``Dad---because you couldn't even become a second-rate pianist, you had an inferiority complex toward Mom, and I was sad that you never stopped her violence against me. You always pretended not to see, and the sound of you running away up the stairs, \textit{tontonton}, I found unbearably sad. And I hated how you tried to compensate by giving me money like a fine. That's just self-satisfaction; it doesn't save anyone, and it felt like you were putting a price on my sadness—I really hated it. Thinking that half of me is you made me feel pathetic, and I wanted to throw away the other half of my blood too...''

By the middle, Yuzuki was already crying as she spoke.

Sosuke was the one who wept uncontrollably. Ranako bit her lip, her face pale.

Had Yuzuki been suffering that much?---I was once again surprised.

She hated the blood that flowed through her---that suffering seemed to be the exact opposite of what I had felt as a child. I had felt as if no blood flowed through me.

Half of me was salt, and the other half was shadow.

That felt more like sorrow than the pain of existence.

``Even though you're my parents---'' Yuzuki said, crying. ``Even though you're my father and mother, why did you make me hate you---!''

Tears dripped from between the fingers of AGATERAM onto the table.

Then, the expressionless face of Ranako, who had been enduring so desperately, finally crumbled.

``I'm sorry,'' she cried, even with a runny nose. ``I was wrong...''

Finally, she made the biggest puddle of tears on the table among the three of them.


\subsection*{5}

All preparations for the wedding proceeded at a breakneck pace.

The wedding company, sympathizing with Yuzuki's illness, said they could prepare everything in just three weeks and went to great lengths to help us. We sent invitations to classmates from our middle school days.

And on August 25, we held our wedding ceremony.

To the music of Chopin's Nocturne No. 2, Yuzuki entered---

Sosuke pushed her wheelchair. Yuzuki in her wedding dress was inexpressibly beautiful. She was beautifully made up, with a tiara and veil, earrings, and a lustrous pearl necklace. AGATERAM was polished to a gleaming shine, proudly reflecting the light. The wheelchair was decorated with white ribbons, a testament to the meticulous care.

We sang hymns, read from the Bible, prayed, and exchanged vows. Then, the ring exchange---the smaller one was mine, the larger one Yuzuki's. I took her silver arm and put the gold ring on it. By now, AGATERAM was no longer just a machine; it had become a precious part of Yuzuki.

I lifted her veil and kissed her.

---Then came the reception. Old friends showed up one after another.

``Yacchan, Yuzuki-san, congratulations on your wedding!''

Seeing that familiar face, I felt a little like crying.

``Thank you, Shimizu---!''

He was wearing an unfamiliar suit, with that familiar face perched on top. He had grown so much wider that his face looked relatively small. It was a funny build, like the toy ``Blackbeard's Peril.'' From behind his huge frame, little Kobayashi Koyomi popped out.

``Yuzuki, congratulations on your wedding!''

``Thank you,'' Yuzuki said, holding out her hands. Kobayashi shook them, eyes wide.

``Wow~, is this a prosthetic hand? Amazing, so beautiful, so high-tech~!''

Yuzuki wiggled her fingers.

``Hehe, right? It even shoots beams.''

``Really~! It shoots beams~!''

It didn't shoot beams.

Then Aida appeared from somewhere. He seemed to have stopped his flashy ways, now a clean-cut, handsome guy. But the moment he saw Yuzuki, his face reverted to that of a sniveling brat, and he squeezed my hand painfully, tears in his eyes.

``Congratulations on your marriage!''

He really did love Yuzuki... I felt a little sentimental, but then I remembered he'd been happily dating and enjoying life. I'd been sentimental for nothing.

``Yuzuki~! Congratulations on your wedding~!''

The girls surrounded Yuzuki.---Seeing that, I felt a little nervous. They were the ones who had bullied Yuzuki in the past. I watched anxiously, but it seemed I'd worried for nothing. Yuzuki and they were all chatting happily.

At some point, Yuzuki had become good friends even with those who had bullied her.

Yuzuki is amazing---I thought again.

A giant wedding cake arrived.

As Yuzuki and I cut the cake together, I thought, \emph{huh}.

We're being ``story-fied'' like crazy.

Yuzuki, who had so strongly rejected being ``story-fied'' or ``consumed,'' was now, with death approaching, speeding through her life, collecting all the plot points and trying to wrap things up.

Did Yuzuki herself realize this?

Knowing Yuzuki, she surely did. She was doing it knowingly.

Why was that---?

I thought about it, but in the end, I was too caught up in the fun of the lively reception to remember.

My eyes met my father's. He was, as always, quietly at the edge of the venue, like a shadow. He might have always been like that, by my side.

Thinking that, the feeling that welled up in my chest was not unlike love.


\subsection*{6}

That night---we sat side by side on the sofa at home, drinking cocoa together.

Some of the reception's alcohol still lingered, and I felt a little lightheaded.

``I guess we don't need a honeymoon---'' Yuzuki leaned on me. ``Our trip to Warsaw felt like our honeymoon.''

``I feel the same way. Though the order is reversed.''

``Order doesn't matter.''

Yuzuki said, then suddenly chuckled.

``We really are married. Hey, Yakumo-kun, do you want to be called 'darling'?''

``Huh? I've never thought about it...''

``Really? Dar-ling---''

``...Not bad. Kind of embarrassing though.''

``I felt a little embarrassed too.''

We laughed.

``Let's get to bed soon---''

We brushed our teeth in the bathroom. I helped her get ready for bed, pushing her wheelchair to her bedroom. Then I carried her like a princess and laid her on the bed. I covered her with the blanket.

``All done. Well then, good night---''

``Wait.''

She grabbed my sleeve as I was about to leave. She frowned.

``Huh---wait... really?''

``Huh---what?''

``What do you mean, 'what'---we're married, you know?''

Yuzuki's face was red.---I thought for a moment, then \emph{oh}, and felt my cheeks grow hot too.

Then, we slept in the same bed.

Yuzuki's white body was so beautiful, so fragile, as if it might shatter into salt crystals, and it scared me---

In the middle of the night---I woke from a sad dream to find Yuzuki crying.

She was facing away from me, her white back to me, sobbing quietly.

``...Yuzuki? Why are you crying---?''

Her shoulder twitched. But she kept crying.

``...I remembered... something from the past...''

I gently rubbed her lower back with my right hand, as if to soothe her. She took my left hand with her cold AGATERAM, held it like a pillow, and clutched it to her chest.

``Yakumo-kun, why didn't you come back then...?''

``Then...?''

``When we were in third grade, that time we said we'd run away together...''

I remembered. In a rush, the memory returned. The day I'd packed but never went. My heart pounded. Yuzuki's heart beat sadly.

Why hadn't I gone that time? I thought. There was no clear reason, I felt. It was just a vague unease that had overwhelmed me---

``Hey, why didn't you come...? I waited with Anpanman and the 1.5 million yen, for so long... It got dark, and the snow started to fall... I was cold, so scared, so lonely... And still, I waited alone, for so long...''

Yuzuki's body was trembling. Thinking she had been trembling like that back then, too, guilt crushed my heart.

``...I'm sorry,'' was all I could say.

``I wanted you to come so much... I felt like my heart was dying... I still dream about that day... And every time, I feel so lonely and scared, I can't bear it, and I wake up crying... Do you understand this feeling, Yakumo-kun...?''

I held Yuzuki's crying back in my arms.

I wished I could travel back in time and go help Yuzuki.


\subsection*{7}

One night, I had a strange dream.

In a pitch-black room with curtains drawn, I was watching TV. The footage seemed to be shot on a camcorder, shaking badly. On the screen was Shimizu, in his wedding suit, laughing \textit{wahaha}. Then Aida appeared, also in his wedding suit, with his thumbs in his ears, wiggling his fingers, blowing on a toy whistle. Next, Yuzuki's former manager, Kitajo Takashi, appeared, taking pictures with a crazed smile, the flash blinding.

---I don't remember what happened after that.

When I woke up, I felt oddly refreshed, and I had tears on my face.

I couldn't imagine how it could have ended that way from that beginning.

What a strange dream.

---Three days later, Yuzuki's arms quickly crumbled, and she lost AGATERAM.

It was the middle of September, when the air was growing chilly.

Yuzuki entered a hospice.


\subsection*{8}

The hospital room was full of flowers. Because I'd kept bringing them.

``You bought more? Honestly, what am I going to do with you...''

Yuzuki smiled, exasperated, forgiving my odd behavior.

The salt disease was progressing rapidly. Her arms had crumbled to the middle of the forearms, and her legs to about one-third of the thighs. The void left by her lost limbs brought back that special phantom pain. It hurt so much, it was so hard, that I bought bouquets for Yuzuki, just as I had picked flowers for my mother as a child. Yuzuki understood.

Yuzuki said she was scared to be alone at night, so I stayed by her side every night until she fell asleep. I wanted to hold her hand, but Yuzuki no longer had one. So, as if putting a child to sleep, I stroked her head or gently patted her belly with a steady rhythm.

When Yuzuki fell asleep, the hospital room became like the bottom of the deep sea. When I drew the curtains, moonlight streamed in, illuminating the flowers in their many vases like quiet flames.

I prayed the flowers would protect and warm Yuzuki, then went to my separate room for the night.


\subsection*{9}

In October, Yuzuki began to suffer from pain.

In addition to the phantom pain, the salinification of her internal organs had begun, causing excruciating pain.

By then, Yuzuki had significantly restricted visiting hours by her own will. During the times I wasn't there, Ranako was constantly at her bedside.

Once, I arrived at the hospital room a little earlier than visiting hours.

Yuzuki was screaming. A terrible scream.

``---It hurts! It hurts! It hurts! Mom, help me, it hurts!''

I froze at the door. I had never heard Yuzuki scream like that. I turned back and went to the room at the proper time later.

Yuzuki sat beautifully in the soft autumn sunlight, even wearing a little makeup. As if the scream I'd heard was just a bad dream.

---Ranako grew thinner every time I saw her.

Dark circles formed under her eyes, her skin was rough, her nasolabial folds deepened, and her gray hair became more noticeable.

As we drank coffee together in the hospice lounge, she said:

``Yuzuki doesn't want to show you anything ugly. So all the ugly parts are my responsibility. ...This must be Yuzuki's revenge.''

I touched the wedding ring on my finger and thought, perhaps that was true.

Yuzuki was trying to take revenge by making Ranako witness her suffering as she died.---But it was a revenge for forgiveness. Just as a court sentences a criminal and then grants forgiveness, Yuzuki was trying to forgive Ranako. By confronting her with pain and hatred, by wallowing in the ugliness, she was trying to extract pure forgiveness from the depths of that dark abyss.

So---

``That is her love. It means she wasn't cold enough to become indifferent to her family. Yuzuki always wanted to be loved by you---''

``...I see. How should a mother treat a daughter who is more clever and accomplished than herself...''

``It's simple---'' I answered immediately. ``Please touch her where it hurts.''

Ranako nodded, her eyes red, and went to Yuzuki, discarding her coffee cup.

In my mind, I pictured the double window.

Yuzuki playing the piano, and Ranako stroking her head---


\subsection*{10}

October 17---Chopin's death anniversary had come again.

Yuzuki's eyesight had deteriorated, her hearing too, and she seemed dreamy in the morning light, half like a spirit. The Yuzuki who had been so full of life now seemed to be slipping out through the cracks, and in the spaces left behind, it was as if spirits and flower sprites had entered, giving her an ethereal quality.

Yuzuki looked at me with hazy eyes and said:

``---I've been thinking lately about where I want to die.''

I felt as if my breath had stopped.

``Since dying is a very sad thing... I want to die in the place I love most.''

``...Yeah.''

``After thinking hard about it---I think I like the secret base the most. The time I spent with you inside that old, beat-up bus is so dear to me---''

Remembering those days of childhood, filled with so much misery and so much happiness, I couldn't help but tear up.

Yuzuki said:

``...Take me there. Today is my last day.''

\subsection*{11}

I carried Yuzuki on my back and left the hospice.

It was a calm, clear autumn day. We walked slowly along the avenue of ginkgo trees, now dyed yellow. The gentle wind carried the soothing scent of fallen leaves. Yuzuki's body was very light. She had lost most of her arms and legs, so it was only natural. The wedding ring on its necklace felt hard against my spine. Yuzuki was still soft and warm, and I couldn't believe she would die soon. When the wind suddenly stopped, a gentle scent of freesia wafted by. Until the very end, Yuzuki had tried to stay beautiful.

Her quiet breaths, like sleeping, fell on my neck. It was such a warm, sunny day that I kept thinking she might have simply fallen asleep.

With each step, there was a \textit{saku, saku} sound. It was the sound of salt rubbing together inside Yuzuki's body.

It sounded like walking through snow---the illusion of winter cold made me involuntarily shiver.

``Let's take a detour---''

I turned the corner again and again, as if reaching the abandoned bus would mean Yuzuki's death.

Yuzuki said nothing. She only gazed with eyes clear as glass beads at the apartment we'd lived in together, her family home, Rhythm the golden retriever, the elementary school, the park where we used to play, the scenery of our hometown---all the things we encountered on our detours---with loving eyes.

\emph{Thank you, goodbye.} She must have repeated those words in her heart many times over.


\subsection*{12}

The broken bus was still there, in the same place, after eleven years.

Its familiar, endearing face was buried among the withered grass.

It had been an old bus to begin with. Eleven years hadn't changed its appearance dramatically. Time passes very slowly for very old things. Like Machu Picchu, the Mona Lisa, or the Parthenon.

Stepping inside the bus felt like stepping into that slow time.

Warm sunlight streamed through the windows. Dust danced in the air, sparkling. The bench still had its cover and cushions. The comic books were still there, in their paper bags.

Distant memories came vividly alive. I could almost see my third-grade self and Yuzuki lying on the bench, playing happily.

We sat on the bench, our grown bodies making it feel smaller than in our memories.

Yuzuki had lost her strength and couldn't even sit up straight. She leaned on me, breathing quietly. I held her waist gently.

``...It's nostalgic.''

``It's nostalgic.''

We talked about memories, bit by bit. From the time we met, slowly, as if tracing our lives. When the memories caught up to the present, we went back to the beginning again. We were spinning through time. While we did so, before we knew it, a long, long time had passed, and we felt like we had become old. Both Yuzuki and I had become wrinkled old people. Having talked so earnestly all this time, we had become old people, bent at the same angle---

Of course, that was just an illusion. Time was passing moment by moment.

``My eyes... are going dark. Yakumo-kun... let me see your face... one last time...''

I faced Yuzuki. The sunlight through the window had already turned to dusk.

Her eyes, swaying as if in confusion, reflected the red light.

``Can you see...?''

``I can... I love... your face, Yakumo-kun...''

Yuzuki's right arm lifted. I knew she was trying to touch my cheek with her nonexistent hand.

``Ah---''

She let out a sad sound. I knew she had gone blind. Her eyes wandered, blinking repeatedly.

``I've... gone into the night... a little early...''

Yuzuki closed her eyes and rubbed her right eyelid against my nose.

``Good... I can... still picture your face... clearly...''

Then quietly, night fell. It became pitch dark, and faint moonlight streamed through the window. Yuzuki's white form seemed to glow faintly.

``It's strange...'' Yuzuki said, with a sigh-like voice. ``So much has happened... and now... everything... seems so... deeply nostalgic... As if... I've become a baby again... Do people... return to... the place they were born...?''

It was as if with every breath, life was leaving her body.

Yuzuki was going to die---Suddenly, I felt as if the entire night sky was pressing down on me.

I couldn't live without Yuzuki---

I finally realized it.

``Yakumo-kun... will you be okay... without me...?'' Yuzuki asked, truly worried. ``Can you... live alone...? Won't you get hungry... and cry...?''

Tears spilled down my cheeks. But I held back so as not to sob, so my body wouldn't shake. I wanted to let Yuzuki go in peace---With that thought, I hid my tears. Fortunately, Yuzuki could no longer see them.

``I'll be fine. Don't worry.''

Yuzuki smiled softly.

``...Good...''

I stroked Yuzuki's body.

``Are you okay? Does it hurt anywhere?''

``I'm fine... I'm fine...''

``Yakumo-kun...''

``What is it, Yuzuki...?''

``Yakumo-kun...? Yakumo-kun...?''

I realized with a start. Yuzuki had gone deaf.

Yuzuki began to cry.

``My ears... no... only my ears... only my ears... I don't want this...! I can't... hear... music anymore...!''

The image of Yuzuki, frightened by the dandelion fluff, flashed through my mind.

I finally understood that there are things in this world more painful than death.

I held Yuzuki's trembling body as I cried. Her trembling transmitted to me, and I felt I might freeze too.

``I'm scared... Yakumo-kun... I'm scared... I don't want to die... I don't want to die...!''

I rubbed her shaking shoulders desperately. Grains of salt fell, one by one.

Yuzuki kept crying.

``Yakumo-kun... why... didn't you come...! I waited... I waited... I waited so long... why... why...!''

I knew she was talking about the time we'd planned to run away. That night had remained deep in Yuzuki's heart. She had been so afraid of the loneliness of that night.

``I'm sorry... Yuzuki... I'm sorry...!''

I could see little Yuzuki, crying alone inside that bus.

I wanted to go back in time and save that lonely, white-shadowed girl.

I wanted to somehow enter that terrible night and save her.

I cried and prayed. God, if you exist, please help her. Please have mercy. Save Yuzuki from this most terrifying night---

Then, one thing occurred to me.

Yuzuki, please listen. Please, listen to me...

I tapped Yuzuki's shoulder rhythmically with my left hand.

I kept tapping, over and over...

``Ah---'' Yuzuki stopped crying and murmured. ``The \textit{Barcarolle}---!''

I had conveyed to Yuzuki, with my left hand, the accompaniment that evokes a boat swaying on the canals of Venice.

As best I could remember, I continued the same motion.

Nocturnes, études, polonaises, mazurkas, sonatas, ballades, scherzos...

I was surprised myself at how much I remembered. I had been watching Yuzuki play all this time. What had always saved me in my hardest times was Yuzuki's music. Just as light in darkness is immediately burned into the eyes, Yuzuki's performances were vividly burned into my heart.

``Yakumo-kun... I can hear... I can hear it... I can hear my music...! Not only that... all the music I've ever heard... I can hear it... all the music... that people made... and played... with love... The darkness now... it's so bright...!''

``Yuzuki...''

``The music came to save me...! The things I truly loved... in the end... they really came... to save me...''

``Yuzuki---'' I could no longer be brave. ``Don't leave me...''

``Yakumo-kun... bye-bye... I love you...''

``Yuzuki.''

``Don't... catch... a cold...''

And then, Yuzuki drew her last breath.

In the darkness, her heartbeat gone, I cried alone.

I heard the sound of Yuzuki's body crumbling into salt.

Frightened, I closed my eyes, covered my ears, and cried like a child.

\fullpageimage[]{13.jpg}
% ==============================
%        Chapter 9
% ==============================
\clearpage
\thispagestyle{plain}
\phantomsection
\begin{center}
    \large\bfseries Chapter 9
\end{center}
\addcontentsline{toc}{chapter}{Chapter 9}
\label{sec:ch9}
\vspace{1cm}

\subsection*{1}

Into the vase that had held my mother's salt, we now placed Yuzuki.

Ranako and Sosuke wept before it. Ranako's face was gaunt with exhaustion; she pressed a handkerchief to her mouth, tears streaming down her cheeks.

``Uuu... I'm sorry, Yuzuki... I'm sorry...''

She apologized over and over. Watching her, I felt nothing. My tears had long run dry.

``I have a message from Yuzuki.---She said she forgives you both. She said she loves you.''

When I relayed this, Ranako collapsed to the floor and wept.

---It seemed my father had arranged the funeral.

Neither I nor Ranako nor Sosuke could move, overwhelmed by grief.

My classmates and friends all seemed to have come, but I don't remember any of it. I was empty; some barely-functioning system was moving my body semi-automatically.

When I came to, I was inside the abandoned bus, holding the vase containing Yuzuki's salt.

It was the middle of the night. The winter moon cast a cold light inside the bus.

I had lost time. When I touched my chin, I found my beard had grown.

I sat there in a daze for a while. My heart felt covered with a thin gray cloud. I couldn't think, couldn't feel---a sad tranquility, like an old city slowly being buried under volcanic ash.---But the moment I thought of Yuzuki's death, the overcast sky turned pitch black, thunder roared, the sea groaned and swallowed the ash-covered city. The void left by losing Yuzuki assailed me with fierce pain. Unable to bear it, I stopped using my mind.

I fled the pain of living, seeking salvation in death's peace.

I sat there like a poorly-made doll, sometimes leaking tears like a leaking pipe. Glancing beside me, Yuzuki wasn't there---

In the moonlight, the salt crystals I'd missed glistened.

I picked up the crystals at the rate of about one per hour.

That became the rhythm of my existence.

At dawn I'd return to the apartment, eat something, sleep, and at night I'd snap awake. The darkness terrified me, made me restless. I'd hug Yuzuki's vase and escape to the old bus, finally able to sigh and zone out.

That was a finished world. The other side of Mount Adatara, where a giant bomb had fallen the day my mother died. The gray world where cold snow fell on the day of the earthquake.

What was I supposed to do there?

Everything was hollow and sad---how was I supposed to live?

...I understood nothing, just picked up salt, one grain per hour.

When I dropped a salt crystal into the vase, it made a faint, faint sound. Like a tiny bell ringing far, far away. I was a sandglass, a scrap of an instrument, just for making that sound.

Sometimes I'd fall asleep in the bus. I always dreamed of Yuzuki. Happy dreams. In the warm sunlight shining into the bus, we'd drink tea, read manga, spend peaceful, drowsy time together. When I said ``kiss,'' she'd blush and shake her head; when I said ``French kiss,'' she'd make a complicated face; when I said ``smooch,'' she'd happily do it. Such embarrassingly happy dreams.

When I woke up, I'd remember that Yuzuki was nowhere in the world, and feel lost. The sadness and loneliness of being left behind would make me cry.

Even so, I began sleeping in the bus just to have those happy dreams.

No matter what pain I'd have to endure afterward, I wanted to see Yuzuki again.


\subsection*{2}

I was woken by a rough sound.

Lifting my groggy head, I met the blank stare of a man.

He was in blue work clothes, a towel around his neck, with a beard like a thief's.

``Whoa---'' He cried out in surprise and turned to shout: ``There's a vagrant here!''

As I stared blankly, he stomped in with his work boots. ``Oh, you're young... Can't sleep here, it's private property! Come on, out, out!''

I was quickly chased out of the bus. Outside were several other workers. A rusty gate was open, and a large transport truck and a rough-terrain crane truck were parked on the gravel lot. The truck's door bore the company name ``OMOYA Construction.''

The crane truck is for lifting heavy loads, but I knew nothing of that; I just watched the men work from a distance, dazed.

The thief-bearded man came over, shooing me away.

``Hey, it's dangerous, move back, move back!''

As I stepped back, the crane let out a low groan. With a creaking, crushing sound, the bus was slightly lifted---it was only then that I snapped back to myself. Filled with an indescribable anxiety and panic, I asked the thief-bearded man, my voice trembling:

``---Wh-where are you taking the bus? What are you going to do with it!''

``Isn't it obvious! We're hauling it away!''

``Wh-where? What are you going to do with it?''

``We're handing it over to another contractor in the city. What they do with it, I don't know. Probably scrap it?''

Scrap it...? That bus, steeped in memories with Yuzuki, was going to be scrapped...?

Just imagining it made tears flow uncontrollably.

``No! Please, don't! Don't scrap it!''

The man, a little intimidated by my sudden outburst, frowned.

``I get how you feel, but you were trespassing!''

He thought I was a vagrant.

``It's not that... that's not it... Please...!''

How could I convey such a long story and such complex feelings?

I grabbed at him desperately; he yelped and shook me off. Then he shouted:

``I told you, it's out of the question! I've been through plenty myself, so I get it, but not working is your own choice---don't interfere with other people's work!''

The bus rose into the air. The shadow was torn from the bus, falling far away.

``Please... don't take it... don't take it...''

My mind was sluggish; I could only repeat myself, crying like a small child.

The transport truck carried the bus away.

The void left where the bus had been hurt unbearably.


\subsection*{3}

I became a shut-in. Everything in the world was unbearably painful.

I didn't look at the internet or watch TV. I turned off my phone, removed the doorbell to refuse visitors, and put a ``Moved'' sticker on the mailbox. I stocked up on preserved food and stayed shut in my dark, curtained room.

There were no seasons, no morning or noon---only an extended, room-temperature time.

I stayed in bed all the time, eating only about once every three days. My throat was terribly dry, so I drank tap water until I vomited. Even when I wasn't thirsty, I felt thirsty.

By my pillow, I had pinned several drawings of the secret base's abandoned bus with thumbtacks. So that I could dream of Yuzuki, even a little. I lived only to wake up and see the next dream.

Every time I opened the curtains, the seasons had changed. Like a vivid magic trick, the scenery of spring, summer, autumn, and winter swapped within the window frame. I watched without emotion. Like an audience member who, no matter how beautifully a magician makes flowers bloom in the air, can only think of popping popcorn at a steady rhythm. No, perhaps the audience was the seasons. When the curtain opened, I wasn't there. A disappearance trick.

Everything passed by me, or through me.

As if there was a fatal hole somewhere in my body. Whatever I saw, ate, or thought, it all fell through the hole, and I just kept shriveling up.

When I slept, I turned to the right. I could see the corner of the room. Even though the room was sealed, dust and fine sand gradually accumulated in the corner.

A little desert, I thought. Deserts are born in corners like that, spread when no one's watching, eventually cover and swallow everything. And so the windowless room becomes a moonless desert.

I was a castaway who had lost all sense of direction. No stars in the night sky, no turbaned desert people passing by, not even camel dung. Just waiting to dry out slowly.

With the same frequency as opening the curtains, I'd open my MacBook.

Each time, there was an email from Furuta.

``How are you?''

``Are you writing?''

``I want to read your new work, Yakumo-kun~''

``Aren't you writing?''

``You haven't given up on writing, have you?''

``Don't give up on writing! No no no no!''

``You have talent! Write novels!! For me!!!''

I couldn't tell if he was indifferent, loving, or just selfish.


\subsection*{4}

Before I knew it, two years had passed. The flow of time was frighteningly fast.

I'd aged for nothing, accumulated nothing---except a lot of trash.

I turned on my smartphone for the first time in two years. Many messages had come. The call history was mostly Furuta and Shimizu in a ratio of about 1:9. Looking at the messages, I saw that Shimizu had been very worried about me. She'd sent short messages almost every day like ``You okay?'' or ``Everything alright?'' What must it feel like to keep writing letters for two years without them being read...?

One message caught my eye:

``I'm marrying Kobayashi Koyomi!''

The date was a year ago. She'd sent many messages asking me to come to the wedding, but I hadn't even seen them. Four months ago, the wedding had already taken place.

---Just then, my phone rang.

It was Shimizu. She must have noticed the messages were read and called.

The ringing continued... My heart pounded. My hand holding the phone trembled.

Two years---I'd lived without speaking to anyone. I'd already forgotten how to make my voice.

Feeling sorry for Shimizu, I turned off the power.

A terrifying silence descended---


\subsection*{5}

My food ran out. I lay on the bed, idly wondering what to do.

To live, I had to go outside. But I'd become unable to. The grief of losing Yuzuki hadn't healed at all, and two years had weakened me entirely; I didn't feel I could survive outside.

Anyway, what was I supposed to live for?

I'd only been able to go on because Yuzuki existed. As long as she was alive somewhere in the world, playing the piano---that alone saved me and gave me meaning.

Without Yuzuki, I had no reason to exist.

---I thought, I'll die.

I felt no fear. The boundary between life and death seemed nonexistent.

I bit off the tip of my tongue. The pain was dull, as if numbed. The taste of blood spread in my mouth. I was surprised at how much blood there was.

I stood up and went to the bathroom. I looked in the mirror---a hideous man stared back.

Was this me? I thought. My hair was overgrown, my beard wild, my skin pale as a corpse; despite sleeping so much, there were dark circles under my eyes, my gaze was hollow, my cheeks sunken, and I could see the shape of my skull. I stuck out my tongue; blood dripped from its tip onto the white sink.

There's no one uglier than you... I said to myself in my mind.

It would be better for the world if I just died...

I took a razor from the drawer. I held it in my right hand and pulled out my tongue with my left.

I pressed the cold blade to my tongue, and with one motion---

Just then, I heard an auditory hallucination.

``Tokatonton''---

Suddenly, I lost motivation. Dazai Osamu's words came back to me.

``Thinking of suicide, tokatonton.''

I didn't know why I heard that hallucination and lost motivation like a character in a story. It was irrational---but since it was gone, it was gone. I lay on the bed, deciding to wait for starvation. The lukewarm darkness, the same temperature as my blanket, slowly began to swallow me---

``Tokatonton''---

Huh? Even starving to death is out? You can't get motivated to do nothing? I tilted my head, got up again, and forced myself to think of something to do.

Suddenly, I remembered Yuzuki's video. If I was going to die, I wanted to watch it first---

I took out the camera from the closet and struggled to connect it to my unnecessarily large LCD TV. Whether due to malnutrition or brain atrophy, I couldn't connect a simple cable, gritted my teeth, got frustrated, lay down, got up, and started over. I fumbled like that for a while, but strangely, I didn't hear ``Tokatonton.''

Finally, I got it connected.

The video began to play.

---Yuzuki smiled. She smiled and waved at me.

She was sitting at the living room table. The background was a big window, a winter blue sky, white walls, a wall clock like an orange cross-section. It was early December, right after Miller's visit, in Milan. ``...It's pretty hard, I'm shaking a bit.''

My voice. Yuzuki laughed softly.

``Then today is practice.''

She added sugar and milk to her mug of coffee and stirred. I, with my pointless artistic aspirations, filmed it from above.---And then her smiling face appeared again. She drank her coffee. The camera slowly moved to capture her profile.

I loved Yuzuki's profile.

The scene changed to Yuzuki's back. The strings of her blue apron tied in a bow. The nostalgic rhythm of \textit{sutoton-toto}. When I crept up with a lecherous step, she noticed and half-turned, laughing: ``What now?'' I filmed her, finding pointless beauty in the cut vegetables. Yuzuki was smiling.

The scene changed again; Yuzuki was walking. She moved lightly through the streets of Milan. The sun reflected off the windows of the houses, dazzling. The wind softly swayed her jet-black hair. She was just walking, but it was like a movie---

---No, it actually \emph{was} a movie.

I wasn't the director. I was just the cameraman; the director was Yuzuki.

The footage had already been edited from the moment it was shot.

Yuzuki, pretending to just be filmed, had been making a film with a message.

The message was simple and clear, anyone could read it. Whenever I pointed the camera at her, Yuzuki would smile. And as long as she was on screen, she'd keep a gentle smile. Always. Even as time passed, even as her fingers disappeared, as her hands, as her feet disappeared, that smile never disappeared from the frame.

``I am happy''---

That was what Yuzuki wanted to convey in this film. From the very beginning, she had been sending that one thought to me, watching the video at this very moment.

Like the end of \textit{Cinema Paradiso}, a film of kisses continuing forever, it's a film of Yuzuki's smiles, continuing like a gentle rain---

I cried. I thought I cried, but no tears came. I had become so dry that not a single tear would fall. The room had become a desert, my body a mummy.

The film was nearing its end---

The wedding footage appeared. Shimizu had filmed it. The ring exchange, the vows, the cake cutting, our friends who'd grown up and come in suits...

The happy wedding ended in no time. I had no memory of filming after that. I'd been absorbed in our happy married life, and then Yuzuki went to the hospice.

The screen went black.

The happy movie was over. The cinema was closing---

Just as I heard that phantom announcement.

Suddenly, a light flicked on inside the screen.

Yuzuki appeared, in pajamas. She was sitting elegantly in her wheelchair in front of an electric piano. In her hand was a remote control for the light. Her gold wedding ring gleamed. She said:

``Good evening, Yakumo-kun. And, probably, long time no see---''

I looked away from the TV, to the right. The spot in front of the electric piano was empty. The clock like an orange cross-section hung on the wall. I turned back to the screen.

``I am happy now, thanks to you, Yakumo-kun---'' Yuzuki smiled softly, tilting her head. ``And you, Yakumo-kun...?''

I stared at the screen without blinking. The clock like an orange cross-section showed a different time...

``If you're happy now, please turn off the screen right now. And forget about me, blankly. Like a cute chicken that forgets everything after three steps... And please, keep smiling forever. Now, go ahead---''

Yuzuki sat beautifully, still as a statue, waiting for me to turn off the screen.

I couldn't move at all. I wasn't happy at all. I no longer knew what happiness was.

Yuzuki sighed and began to speak.

``---If you're still watching, it means you haven't found happiness... The only thing I, as I'm dying, wish for is your happiness. That's all...''

Yuzuki lowered her head. I felt sad. I'm sorry, Yuzuki...

``I think I can see you. In a dark room, all alone, not eating properly, all skin and bones, hunched over...''

Yuzuki stared straight at me. I felt as if she saw through my ugly self, and I squirmed in shame.

Her expression softened. It was a sad, loving expression.

``You really can't do without me, can you... That actually makes me a little happy, though I know it's selfish---but I can't say that anymore. I can't be by your side. You have to live in a world without me. Did we say a proper goodbye? ...Maybe we didn't. Then let's say a proper goodbye now---''

Yuzuki paused the video, moved the TV, and instructed me to sit farther away. And to put the vase containing her ashes nearby. I did as she said.

``---Are you ready? Now, I will perform a miracle. A once-in-a-lifetime miracle. After that, you and I will say goodbye forever. So, with all your being, watch and feel it.''

She took a deep breath. Then she said:

``---Now must be your hardest time. You must be in your most terrifying night. I want to save you from it, and bring you back to a bright world. So now, I will go back in time to save you.''

The light clicked off, and the screen went dark.

Then, an orange, circular light appeared.

I gasped.

Yuzuki was in the room.

In front of the piano, lit by a small desk lamp, she smiled gently.

``I've come to save you, Yakumo-kun.''

She smiled like a boy.

``Yuzuki---'' I reached out involuntarily. She stopped me.

``Don't move, or the magic will break---''

The magic would break---indeed, she was right. By narrowing the light to a circle, she had simply blended the TV's rectangular frame with the darkness, making it look as if she was inside the screen.

A simple trick---but for me, it was nothing short of magic.

Yuzuki had truly come back in time to save me.

``---Now, I'm going to play my final performance. My very last, just for you. AGATERAM isn't working well anymore, so it won't be a beautiful piece.---So, what I'm about to play is a happy song, one that will make you cheerful and able to walk again. Actually, it's a song I wrote when I was five. Critics would call it terrible, but I love this terrible song. Just as I love terrible you---''

Yuzuki began to play. A happy, bouncing melody flowed out.

It was a very strange tune. But it was so happy, so full of energy, it made you want to march on forever---. AGATERAM sometimes acted up, mixing in noise. But Yuzuki effortlessly incorporated it into the music, making it even more cheerful. As if even a lonely kid sulking in the corner of the classroom could become an important friend.

It was like playing a toy piano. Nothing difficult at all. But it was so clear and nostalgic, a music that seemed to reach heaven---

---Just then, the strange dream I'd seen before flashed through my mind.

In a dark room with the curtains drawn, I was watching TV.

On the screen was Shimizu, in his wedding suit, laughing \textit{wahaha}. Then Aida appeared, also in his wedding suit, with his thumbs in his ears, wiggling his fingers, blowing on a toy whistle. Next, Yuzuki's former manager, Kitajo Takashi, appeared, taking pictures with a crazed smile, the flash blinding.

I remembered the rest of that dream.

I found myself inside the TV screen, becoming the cameraman.

A vast flower field stretched as far as the eye could see. The sky was a mysterious strawberry-milk color.

Behind Shimizu, Aida, and Kitajo stretched a long line of people---Kobayashi Koyomi, Sakamoto, and many others, all in their wedding attire. Each one, upon reaching the camera, did something strange, like a little monster, making funny noises and movements, then rejoined the procession. The line went on and on: my first-year homeroom teacher, Sumida-sensei, Sekigahara who'd wanted to escape Fukushima, the psychiatrist who'd gone a little crazy, the burly Chopin player from Warsaw airport, Rhythm the golden retriever---everyone was there. They marched on, full of energy, like strange, cheerful monsters. Everyone was smiling.

``Yakumo-kun!''

I turned around. Yuzuki was there. Smiling. In her wedding dress, seated in her wheelchair.

``Let's go!''

I also smiled, tossed the video camera aside, and carried Yuzuki like a princess, joining the procession. Then the Great Demon King, Shimizu, appeared with a \textit{wahaha}, effortlessly carrying a ridiculously large grand piano. Yuzuki smiled happily, tapping the keys with AGATERAM. A cheerful, beautiful music flowed out, and we became a festival procession, marching on forever---

The music Yuzuki played in my dream was the same as the one in the video.

I realized it. I'd had that dream three days before Yuzuki lost AGATERAM---that night, I must have heard her playing in my sleep. And the music had entered my dream, showing me that strange vision. Like a fossilized myth coming back to life, the forgotten dream vividly revived---

Aida's whistle was loud even in the dream. Kobayashi Koyomi stretched her short limbs, a little Godzilla. Roppongi-senpai had powered up to become Senbon-gi, closer to a Gundam than a human. Everyone was being silly, having fun. The beautiful woman in the twelve-layered kimono was Haruhime from the Uname Festival; the one flirting with her was probably her fiancé, Jiro. Next to me was Urashima Taro,\footnote{Urashima Taro (浦島太郎) is a Japanese folk hero who visits the Dragon Palace under the sea.} riding a turtle that swam through the air, waving a fishing rod dangerously. The seven dwarfs surrounded a cute Snow White. Looking closer, even various fictional characters were marching with us. Characters from \textit{JoJo's Bizarre Adventure} were posing with their signature moves, and the Disney characters moved with their usual cheerful animation. Everyone was smiling.

Beside them, coughing, was a thin, sharp-nosed, brown-haired man who couldn't quite keep up with the manga characters but marched on with vigor. I was shocked.

Ah, it's Chopin---

The woman with him was probably George Sand, and their children, Maurice and Solange. They'd had a complicated relationship of love and hate, but now they seemed to be getting along, marching together happily. Everyone was smiling.

Behind them, a young but brave soldier, wearing a helmet and carrying a gun.

A boy who had fought in the Warsaw Uprising.\footnote{The Warsaw Uprising of 1944 was a major World War II operation by the Polish resistance to liberate Warsaw from German occupation.}

Behind him, more Polish soldiers and civilians who had died in the uprising followed with their guns. The boy at the front looked proud. They were gentle, noble Polish warriors, letting the child lead. A rumbling noise came from behind; it was the city itself, marching. The old town of Warsaw, destroyed by the German army. The beautiful, nostalgic streets marched alongside the dead, as if clinging to them. I thought, oh good. They were still with their homeland. Everyone was smiling.

Behind them were Japanese people. I didn't know them personally, but I recognized them from the internet and news footage.

The people who died in the Great East Japan Earthquake.\footnote{The Great East Japan Earthquake and tsunami of March 11, 2011, caused widespread devastation and loss of life in the Tohoku region.}

People who couldn't escape the tsunami, people who died trying to save others, people who died from illness after the disaster, people who took their own lives from grief---everyone was there. Now they were marching happily. The destroyed landscape of their homeland marched with them, protecting them. Everyone, everyone was smiling.

\textit{Wahaha!} laughed Shimizu. Yuzuki played the piano. The living, the dead, the stories---all marched together. We were a cheerful night parade of a hundred demons. A bright, lively procession of prayer, marching to the ends of the earth---!

Suddenly, a cliff appeared before us. The ground split from left to right.

``Piece of cake,'' Shimizu said. I nodded.---We jumped lightly. At that moment, Yuzuki's body floated up. The piano floated too. Yuzuki rose into the air, still playing. Chopin floated. George Sand and her children floated. The people who died in the Warsaw Uprising, the Warsaw Old Town, the people who died in the earthquake, and their homeland scenery---all floated.

The dead soared into the sky, looking comfortable and happy---

We, the living and the stories, landed on the other side of the cliff and stared up at the sky in awe.

The music was reaching its climax.

The sky split open, and a dazzling light shone through.

Heaven, I thought. Everyone was going to heaven.

\textit{Woof woof woof!} A dog ran through the sky to Yuzuki. It was Melody, the golden retriever who loved Yuzuki so much she'd wet herself. Melody flew alongside Yuzuki. Yuzuki began to smile. Then a grand piano and a gentle, dignified woman flew in.

It was Tanaka Kiyoko-sensei.

She played a crystalline piano tone in harmony with Yuzuki's music. She'd been unable to play due to collagen disease, but now she could play freely, looking so happy. Yuzuki was overjoyed to play with her beloved teacher.

Then even Chopin-sensei appeared. On his Pleyel piano. The genius of improvisation, he gently added his touch, making Yuzuki's five-year-old composition even more joyful and beautiful.

Other people with instruments gathered one after another, forming a free concerto.

A great \textit{Baooon!} sounded, and a huge white creature appeared from the horizon.

A giant pure white whale---! I shivered. It was the sand whale I'd drawn with Yuzuki's salt. The whale carried the dead on its back, taking them first-class to heaven. I worried about Yuzuki a little. Was there electricity in heaven? Could she read sheet music at night? I picked a lily of the valley and tossed it into the air; it spun and flew, becoming a Warsaw streetlamp, lighting everyone's way. In heaven, anything goes. Everything is okay.

A majestic music resounded through the sky. Music that only the dead could play.

What a noble, gentle music it was.

It must be the nobility and gentleness of the dead's souls.

The dead are saved by the beauty of their own souls---

Yuzuki would surely meet my mother in heaven. In a warm, sunny place. And they'd share the oranges growing nearby.

We on the ground waved cheerfully, saying goodbye. Sad and lonely, but everyone was smiling. Everyone smiled, praying for the dead's peace.

We started walking again, another lively procession of prayer. Then:

``Tokatonton''---

Huh? I turned around and saw a carpenter building a new house.

\textit{Tokatonton-tokatonton-tontotekan}---

A strong, rhythmic sound. We became a strange orchestra, marching through an endless flower field.

\textit{Wahaha-pyupiro-gao-skatan-bogubogu-wellington-chararira-sutoton-toto-sukoton-pappa-shassha-pirupiru-kachanko-tonton-pururuwoon-pyon-katakata-totekankan-tokatonton-tontotekan}---!

The noise of everyday life became our prayer music.

And a new town began to chase us---

Then I woke up, feeling refreshed, tears on my face.

As if something I thought could never be saved had been saved. I had already been saved when I heard Yuzuki's performance in my dream. I'd just finally remembered it.

Yuzuki finished playing and turned to me. Her expression was that of a goddess, full of love.

``---How was it? I played with all my heart, so that you'd feel better. So that you'd be saved.''

I nodded. Over and over. But no tears came.

As if she'd seen right through me, she said:

``If you're so weak that you can't even cry, please take a pinch of my salt. I want to become your tears. Salty tears, washing away your grief, carrying you into tomorrow. I want to let you breathe in a bright world again. And I want to become your life---''

I looked at the vase containing Yuzuki's salt. And---I took a pinch.

I brought my trembling fingers to my tongue.

Salty. And painful.

---I remembered the taste of tears.

In that moment, all the memories of tears I'd ever shed came back. Sad tears, frustrated tears, happy tears, tears of surprise, bittersweet tears, even the tears from yawning---all of them came back, nostalgic and dear. As if my empty head were suddenly filled with a brilliant bouquet.

And as if blinded by the too-bright light, I cried. I sobbed.

I cried so much it felt like I'd flood the desert of a room.

``Now you'll be okay---'' Yuzuki said, as if she could really see me. ``You can live strongly toward the future now.---Well then, I'll return to the past.''

The desk lamp clicked off, plunging the room into darkness.

Then the room light came on, and Yuzuki was on the other side of the screen again.

She said, looking lonely:

``The miracle only happens once, so you can't watch this again, okay? Well then, this is---''

Just then, on the other side of the screen, I shouted something loudly. Yuzuki flinched and turned toward the bedroom. After a surreal pause, she turned back and laughed.

``Yakumo-kun, was that sleep-talking? What did you say?''

I remembered. In my dream, I had shouted:

``Yuzuki's piano is the best!''

Yuzuki smiled and waved.

``Well then, bye-bye, Yakumo-kun. Take care. Don't catch a cold. When you become a wonderful old man, come see me---''

And the video ended.

I thought, I will live.

And---I'll write a novel. I'll write about Yuzuki.

About Yuzuki, who had just saved my heart.


\subsection*{6}

I made up my mind and stepped outside---

The light stung my eyes. I squeezed them shut and raised my hand to shield my face. Still, tears spilled from the brightness. I opened my eyes little by little---

A blue sky.

Petals danced in the air. Cherry blossom petals. They had been carried by the wind from the neighbor's house. I felt the scent of the season so strongly it was overwhelming.

The smell of spring---it was already the third spring since I'd shut myself in.

I took a step forward. My vision swam; I staggered and put my hand against the wall. My arms, nothing but skin and bone, felt so thin they might break. My knees shook. I took another step. My feet were heavy as lead. I braced myself against falling. I felt that if I fell here, I'd never get up again. I moved forward, little by little, as if wading through mud. I reached the street. A car sped past. Even though it was far away, it startled me so much I nearly fell. I couldn't keep up with the speed and information overload of the world.

The fear, my own patheticness, the dazzling light---tears streamed down my face.

Even so, I gritted my teeth and stepped forward, finally reaching the nearby FamilyMart. It was only a hundred meters, but I felt like I'd walked all the way from Nara to the capital.

The automatic door opened. The clerk who said ``Welcome!'' did a double-take at the sight of me. Of course. My hair was long enough to cover my face, I was nothing but skin and bones, and I was gasping like I was about to die. There couldn't be a creepier customer. I looked around and went to the bento section at the back. The other customers stared in shock.---I didn't care. I'm a monster. Ugly, but saved by Yuzuki.

I filled a basket with as much food as it could hold and brought it to the register. The female clerk with three piercings in her left ear, half-stunned, half-frightened, scanned the barcodes one by one.

I had something to tell her. But I'd forgotten how to speak.

I practiced on the spot. A sound like wind through a hollow tree came from my throat. Once more. This time, an oboe-like ``La'' sound. Tuning complete.---I said:

``Please give me a FamilyMart fried chicken.''

\subsection*{7}

Ramen, soba, sandwiches, pork cutlet curry, pork bowl, chicken rice...

I ate everything. My stomach had shrunk from years of undereating, but I managed to stuff it down over time. I had to restore my withered body.

After eating, I sat at my desk and opened my MacBook. I typed an email to Furuta:

``I'm sorry for worrying you. I'll write again. Please be the first to read it.''

My bangs were in the way... I went to the bathroom and roughly chopped my hair with the razor I'd tried to cut my tongue with. I could fix it later. Now I wanted to write as soon as possible.

By the time I got back to my desk, Furuta had already replied:

``Thank you!! I'll be waiting!!''

Tears welled up. I also sent a message to Shimizu:

``Shimizu, sorry for worrying you. Sorry I couldn't make it to your wedding. Thank you for all the messages. I'll pick myself up little by little, and someday I'll come see you.''

The reply came quickly:

``Okay! I'll be waiting, Yacchan!''

Everyone around me was so kind.

I began to write.

I started from the day I was told my mother had salt disease. More than two years of absence had rusted my hands and brain; I could only write like a grade-schooler. I couldn't even beat my first-grade self. The particles were off, the sentences didn't flow. It was as hard to read as if a nasty stone was waiting to trip the reader with every step. No emotion. It was like a complicated washing machine manual disguised as a novel.

At this rate, Yuzuki's kindness wouldn't come through. I rewrote it again and again.

Starting the next day, I also rebuilt my life. I opened the windows, cleaned the room, ate proper meals, and took two walks a day. I spent all my remaining time on writing.

Seasons passed in an instant.

And then, October 17---Yuzuki and Chopin's death anniversary arrived. The third anniversary of Yuzuki's death.

A message came from someone dear. From Emil, the maker of AGATERAM.

After his usual lengthy greeting, he included a URL. I clicked it and was taken to a video-sharing site. It was a mechanism that allowed only a select few to view the video. Only I could see it. With my heart pounding, I played the video.

Yuzuki appeared---

She wore a white dress and AGATERAM. She bowed, placed the Anpanman doll on the piano, and played the \textit{Barcarolle}---a beautiful sound that stole my soul. When she finished, she stood and bowed. My eyes filled with tears, and I couldn't help but applaud.

The footage cut to a commemorative photo of me, Yuzuki, Miaha, and Emil. Everyone looked happy. Miaha hugged Yuzuki tightly from the side, her sky-blue eyes sparkling. Yuzuki gently held her shoulder with AGATERAM.

The camera slowly zoomed in on Miaha's face. Her face faded, replaced by that of a slightly older girl.

A beautiful girl. Wavy blonde hair, pale skin, cute freckles on her cheeks, and sparkling sky-blue eyes---

I gasped. It was Miaha, three years older.

The camera slowly zoomed out.

I gasped again. On Miaha's arms, AGATERAM gleamed. It was a new model, more refined and beautiful. In those arms, she held the charm Yuzuki had given her---the Anpanman doll. Her dress was the same sky-blue as that day.

Miaha smiled and said, in slightly awkward Japanese:

``Hello, long time no see. I am Miaha Kamiński. I am 9 years old. I still can't believe three years have passed since then. On that day three years ago, Yuzuki-sensei gave me love and courage. After that, I was able to go to school regularly, studied hard, and made many friends.

And my father made AGATERAM for me. Every day since then, I have practiced the piano I love. At first, it didn't go well at all, and I felt discouraged, but whenever I looked at Anpanman, I felt encouraged and could keep smiling. Now, I love piano more than three meals a day. I love it.

It's all thanks to Yuzuki-sensei---When she passed away, I was very sad. I cried for days and days. But I also think this---''

Miaha placed her silver palm over her heart.

``Sensei is always inside me---

Inside me, she always plays a beautiful piano melody.

If I believe that, every wind will become sensei's music.

Today, I will play with a prayer, so it reaches sensei in heaven---''

Miaha placed the Anpanman doll on the grand piano and sat down.

On her left, transparent light streamed through the window behind her.

Miaha took a deep breath. Then AGATERAM began to move smoothly.

A forte, powerfully launching into the silence---

A black-key accompaniment, like a swaying boat---

Tears overflowed from my eyes.

The \textit{Barcarolle}---

A beautiful melody flowed out. Clear, pure, lovely notes. Pearls of sound, like tears, sparkling and rising into the sky.

Miaha, born without forearms, must have struggled greatly to operate AGATERAM. Without the memory of moving her forearms, controlling the myoelectric signals must have been difficult. But she showed no sign of that struggle. She played with a natural beauty, as if born with the piano.

Yuzuki's figure flashed in my mind. When I first met her, playing the piano like a scrap of spring sky whimsically descended---

It overlapped with Miaha's figure.

Miaha must have listened to Yuzuki's performance again and again. Praying that someday she could play as beautifully, she had practiced every single day.

That was why her sound carried Yuzuki's sound. And Yuzuki's sound had been inherited from Tanaka Kiyoko-sensei. She, too, had inherited it from someone. Like the clear meltwater of Mount Fuji, polished by the earth, a sound polished by long, long time, a prayer---

It was as if fate itself were sounding, I thought.

Surely, it reached Yuzuki in heaven.


\subsection*{8}

I wrote a long thank-you email. She would probably have a friend translate it at length. And since the kind, tearful Emil would surely shed tears, imagining it made me smile.

The next day, a long reply came, asking me this:

``I would like to make this video public. I think it could give courage and hope to many people.---However, I have one concern. That is, the video might become commercial. AGATERAM will finally be released next month. This video would be a powerful promotional tool. But I also don't want to commercialize Yuzuki-san. I'm in a dilemma...''

I remembered the CD cover of ``SADNESS'' after the earthquake. The disaster had been packaged, commercialized, consumed. A similar situation was at issue here.

Was it okay to make money off this beautiful relationship between Yuzuki and Miaha?

But I answered immediately:

``Please make it public. You can use it commercially. AGATERAM is a work that can give hope to many people, beyond being just a product. Yuzuki received hope from AGATERAM. And so did I. Yuzuki would be proud that it can reach so many people. No promotional fee is needed. It's a wonderful thing.''

---That was good. To send something out is to risk being consumed. Important feelings can never reach only their intended recipients. They might generate money along the way, or be spat on by someone heartless.

But that's fine. Like a bouquet hiding a cannon, let it fly far while accepting being consumed. And let it become a weapon for someone facing hardship and despair.

---When the video was made public, it became an instant topic and spread unbelievably.

Many comments came from around the world.

``I lost my hand in an accident last month and was filled with sorrow. But now I'm full of hope!''

Reading that English comment made my heart warm. When AGATERAM went on sale, it sold out in an instant. An affordable, beautiful silver hand was extended to those who needed it. That hand would surely be extended to someone else next. Copernicus Technologies later made substantial donations to the Great East Japan Earthquake, the Kumamoto Earthquake, and other organizations.


\subsection*{9}

During all this, the title of my novel was decided.

Like a single bubble rising slowly from the depths of water, the title naturally emerged from deep within my mind.

``I Want to Become Your Tears''

It was the words Yuzuki had said to me. I understood that was the core of this story.

After that, words came naturally. I just had to reconstruct a story that already existed somewhere. But still, I wrote with difficulty and pain. I wrote crying many times. I wrote like a prayer. That Yuzuki's soul be saved. That the souls of the dead be saved. That the hearts of those who read this story be saved. Writing with that prayer was, above all, my own salvation.

And in January---the manuscript was finally complete.

I trembled with a sense of accomplishment. I had finally written the story I needed to write. Only after finishing did I understand why Yuzuki had so rapidly ``story-fied'' her life at the end.

Stories can truly save people.

By ``story-fying'' herself, Yuzuki had saved herself and even saved me.

I felt myself slipping into deep exhaustion, but I pulled myself together and revised it over and over. I polished it until I could do no more with my current ability.

Then, suppressing my pounding heart, I sent it to Furuta.

``Thank you for waiting. It's finally done. Please be the first to read it.''

A reply came immediately.

``I've been waiting!! Thank you!! Thank you!!''

What a nerve-wracking time. An hour passed, two hours---

He was taking much longer than before. I fidgeted for three hours, and finally the reply came.

``Wonderful! It's a masterpiece! I'll enter it in the new writer award. It's sure to win!''

I let out a breath.

``Thank you!''

``Thank you. Being the first to read this story is my pride!''

My chest grew warm. I leaned my head against the chair and looked at the ceiling. Outside the window was a clear winter sky. A warm, gentle sunlight streamed in.

For the first time in a long while, I had nothing to do.


\subsection*{10}

I walked around uselessly, thinking about what my next novel would be.

---Then, I suddenly thought.

What had happened to that abandoned bus?

I contacted the company, OMOYA Construction, and learned its surprising destination.

I began preparing to go see it, waiting for the right time. I couldn't start my next novel, but I wasn't anxious. I believed it would win.

Even if it didn't, I could just try again. Like the streets of Warsaw, like the disaster areas, from zero, from minus, I could start over as many times as I needed---

And so the story returns to the beginning.

% ==============================
%        Epilogue
% ==============================
\clearpage
\thispagestyle{plain}
\phantomsection
\begin{center}
    \large\bfseries Epilogue
\end{center}
\addcontentsline{toc}{chapter}{Epilogue}
\label{sec:epilogue}
\vspace{1cm}

Deep, deep down into the sea we dove. The seawater gradually grew darker, richer in color---

Eventually, we reached the seabed. Depth: 30 meters. Because of the overcast sky, visibility was poor. Detritus drifted like a snowy night.

I signaled to Shimizu and turned on my flashlight.

In the sudden beam of light, the form of an abandoned bus, sunk on the seabed, was illuminated. The bus loomed large and black, covered in barnacles, giving it a slightly eerie appearance.

I made an X with my arms in front of my chest: not this one.

Shimizu nodded and swam to the right side of the bus---

Another abandoned bus appeared.

``Artificial reefs''---

That was the true identity of these buses sunk on the seabed. Man-made objects are submerged in the sea to encourage fish breeding by providing them with habitat. At this location, five buses had been sunk, aligned with their heads facing west-southwest. Fish I didn't know the names of lay motionless inside the buses.

---There.

I found it. At the far end, I spotted that nostalgic abandoned bus. Its once-endearing face was now buried under shells and seaweed, as if sleeping.

I signaled to Shimizu; he nodded and stopped where he was, letting me go alone. I nodded back and swam to the entrance on the left side of the driver's seat. All windows and doors had been removed. The entrance gaped open, a hollow square.

I closed my eyes and let my thoughts drift, then slowly entered.

---Inside the bus, it was midnight.

A cold, silent winter's night---

Cold moonlight streamed through the windows. Snow fell silently, fluttering.

A girl was sobbing softly. At the bottom of her most terrifying night, all alone.

Like a lonely white shadow forgotten in the gaps of time---

I had returned to my child self. With each step, the wooden floor made a \textit{kotsu, kotsu} sound.

When I stood before Yuzuki, she looked up with a start.

``Sorry to keep you waiting---'' I smiled. ``I've come back through time to save you.''

Yuzuki's eyes widened. ``What do you mean---?''

``It's a little trick only novelists can use. Just like you used a trick to come save me from the future. I used a little trick to go back to the past and come save you. Actually, I wrote that whole long novel just for this---''

Yuzuki tilted her head slightly and said:

``...I don't really understand, but somehow, it feels like I do...''

I held out my hand and said:

``Let's go. Wherever, together---First, um...''

``Lake Inawashiro.\footnote{Lake Inawashiro (猪苗代湖) is a large lake in central Fukushima Prefecture, the fourth largest in Japan by surface area. It is a popular tourist destination.}''

Yuzuki said that, smiling, and stood up.

And she took my hand.

---

Back on the boat, Shimizu and I dried our wet bodies and lay down on the deck.

Shimizu's prosthetic leg gleamed silver. It was AGATERAM Leg, a new model from Emil and his team. Thanks to it, Shimizu had been able to dive with me. I reflected on how strange the alignment of the stars can be.

The sky had cleared, and the weather was calm. A flock of seagulls flew overhead.

``Shimizu---'' I said. ``Thank you.''

``It's fine. We're friends, aren't we?''

``Thank you for being my friend.''

Shimizu nodded. After a long pause, he murmured:

``...I hope your novel gets published.''

``Me too. But even if it doesn't, that's okay too.''

``Huh? Really?''

I chuckled at his puzzled expression and said:

``If that happens, I'll just write a harmless romantic comedy.''

The flock of seagulls flew off toward the other side of the sea.

Then we returned Yuzuki's salted remains to the sea. Just as with my mother, we scooped it up little by little and scattered it, and finally floated many flowers on the water. The emptiness of the now-empty vase no longer brought me pain.

Watching the flowers drift slowly away on the waves, I let my thoughts drift toward the future.

What would become of me from now on? What would become of Fukushima?

I didn't know. I didn't know, but I hoped things would work out. I hoped Fukushima's fish would be bought again. I hoped the broken towns would be rebuilt. I hoped the children would return to the elementary schools that had lost their students. I hoped the wounds of losing something precious would heal.

I thought, surely it would be okay. The emptiness left by loss doesn't ache forever. \textit{Zal} takes a long time to heal, but eventually it does. And it is filled with prayer, and something more beautiful than before emerges.

Like Shimizu hitting a home run with a prosthetic leg.

Like the resurrected Warsaw Old Town.

Like Yuzuki playing with AGATERAM.

Like the tone Miaha plays.

And it becomes a new hope for someone unknown to me yet.

Who would read the novel I had written? I hoped many people would. I hoped it would flow away like music, half consumed and half leaving something behind. If it could save even a little of what only stories can save, and make the reader's tomorrow even a little brighter, nothing could make me happier.

I wanted to become a part of someone's soul, their flesh and blood---not in such a grandiose way.

I, too, like Yuzuki, felt this way:

I want to become your tears.

—— THE END ——


% ==============================
%        Afterwards
% ==============================
\clearpage
\thispagestyle{plain}
\phantomsection
\begin{center}
    \large\bfseries Afterwards
\end{center}
\addcontentsline{toc}{chapter}{Afterwards}
\label{sec:afterwards}
\vspace{1cm}

In publishing this work, I received help from many people.

I would like to take this opportunity to express my heartfelt gratitude to everyone at Shogakukan Inc.,\footnote{Shogakukan Inc. (株式会社小学館) is a major Japanese publisher of magazines, manga, and books, founded in 1922.} to director Iso Mitsuo\footnote{Iso Mitsuo (磯光雄) is a Japanese anime director and screenwriter, known for works such as \textit{Dennō Coil}. He served as a guest judge for the 16th Shogakukan Light Novel Award.} for his kind words on the obi as a guest judge for the 16th Shogakukan Light Novel Award,\footnote{The Shogakukan Light Novel Award (小学館ライトノベル大賞) is a literary prize for new light novel authors, sponsored by Shogakukan.} to my editor Hamada for his invaluable support, to illustrator Yanagi Sue\footnote{Yanagi Sue (柳すえ) is an illustrator known for her delicate and beautiful artwork.} for the beautiful cover and interior illustrations, to Imai Riko, a participant in the 18th Chopin International Piano Competition,\footnote{The Chopin International Piano Competition (ショパン国際ピアノコンクール) is one of the world's most prestigious piano competitions, held every five years in Warsaw, Poland.} for sharing her insights on the competition, to my friends who read the work and encouraged me, to the great predecessors who paved the way, and to the characters of this story.

I will continue to do my best to write many more interesting novels from now on.

Thank you all so much.

I should add that while the high school names appearing in the Koshien-related episodes are real, there is no real-life model for Shimizu, the batter with a prosthetic leg.

\textbf{References:}

--- \textit{The Heart of the Genius Chopin: Chopin's Letters}, Frédéric Chopin, translated by Harada Mitsuko (Daiichi Shobō)\footnote{Daiichi Shobō (第一書房) was a Japanese publishing house active from the early to mid-20th century.}

--- \textit{Chopin: The Cannon Hidden Among Flowers}, Suzumi Yoshiyo (Iwanami Junior Shinsho)\footnote{Iwanami Junior Shinsho (岩波ジュニア新書) is a series of compact books aimed at younger readers, published by Iwanami Shoten.}

--- \textit{Villon's Wife}, Dazai Osamu (Shinchō Bunko)\footnote{Dazai Osamu (太宰治, 1909–1948) was one of Japan's most celebrated novelists, known for works such as \textit{No Longer Human} and \textit{The Setting Sun}. \textit{Villon's Wife} (ヴィヨンの妻) is one of his later short stories. Shinchō Bunko (新潮文庫) is a major Japanese paperback imprint published by Shinchōsha.}

\end{document}